% Shared mathematical module. Input exactly once in each host document.
% No document class, package loads, macros, or theorem environments are defined.
% Requirements and bibliography entries: INTEGRATION_MAP.md.
% Source locators and snapshot provenance: ES_NIEMEIER_PROVENANCE.md.
% All labels and citation keys introduced here use the esn: namespace.

\section{The modulo-840 lattice, triality, and exact shell maps}
\label{esn:sec:shared}

This section constructs the original modulo-$840$ lattice $N$ with
root system $A_5^4D_4$, its full integral isometry group, and the exact
maps from denominator shells and finite prime samples to the retained
residue coordinates. The original data and source labels are recorded in
\cite{esn:Provenance}. The shared continuation proves the index-six Leech
neighbour $\Lambda$ and the marked isometry $\mathcal L$ into the three
octonion sectors. In the specified Cayley frame, put
$\tau(q)=\operatorname{Re}((q_1q_2)q_3)$ and $F=\tau\circ\mathcal L$;
the Albert determinant is $\lambda^3-\lambda(y,y)+2F(y)$.
The value-generated groups have quotient
$\mathcal V_F(\Lambda)/\mathcal V_F(N)\cong\mathbb Z/12$
(Theorem~\ref{nca:thm:value-modules}). The three residue-cycled neighbours
meet in $J$, with determinant $81$ and index $9$ in each; the fixed
fourth extension has exactly $72$ roots of type $A_2^{12}$
(Theorems~\ref{esrf:thm:common-J} and~\ref{esrf:thm:fourth-extension}).
The common lattice and this fourth extension have the exact same
value-generated group
$\frac1{27}\mathbb Z\oplus\frac{\sqrt2}{216}\mathbb Z
\oplus\frac{\sqrt3}{18}\mathbb Z$.
The inclusions from $\mathcal V_F(N)$ and into $\mathcal V_F(\Lambda)$
have matrices $\operatorname{diag}(1,4,1)$ and
$\operatorname{diag}(1,1,3)$ and quotients $\mathbb Z/4$, $\mathbb Z/3$
(Theorems~\ref{jmv:thm:JM-values} and~\ref{jmv:thm:chain}).
Wilson's construction \cite{esn:Wilson2009} is realized with both metrics
and the direction $c_{\rm W}=c^{-1}$ retained. Its fixed Gram matrix is
$3G_{E_8}$, its complement is explicitly isometric to $A_2\otimes E_8$,
and the index is $3^8$, with quotient $\mathbb F_3^8$. Its literal
intersections with $\mathcal L(N)$ and $\mathcal L(\Lambda)$ are the
same rank-four lattice of Gram matrix $6G_{D_4}$, minimum squared norm
$12$ and value-generated subgroup $4\sqrt2\mathbb Z$ for $\tau$
($8\sqrt2\mathbb Z$ for $2\tau$). This is not the actual image:
$12\sqrt2$ is not attained by $\tau$ on that intersection.
All coordinates, multiplication conventions, domain lattices and maps
are specified in the complete proofs below.

\subsection{The finite multiplicative carrier and its four fibres}
\label{esn:subsec:residues}

For an odd prime \(q\) and a unit \(r\) modulo \(q\), write
\(\left(\frac rq\right)=1\) when \(r\) is a square modulo \(q\),
and \(-1\) otherwise.  Let
\begin{equation}
 \begin{split}
 V_{840}&=\{r\in(\mathbb Z/840\mathbb Z)^\times:
                                      r\equiv1\pmod {24}\},\\
 \beta(r)&=\left(\frac{1-\left(\frac r5\right)}2,
                  \frac{1-\left(\frac r7\right)}2\right)
                          \in\mathbb F_2^2,\\
 S_{840}&=\ker\beta,\qquad F_\eta=\beta^{-1}(\eta)
                                      \quad(\eta\in\mathbb F_2^2).
 \end{split}
 \label{esn:eq:residue-carrier}
\end{equation}
Multiplication is used in \(V_{840}\) and addition in \(\mathbb F_2^2\).

\begin{theorem}[The marked cyclic fibre]
\label{esn:thm:cyclic-fibre}
The map \(\beta\) is an onto group homomorphism, each \(F_\eta\) has
six elements, and
\begin{equation}
 \begin{split}
 S_{840}&=\{1,121,169,289,361,529\}=\langle289\rangle,\\
 (289^j)_{j=0}^5&=(1,289,361,169,121,529)\pmod {840}.
 \end{split}
 \label{esn:eq:ordered-fibre}
\end{equation}
Each fibre is a torsor under multiplication by \(S_{840}\).  In
particular the marking
\begin{equation}
 \kappa:S_{840}\xrightarrow{\ \cong\ }\mathbb Z/6\mathbb Z,
 \qquad r=289^{\kappa(r)}\pmod {840},
 \label{esn:eq:kappa}
\end{equation}
is an isomorphism of groups.
\end{theorem}

\begin{proof}
The Chinese remainder map modulo \(24,5,7\), whose moduli are pairwise
coprime, identifies \(V_{840}\) with
\((\mathbb Z/5)^\times\times(\mathbb Z/7)^\times\).  The powers of
\(2\) modulo \(5\) are \(1,2,4,3\), and the powers of \(3\) modulo
\(7\) are \(1,3,2,6,4,5\); these prove the respective cyclic orders
four and six.  In a cyclic group of either even order, the squares are
exactly the even powers.  Therefore each Legendre character is the
homomorphism recording exponent parity and is onto.  Their product is
onto \(\mathbb F_2^2\).  Its kernel has
\((4/2)(6/2)=6\) elements and every fibre has that same cardinality.

Successive multiplication of the six displayed residues by \(289\)
gives the next entry and then returns from \(529\) to \(1\).
They are distinct.  Moreover \(289\equiv1\pmod {24}\),
\(289\equiv4\pmod5\), and \(289\equiv2\pmod7\), so it belongs to
the simultaneous square kernel.  Its six powers therefore exhaust that
kernel.  If \(r,s\in F_\eta\), then \(sr^{-1}\in S_{840}\) and
\(s=(sr^{-1})r\), uniquely.  This proves the torsor assertion and all
claims about \(\kappa\).
\end{proof}

Fix the coordinate order
\[
 (\eta_1,\eta_2,\eta_3,\eta_4)=(00,10,01,11)
\]
and choose one origin \(f_i\in F_{\eta_i}\) in each fibre, with
\(f_1=1\).  The six ordered elements of its fibre are
\(f_i289^j\), \(0\le j\le5\).  These origins are part of the
coordinate identifications used below; no selection of them is inferred
from an individual denominator certificate.

\subsection{The root lattice and its exact discriminant form}
\label{esn:subsec:root-and-discriminant}

For a finite set \(F\), give the free abelian group
\(\mathbb Z[F]=\bigoplus_{r\in F}\mathbb Ze_r\) the metric
\((e_r,e_s)=1\) for \(r=s\) and zero otherwise.  Put
\[
 A(F)=\left\{\sum_{r\in F}n_re_r:\sum_{r\in F}n_r=0\right\}.
\]
The ordered fibre coordinates identify each \(A(F_{\eta_i})\)
isometrically with
\[
 A=\{(n_0,\ldots,n_5)\in\mathbb Z^6:\textstyle\sum_j n_j=0\}.
\]
In a separate Euclidean space with orthonormal basis
\(e_1,e_2,e_3,e_4\), define
\begin{equation}
 D=\{x\in\mathbb Z^4:x_1+x_2+x_3+x_4\equiv0\pmod2\},
 \qquad R_{\mathrm{ES}}=
       \left(\bigoplus_{i=1}^4 A(F_{\eta_i})\right)\mathbin{\perp}D.
 \label{esn:eq:root-lattice}
\end{equation}
All five summands are orthogonal with precisely their displayed metrics.
The symbol \(R_{\mathrm{ES}}\) denotes this lattice; the unadorned
\(R\) later used for an integer shell modulus has a different declared
domain.

For a full integral lattice \(L\) in its real span, let
\[
 L^\vee=\{x:(x,L)\subset\mathbb Z\},\qquad D_L=L^\vee/L.
\]
If \(L\) is even integral, its discriminant quadratic form is
\(q_L(x+L)=(x,x)\pmod {2\mathbb Z}\).  This is well-defined because
replacing \(x\) by \(x+\ell\) changes its squared norm by
\(2(x,\ell)+(\ell,\ell)\in2\mathbb Z\).

\begin{theorem}[Root metric and marked discriminant]
\label{esn:thm:root-discriminant}
The lattice \(R_{\mathrm{ES}}\) is even integral of rank 24 and has
Gram matrix
\[
 G_{R_{\mathrm{ES}}}=G_A^{\oplus4}\oplus G_D,
\]
where
\[
 G_A=\begin{pmatrix}
 2&-1&0&0&0\\-1&2&-1&0&0\\0&-1&2&-1&0\\
 0&0&-1&2&-1\\0&0&0&-1&2
 \end{pmatrix},\qquad
 G_D=\begin{pmatrix}
 2&-1&0&0\\-1&2&-1&-1\\0&-1&2&0\\0&-1&0&2
 \end{pmatrix}.
\]
Their determinants are 6 and 4, respectively, and
\(\det R_{\mathrm{ES}}=6^4\cdot4=5184\).

In each identified \(A\) summand use
\(\lambda=(5,-1,-1,-1,-1,-1)/6\) as the discriminant generator.
For \(D\), use the representatives
\begin{equation}
 v=(1,0,0,0),\qquad s=\tfrac12(1,1,1,1),\qquad
 c=\tfrac12(-1,1,1,1)
 \label{esn:eq:d4-representatives}
\end{equation}
and the exact marking \(10\mapsto v+D\), \(01\mapsto s+D\),
\(11\mapsto c+D\).  Then
\begin{equation}
 \begin{split}
 D_{R_{\mathrm{ES}}}&=(\mathbb Z/6\mathbb Z)^4\oplus\mathbb F_2^2,\\
 q_{R_{\mathrm{ES}}}(x_1,x_2,x_3,x_4;y)
   &=\frac56\sum_{i=1}^4x_i^2+\mathbf1_{y\ne0}
                                              \pmod {2\mathbb Z}.
 \end{split}
 \label{esn:eq:full-discriminant}
\end{equation}
\end{theorem}

\begin{proof}
For a six-element ordered fibre the five differences
\(e_j-e_{j+1}\), \(0\le j\le4\), form a basis of its augmentation
lattice: the coefficients of a zero-sum vector in this basis are its
first five successive partial sums.  Dot products give \(G_A\).
The determinant \(d_n\) of the analogous tridiagonal \(n\)-by-\(n\)
matrix satisfies \(d_n=2d_{n-1}-d_{n-2}\), \(d_0=1,d_1=2\), by
first-row expansion; induction gives \(d_n=n+1\).
The four vectors
\(e_1-e_2,e_2-e_3,e_3-e_4,e_3+e_4\) lie in \(D\) and their
coordinate determinant has absolute value 2.  Their integer span
therefore has index 2 in \(\mathbb Z^4\).  The parity homomorphism
\(\mathbb Z^4\to\mathbb F_2\), \(x\mapsto\sum x_i\pmod2\), is
onto and has kernel \(D\), so this span equals \(D\).
Their dot products give \(G_D\), whose determinant is the square of
the coordinate determinant, namely 4.  Integer coordinates make all
pairings integral.  In either lattice the squared norm is congruent
modulo two to the coordinate sum, hence is even.  Rank and determinant
now follow from the orthogonal direct sum.

The vector \(\lambda\) has sum zero and integral pairing with each
augmentation root.  Its norm is \(5/6\), and \(k\lambda\in A\)
holds exactly when \(6\) divides \(k\).  Thus it has exact order six
in \(D_A\).  In an integral basis of an integral lattice with Gram
matrix \(G\), the dual coordinates are \(G^{-1}\mathbb Z^n\);
therefore \(|L^\vee/L|=|\det G|\).  Since \(|D_A|=6\), this
element generates the full discriminant group and gives its stated form.

A vector pairing integrally with all \(e_i-e_j\) has common
fractional parts in its four coordinates.  Pairing also with
\(e_i+e_j\) forces this common fractional part to be either zero or
one half.  Conversely every entirely integral or entirely half-integral
vector pairs integrally with every root of \(D\).  There are two
integer classes according to sum parity, and two half-integer classes
according to the parity of their difference from \(s\).  The four
representatives \(0,v,s,c\) are therefore all the classes, and their
squared norms are \(0,1,1,1\).  Their addition is as marked because
\(v+s-c=(2,0,0,0)\in D\).  The dual and quadratic form of an
orthogonal sum are the direct sums of those of its summands, which proves
\eqref{esn:eq:full-discriminant}.
\end{proof}

\subsection{The index-72 glue and all glue-class minima}
\label{esn:subsec:glue}

Define the binary and ternary codes
\begin{equation}
 \begin{split}
 C_2&=\operatorname{span}_{\mathbb F_2}
       \{(1111;00),(1100;10),(1010;01)\},\\
 C_3&=\operatorname{span}_{\mathbb F_3}
       \{z_1=(1,1,1,0),\ z_2=(1,2,0,1)\}.
 \end{split}
 \label{esn:eq:codes}
\end{equation}
The semicolon distinguishes the four binary A5 coordinates from the
two D4 discriminant coordinates.  Let
\begin{equation}
 \begin{split}
 \mathcal H&=\{(3u+2z;y):(u;y)\in C_2, z\in C_3\}
                                  \subset D_{R_{\mathrm{ES}}},\\
 N&=\{x\in R_{\mathrm{ES}}^\vee:
                                x+R_{\mathrm{ES}}\in\mathcal H\}.
 \end{split}
 \label{esn:eq:glue-and-overlattice}
\end{equation}
Here \(3u+2z\) is computed coordinatewise modulo six using the
integer representatives of \(u\) and \(z\).

\begin{theorem}[Even unimodular completion with exactly 144 roots]
\label{esn:thm:niemeier}
The subgroup \(\mathcal H\) is isotropic of order 72.  Its inverse
image \(N\) is a positive-definite even unimodular lattice of rank 24,
with \([N:R_{\mathrm{ES}}]=72\).  Among the 71 nonzero glue classes,
46 have minimum squared norm 4 and 25 have minimum squared norm 6.
The roots of \(N\) are exactly those of \(R_{\mathrm{ES}}\), with
root system \(A_5^4D_4\): 30 roots in each A5 summand and 24 in D4.
\end{theorem}

\begin{proof}
The binary generators are independent, hence give eight words.  Listing
their sums shows that the seven nonzero words comprise one word
\((1111;00)\) and six words whose first part has weight two and whose
second part is nonzero.  A ternary word is
\[
 az_1+bz_2=(a+b,a+2b,a,b).
\]
Its last two coordinates prove injectivity of the parameter pair
\((a,b)\), so there are nine words.  If \(a=0\ne b\) the third
coordinate alone vanishes; if \(b=0\ne a\) the fourth alone vanishes;
and if both are nonzero exactly one of \(a+b,a+2b\) vanishes.
Thus every nonzero word has weight three, and each of the four
coordinates is the unique vanishing coordinate for two words.

The map \((u_i,z_i)\mapsto3u_i+2z_i\pmod6\) is bijective:
reduction modulo two recovers \(u_i\), and reduction modulo three
recovers \(2z_i\), from which multiplication by 2 recovers \(z_i\).
It is a homomorphism of additive groups, since changing binary or
ternary representatives by their moduli changes \(3u_i+2z_i\) by a
multiple of six.  Its product with the retained binary \(y\)-coordinate
therefore maps \(C_2\times C_3\) injectively onto the subgroup
\(\mathcal H\).  It follows that \(|\mathcal H|=8\cdot9=72\).  For
\(u_i\in\{0,1\}\) and \(z_i\in\{0,1,2\}\),
\[
 \frac56(3u_i+2z_i)^2
 =\frac{15}{2}u_i+10u_iz_i+\frac{10}{3}z_i^2
 \equiv\frac32u_i+\frac43z_i^2\pmod {2\mathbb Z}.
\]
The cross-term is even; a nonzero \(z_i\) has square 1 or 4, whose
contributions in the last expression differ by 4.  The binary
contribution including D4 is therefore zero, 6, or \(3+1=4\), and
the ternary contribution is zero or \((4/3)3=4\) modulo
\(2\mathbb Z\).  This proves isotropy.

The inverse image is a lattice because it is a subgroup between the
full lattices \(R_{\mathrm{ES}}\) and \(R_{\mathrm{ES}}^\vee\).
Isotropy gives \((x,x)\in2\mathbb Z\) for every \(x\in N\), and
\[
 (x,y)=\frac{(x+y,x+y)-(x,x)-(y,y)}2\in\mathbb Z
                                    \qquad(x,y\in N).
\]
Thus it is even integral.  Its index is \(|\mathcal H|=72\).
For a finite-index inclusion of equal-rank lattices, the covolume
changes by that index, and the Gram determinant is the square of the
covolume in orthonormal coordinates.  Hence
\[
 \det N=\frac{6^4\cdot4}{72^2}=1.
\]
The dual-index formula proved above implies \(N=N^\vee\).
Its rank and positive definiteness are inherited from the unchanged real
span and metric.

For \(0\le k\le5\), every vector in the A5 class \(k\lambda+A\)
has the form
\[
 m-\frac{k}{6}(1,1,1,1,1,1),\qquad
 m\in\mathbb Z^6,\quad\sum m_i=k.
\]
Its squared norm is \(\sum m_i^2-k^2/6\).  Since
\(m_i^2\ge m_i\) for every integer \(m_i\), the minimum is at least
\(k-k^2/6\).  It is attained by choosing exactly \(k\) of the six
entries equal to one and the others zero.  Thus
\begin{equation}
 \mu_A(k)=\frac{k(6-k)}6.
 \label{esn:eq:a5-coset-minimum}
\end{equation}
Every nonzero integral D4 class has minimum one, attained by \(v\);
every half-integral vector has norm at least
\(4(1/2)^2=1\), attained by \(s,c\) in their respective classes.
The four A5 minima and the D4 minimum add and are simultaneously attained
because the metric is an orthogonal sum.

For a coordinate with \(k_i=3u_i+2z_i\pmod6\), its minimum
contribution is 0 for \((u_i,z_i)=(0,0)\), \(3/2\) for
\((1,0)\), \(4/3\) for \((0,z_i\ne0)\), and \(5/6\) for
\((1,z_i\ne0)\).  The complete nonzero case table is
\begin{equation}
\begin{array}{c|c|c|r|l}
 \operatorname{wt}(u)&\operatorname{wt}(z)&
 |\operatorname{supp}(u)\cap\operatorname{supp}(z)|&
 \text{classes}&\text{minimum squared norm}\\\hline
 2&0&0&6&2(3/2)+1=4\\
 4&0&0&1&4(3/2)=6\\
 0&3&0&8&3(4/3)=4\\
 2&3&1&24&(3/2)+2(4/3)+(5/6)+1=6\\
 2&3&2&24&(4/3)+2(5/6)+1=4\\
 4&3&3&8&(3/2)+3(5/6)=4.
\end{array}
 \label{esn:eq:glue-minimum-table}
\end{equation}
For each of the six weight-two binary words, four nonzero ternary
words omit a coordinate inside its support, and four omit a coordinate
outside.  This proves the two counts \(6\cdot4=24\); all other counts
follow from the enumerated binary and ternary words.  Thus the total
number at minimum 4 is \(6+8+24+8=46\), and the total at minimum 6 is
\(1+24=25\).  Their sum is 71, leaving exactly the zero glue class.

No nonzero glue class can contain a norm-two vector.  In the zero
class, an A5 norm-two vector has exactly one coordinate 1 and another
\(-1\), so is \(e_i-e_j\), with \(6\cdot5=30\) possibilities.
A D4 norm-two vector has exactly two coordinates of absolute value one,
so is \(\pm e_i\pm e_j\), with \(4\binom42=24\) possibilities.
An orthogonal sum of two nonzero root-lattice vectors has norm at least
4 because each summand is positive even.  These are therefore all
\(4\cdot30+24=144\) roots of \(N\).
\end{proof}

A Niemeier lattice means a positive-definite even unimodular lattice of
rank 24.  The theorem constructs one with root system \(A_5^4D_4\)
directly.  Its root type appears with lambency 6 in the primary umbral
table \cite[Table~3, p.~21]{esn:CDH2014}.  This table reference is not
needed to prove the inverse-image construction or any minimum above.

For comparison with the earlier labeled glue presentation in the source
\cite{esn:ESMain2026,esn:Provenance}, put
\[
 h_1=(1,3,1,1;00),\qquad h_2=(1,4,2,3;11),\qquad
 h_3=(3,0,3,0;01),
\]
and let \(\mathcal H_{\mathrm{old}}\) be their subgroup.  The exact
transport is
\begin{equation}
 P(x_1,x_2,x_3,x_4;y)=(x_3,x_4,x_1,x_2;y),\qquad
 P(\mathcal H_{\mathrm{old}})=\mathcal H.
 \label{esn:eq:old-code-transporter}
\end{equation}
Indeed, if \(rh_1+th_2+wh_3=0\) with \(r,t\in\mathbb Z/6\)
and \(w\in\mathbb Z/2\), the final two coordinates imply that \(t\)
is even and \(w=0\).  The first and third coordinates then give
\(r+t=0\) and \(r+2t=0\) modulo six, whence \(t=r=0\).
Thus \(|\mathcal H_{\mathrm{old}}|=6\cdot6\cdot2=72\).  Directly,
\[
\begin{split}
 P(h_1)&=3(1111;00)+2(2,2,2,0;00),\\
 P(h_2)&=3(0110;11)+2(1,0,2,2;00),\\
 P(h_3)&=3(1010;01).
\end{split}
\]
The displayed ternary parts are \(2z_1\) and \(2z_1+2z_2\).
Consequently \(P(\mathcal H_{\mathrm{old}})\subset\mathcal H\),
and the equal orders prove equality.  This transporter is realized by
the corresponding isometry permuting the full identified A5 summands.

\subsection{The common D4 root lattice, shortest weights, and triality}
\label{esn:subsec:d4-flag}

Use orthonormal coordinates \(\varepsilon_1,\ldots,\varepsilon_4\)
for the ordinary F4 root realization.  Its long roots and the three
parts of its short roots are
\begin{equation}
\begin{split}
 \Phi_D&=\{\pm\varepsilon_i\pm\varepsilon_j:1\le i<j\le4\},\\
 \mathcal V&=\{\pm\varepsilon_i:1\le i\le4\},\\
 \mathcal S_\pm&=
 \left\{\frac12\sum_{i=1}^4\delta_i\varepsilon_i:
          \delta_i\in\{1,-1\},\ \prod_i\delta_i=\pm1\right\},\\
 \Phi_F&=\Phi_D\cup\mathcal V\cup\mathcal S_+\cup\mathcal S_-.
\end{split}
 \label{esn:eq:root-weight-sets}
\end{equation}
For a nonzero vector \(a\), reflection means the exact map
\(s_a(x)=x-2(x,a)a/(a,a)\).  Write \(W_D\) and \(W_F\) for the
groups generated by the reflections in \(\Phi_D\) and \(\Phi_F\),
respectively.

\begin{theorem}[Integral D4 and the marked triality quotient]
\label{esn:thm:d4-f4-triality}
The map
\begin{equation}
 \iota_D:D\longrightarrow\mathbb Z\langle\Phi_D\rangle,
 \qquad(x_1,x_2,x_3,x_4)\longmapsto
                         \sum_{i=1}^4x_i\varepsilon_i
 \label{esn:eq:d4-isometry}
\end{equation}
is an integral isometry onto the long-root lattice.  It extends to duals
and discriminant forms, sending the shortest-vector sets of the marked
classes \(v,s,c\) respectively to
\(\mathcal V,\mathcal S_+,\mathcal S_-\).

The two matrices
\begin{equation}
 R_D=\operatorname{diag}(-1,1,1,1),\qquad
 H_D=\frac12\begin{pmatrix}
 1&1&1&1\\1&1&-1&-1\\1&-1&1&-1\\1&-1&-1&1
 \end{pmatrix}
 \label{esn:eq:d4-triality-matrices}
\end{equation}
are the short-root reflections \(s_{\varepsilon_1}\) and \(s_c\),
where \(c=(-1,1,1,1)/2\) is read in the \(\varepsilon\) coordinates.
They induce \((s\ c)\) and \((v\ s)\) on the three nonzero
discriminant classes.  There is a split exact sequence
\begin{equation}
 1\longrightarrow W_D\longrightarrow W_F
   \longrightarrow S_3\longrightarrow1,
 \qquad W_D\cong\{\text{even signed permutations of four coordinates}\},
 \label{esn:eq:weyl-triality-sequence}
\end{equation}
whose displayed section is \(\langle R_D,H_D\rangle\cong S_3\).
In particular \(|W_D|=192\) and \(|W_F|=1152\).
\end{theorem}

\begin{proof}
The four D4 basis vectors in Theorem~\ref{esn:thm:root-discriminant}
are among \(\Phi_D\), and every vector of \(\Phi_D\) has even
coordinate sum.  Thus the displayed identity of coordinates has exactly
the stated image, preserves the metric, and extends to duals by the
integer-pairing definition.  In a nonzero integral dual class, norm one
forces one coordinate \(\pm1\) and all others zero.  In a
half-integral class, norm one forces all four coordinates to be
\(\pm1/2\).  The difference from \(s\) has even coordinate sum
exactly when the number of negative signs is even.  This proves the
three shortest-vector identifications with the stated marking.

The reflection identity for \(R_D\) is immediate, and
\(H_D=I-2cc^T\) follows by multiplying the four coordinates of \(c\).
Both are orthogonal involutions.  The first preserves the parity of an
integer coordinate sum.  For \(x\in D\), the four coordinates of
\(H_Dx\) are integers because each numerator is congruent to
\(\sum x_i\) modulo two, and their sum is \(2x_1\).  Thus both
preserve \(D\).  Their effects on the representatives show that
\(R_D\) fixes the class \(v\) and exchanges \(s,c\), whereas
\(H_Dv=s\), \(H_Ds=v\), and \(H_Dc=-c\equiv c\pmod D\).
They therefore preserve all of \(\Phi_F\) by the shortest-vector
description.  On the plane with ordered basis \(\varepsilon_1,c\),
their product has matrix
\[
 \begin{pmatrix}0&-1\\1&-1\end{pmatrix},
\]
whose cube is the identity, while both reflections fix the orthogonal
complement.  Consequently \(R_D^2=H_D^2=(R_DH_D)^3=1\).
Every word in two such involutions reduces to
\((R_DH_D)^k\) or \((R_DH_D)^kR_D\), \(0\le k\le2\).
Their discriminant permutations are the six elements of \(S_3\),
so these six words are distinct and give the stated section.

Reflection in \(\varepsilon_i-\varepsilon_j\) swaps coordinates
\(i,j\).  Reflection in \(\varepsilon_i+\varepsilon_j\) swaps them
and changes both signs.  Hence the long-root reflections generate every
permutation and every even sign change, and each long-root reflection
is of this form.  This proves the assertion about \(W_D\) and its
order \(2^3\cdot4!=192\).  Such a reflection acts trivially on
\(D^\vee/D\): for a root \(a\) of squared norm two,
\(s_a(x)-x=-(x,a)a\in D\) for all \(x\in D^\vee\).

Let \(K=\langle W_D,R_D,H_D\rangle\).  Each of its generators
preserves the long roots, so \(W_D\) is normal in \(K\) by
\(g s_a g^{-1}=s_{ga}\).  The even signed permutations act
transitively on each of \(\mathcal V,\mathcal S_+,\mathcal S_-\),
and the displayed section permutes these three sets transitively.
Thus every short root is \(g\varepsilon_1\) for some \(g\in K\),
and its reflection \(gR_Dg^{-1}\) belongs to \(K\).  All generators
of \(W_F\) therefore belong to \(K\); the reverse inclusion follows
from their definitions, proving \(W_F=K\).
An element in the kernel of its permutation action on the three short
sets preserves \(\mathcal V\), so it is a signed permutation matrix.
Preservation of \(\mathcal S_+\) forces an even number of sign
changes, so the kernel is \(W_D\).  This proves exactness, splitting,
and the order \(6\cdot192=1152\).
\end{proof}

\paragraph{The exact geometric datum and weight marking.}
Let \(\mathbb O\) denote the real normed octonion algebra and let
\(J=\mathfrak h_3(\mathbb O)\) be the real vector space of Hermitian
three-by-three matrices, with Jordan product
\(a\circ b=(ab+ba)/2\).  Define
\[
 \mathbb O\mathbb P^2=\{a\in J:a^2=a,\ \operatorname{tr}a=1\},
 \qquad
 \operatorname{Fl}(\mathbb O)=
 \{(a,b)\in(\mathbb O\mathbb P^2)^2:
                                 \operatorname{Re}\operatorname{tr}(ab)=0\}.
\]
For the compact group \(F_4=\operatorname{Aut}_{\mathbb R}(J,\circ)\),
Mare and Willems prove the homogeneous identification
\(\operatorname{Fl}(\mathbb O)=F_4/\operatorname{Spin}(8)\) and
identify its tangent representation as
\(V_8\oplus S_8^+\oplus S_8^-\)
\cite[Definition~2.2, Proposition~2.4, Proposition~2.7, and
pp.~491--492]{esn:MareWillems2013}.  These are the prior geometric
results used here.  With their torus coordinate forms \(L_i\), their
equations (5.4)--(5.6) assign the vector, positive half-spin, and negative
half-spin weights respectively the patterns
\(\pm L_i\), even half-sign sums, and odd half-sign sums
\cite[pp.~510--511]{esn:MareWillems2013}.
The coordinate map \(L_i\mapsto\varepsilon_i\) retains that sign
marking and identifies these three weight sets with those of
\eqref{esn:eq:root-weight-sets}.

For completeness the half-spin weight patterns admit the following
direct calculation.  In the complex exterior algebra
\(\Lambda^*\mathbb C^4\), let \(a_i\) be exterior multiplication
by the \(i\)-th basis vector and \(b_i\) contraction by its dual.
On a basis wedge, insertion and deletion give
\[
 a_ia_j+a_ja_i=b_ib_j+b_jb_i=0,\qquad
 a_ib_j+b_ja_i=\delta_{ij}I.
\]
Thus \(C_i=a_i+b_i\) and \(D_i=\mathrm i(b_i-a_i)\) square to
one and mutually anticommute for distinct Clifford generators.
Writing \(N_i=a_ib_i\), direct multiplication gives
\[
 \tfrac12 C_iD_i=\mathrm i(N_i-\tfrac12 I).
\]
The four commuting infinitesimal torus operators therefore have weights
\(\sum_i(\mathbf1_{i\in I}-1/2)\varepsilon_i\) on the basis
wedge indexed by \(I\subset\{1,2,3,4\}\).  The even exterior
part has an even number of negative signs, and the odd part has an odd
number.  There are eight one-dimensional weight spaces in each part.
For the vector representation the four real coordinate planes
complexify into their two rotation eigenspaces, with weights
\(\varepsilon_i,-\varepsilon_i\).  This proves the displayed
weight formulas directly in the stated marking.

The map proved here is consequently the actual chain
\[
 D\xrightarrow{\ \iota_D\ }\mathbb Z\langle\Phi_D\rangle,
 \qquad
 \{v,s,c\}\longleftrightarrow
       \{\mathcal V,\mathcal S_+,\mathcal S_-\}
       =\{\operatorname{wt}(V_8),
          \operatorname{wt}(S_8^+),\operatorname{wt}(S_8^-)\}.
\]
The same \(S_3\) acts through the explicitly proved permutations of
these three sets.  The companion higher-rung treatment records the
geometric datum and root-coordinate comparison at the respective labels
\begin{center}
\texttt{h24:thm:octonionic-triality}\\
\texttt{h24:prop:two-Coxeter-systems}.
\end{center}
The exact source pointers and passages are included in
\cite{esn:Provenance}.  The integral metric map
in this section is the map on D4 and its dual.  The definitions make no
map from the whole rank-24 lattice \(N\) into the flag tangent space.

\subsection{The full indicated glue stabilizer and its triality quotient}
\label{esn:subsec:glue-gl2-triality}

All vector spaces over finite fields in this subsection retain their displayed
bases.  Write
\[
 j:\mathbb F_3^2\longrightarrow\mathbb F_3^4,
 \qquad j(a,b)=(a+b,a+2b,a,b),
 \qquad C_3=j(\mathbb F_3^2).
\]
The four row forms defining this injection are
\[
 \ell_1=(1,1),\quad \ell_2=(1,2),\quad
 \ell_3=(1,0),\quad \ell_4=(0,1).
\]
They represent all four one-dimensional subspaces of the dual of
\(\mathbb F_3^2\).  Indeed each nonzero row has either first coordinate zero,
in which case its line is represented by \(\ell_4\), or a unique scalar
multiple with first coordinate one, namely one of
\((1,0),(1,1),(1,2)\).

Let \(E\subset\mathbb F_2^4\) be the even-weight subspace and define
\[
 \tau:E\longrightarrow\mathbb F_2^2,
 \qquad \tau(u_1,u_2,u_3,u_4)
       =(u_2+u_4,u_3+u_4).
\]
The binary code in the index-72 glue is exactly
\[
 C_2=\{(u,\tau(u)):u\in E\}.
\]
To verify this, the independent vectors
\(1111,1100,1010\) form a basis of the three-dimensional space \(E\),
and their images under \(\tau\) are respectively
\(00,10,01\).  These are precisely the three stated binary code generators.
In particular \(\tau\) is onto and
\[
 \ker\tau=\langle1111\rangle.
\]
Explicitly, \(\tau(u)=0\) gives \(u_2=u_3=u_4\), and the even-weight
equation then gives \(u_1=u_4\).

For a permutation \(p\) of \(\{1,2,3,4\}\), use the coordinate convention
\[
 (P_pu)_i=u_{p(i)}.
\]
The permutation preserves \(E\) and \(\langle1111\rangle\), so it induces
a unique invertible map \(B_p\) satisfying
\begin{equation}
 B_p\tau(u)=\tau(P_pu)\qquad(u\in E).
 \label{esn:eq:binary-permutation-map}
\end{equation}
For an entirely explicit formula, the two columns of \(B_p\) are
\(\tau(P_p(1100))\) and \(\tau(P_p(1010))\).
The convention implies \(P_pP_q=P_{q\circ p}\).  Accordingly
\(B_pB_q=B_{q\circ p}\); no reversal is suppressed in what follows.

The subgroup of \eqref{esn:eq:glue-and-overlattice} can therefore be written
\[
 \mathcal H=
 \{(3u+2z;\tau(u)):u\in E,\ z\in C_3\}
 \subset(\mathbb Z/6)^4\oplus\mathbb F_2^2.
\]
Let \(G\) be its stabilizer among the transformations
\[
 g=(\varepsilon,p,B),\qquad
 (x;y)\longmapsto
 \bigl((\varepsilon_i x_{p(i)})_{i=1}^4;By\bigr),
\]
where \(\varepsilon_i\in\{1,-1\}\), \(p\in S_4\), and
\(B\in GL_2(\mathbb F_2)\).  Denote the signed permutation matrix on the
first four coordinates by \(M_g\).

\begin{theorem}[An explicit isomorphism to \(GL_2(\mathbb F_3)\)]
\label{esn:thm:glue-gl2}
Restriction to the ternary code gives an isomorphism
\[
 \Theta:G\xrightarrow{\ \cong\ }GL_2(\mathbb F_3),
 \qquad M_gj=j\Theta(g).
\]
For every \(A\in GL_2(\mathbb F_3)\), its inverse image is specified
uniquely by
\begin{equation}
 \ell_iA=\varepsilon_i\ell_{p(i)}\quad(1\le i\le4),
 \qquad B=B_p.
 \label{esn:eq:inverse-gl2-glue-map}
\end{equation}
In particular \(|G|=48\).
\end{theorem}

\begin{proof}
The elements of \(\mathcal H\) of order dividing two are precisely
\((3u;\tau(u))\), and those of order dividing three are precisely
\((2z;0)\).  Every indicated ambient transformation preserves these primary
parts, so an element of \(G\) satisfies \(M_gC_3=C_3\).  The injection
\(j\) therefore yields one and only one \(\Theta(g)\in GL_2(\mathbb F_3)\).
The binary part is preserved exactly when
\(B\tau(u)=\tau(P_pu)\), because all signs become one modulo two.
Surjectivity of \(\tau\) forces \(B=B_p\).

If \(\Theta(g)=I\), the equality \(M_gj=j\) implies
\(\varepsilon_i\ell_{p(i)}=\ell_i\) for every \(i\).
The four row lines are distinct, so \(p(i)=i\).  Since each \(\ell_i\)
is nonzero and \(1\ne-1\) in \(\mathbb F_3\), also
\(\varepsilon_i=1\) for every \(i\).  It follows that \(B_p=I\) and
\(g=1\).  Thus \(\Theta\) is injective.

Conversely fix \(A\in GL_2(\mathbb F_3)\).  The rows \(\ell_iA\) are
nonzero and belong to four distinct row lines.  Since the \(\ell_i\)
represent every row line, each \(\ell_iA\) has a unique expression
\(\varepsilon_i\ell_{p(i)}\), where the nonzero scalar in \(\mathbb F_3\)
is represented by the exact integer sign \(\varepsilon_i=1\) or \(-1\).
Distinctness of the row lines makes \(p\) a permutation.
The resulting signed permutation satisfies \(M_gj=jA\), hence preserves
\(C_3\).  Choosing \(B=B_p\) preserves \(C_2\) as proved above, hence
preserves \(\mathcal H\).  This constructs an element of \(G\) mapping
to \(A\), proving surjectivity and the stated inverse formula.

Finally
\[
 M_{gh}j=M_gM_hj=M_gj\Theta(h)
 =j\Theta(g)\Theta(h),
\]
so \(\Theta(gh)=\Theta(g)\Theta(h)\).
An invertible \(2\)-by-\(2\) matrix over \(\mathbb F_3\) has eight
choices for its first nonzero column and six choices for its second column
outside the span of the first.  Thus its group order is \(8\cdot6=48\).
\end{proof}

The map to the permutation of the four component labels is more precisely
\(g\mapsto p^{-1}\), because a vector supported on component \(k\) is
sent to component \(p^{-1}(k)\).  Its kernel under \(\Theta\) is
\(\{I,-I\}\).  In fact a matrix preserving all four row lines is diagonal
by preservation of \(\ell_3\) and \(\ell_4\), and its diagonal entries
are equal by preservation of \(\ell_1\).  Conversely scalar matrices
preserve all four lines.  Therefore the component permutation image has
order \(48/2=24\), so is all \(S_4\), with
\[
 G/\Theta^{-1}(\{\pm I\})
       \xrightarrow{\ \overline\Theta\ }PGL_2(\mathbb F_3)\cong S_4.
\]
The center of \(GL_2(\mathbb F_3)\) is exactly \(\{\pm I\}\): a
central matrix commutes with \(\operatorname{diag}(1,-1)\), forcing its
off-diagonal entries to vanish, and commuting with
\(\left(\begin{smallmatrix}0&1\\1&0\end{smallmatrix}\right)\) forces
the diagonal entries to agree.

\begin{theorem}[The exact triality quotient and its quaternion kernel]
\label{esn:thm:triality-q8}
The projection \(\pi:G\to GL_2(\mathbb F_2)\),
\((\varepsilon,p,B)\mapsto B\), is onto, and under \(\Theta\) its
kernel is the quaternion subgroup generated by
\[
 I_q=\begin{pmatrix}0&1\\2&0\end{pmatrix},
 \qquad J_q=\begin{pmatrix}1&1\\1&2\end{pmatrix}.
\]
There is a split exact sequence
\[
 1\longrightarrow Q_8\longrightarrow GL_2(\mathbb F_3)
 \xrightarrow{\ \pi\Theta^{-1}\ }GL_2(\mathbb F_2)
 \longrightarrow1.
\]
The splitting is the subgroup generated by
\[
 A_R=\begin{pmatrix}1&0\\0&2\end{pmatrix},
 \qquad A_H=\begin{pmatrix}1&2\\0&2\end{pmatrix}.
\]
It identifies \(GL_2(\mathbb F_3)\) with \(Q_8\rtimes S_3\) for the
conjugation action of this displayed subgroup on this displayed \(Q_8\).
\end{theorem}

\begin{proof}
The nonzero elements of \(E/\langle1111\rangle\) are represented by
the complementary pairs of two-element subsets
\[
 \{12\mid34\},\qquad\{13\mid24\},\qquad\{14\mid23\},
\]
marked respectively by \(10,01,11\) via \(\tau\).
Thus \(B_p\) is exactly the permutation action on these three
partitions, expressed in the displayed binary marking.

Using \eqref{esn:eq:inverse-gl2-glue-map} gives
\[
 M(A_R)x=(x_2,x_1,x_3,-x_4),\quad
 B_R=\begin{pmatrix}1&1\\0&1\end{pmatrix},
\]
\[
 M(A_H)x=(x_1,x_3,x_2,-x_4),\quad
 B_H=\begin{pmatrix}0&1\\1&0\end{pmatrix}.
\]
These binary matrices act respectively as the transpositions
\((01\ 11)\) and \((10\ 01)\), so they generate all permutations of
the three nonzero binary vectors.  There are
\((4-1)(4-2)=6\) invertible binary matrices; a matrix fixing all nonzero
vectors fixes the basis.  Hence this permutation action identifies
\(GL_2(\mathbb F_2)\) with \(S_3\) and proves surjectivity of \(\pi\).
Its kernel therefore has order eight.

For \(I_q\) the row rule gives
\[
 p=(2,1,4,3),\qquad\varepsilon=(-1,1,1,-1),
\]
and for \(J_q\) it gives
\[
 p=(3,4,1,2),\qquad\varepsilon=(-1,-1,1,1).
\]
Both permutations fix each of the three displayed complementary
partitions, so both matrices lie in the kernel.  Direct matrix
multiplication over \(\mathbb F_3\) gives
\[
 I_q^2=J_q^2=-I,\qquad I_qJ_q=-J_qI_q,
 \qquad I_qJ_q=\begin{pmatrix}1&2\\2&2\end{pmatrix}.
\]
The eight matrices
\(I,-I,\pm I_q,\pm J_q,\pm I_qJ_q\) are pairwise distinct by their
displayed entries.  The relations show this set is a group: powers reduce
using \(I_q^2=J_q^2=-I\), and factors can be reordered using
\(J_qI_q=-I_qJ_q\).  These relations identify it with the quaternion
group \(Q_8\).  Since the kernel has order eight, this is the whole kernel.

The matrices \(A_R,A_H\) satisfy
\[
 A_R^2=A_H^2=I,\qquad
 A_RA_H=\begin{pmatrix}1&2\\0&1\end{pmatrix},\qquad
 (A_RA_H)^3=I.
\]
Every word in the two involutions reduces to one of
\((A_RA_H)^k\) or \((A_RA_H)^kA_R\), \(k=0,1,2\), by alternation
and the cubic relation.  Their images under \(\pi\) are the six distinct
permutations generated by \(B_R,B_H\).  Hence the subgroup has exactly
six elements and maps isomorphically onto \(GL_2(\mathbb F_2)\).
This gives the splitting.  For any \(g\), the unique splitting element
with image \(\pi(g)\) leaves \(g\) times its inverse in the kernel;
the intersection of kernel and splitting is trivial.  Multiplication
therefore realizes the asserted semidirect product with the indicated
conjugation action.
\end{proof}

The matrices in the kernel have an additional marked interpretation:
their projective component permutations are the four permutations
\(1,(12)(34),(13)(24),(14)(23)\).  Each fixes all three partitions, and
these are exactly the kernel of the onto action \(S_4\to S_3\), whose
kernel has order \(24/6=4\).  Thus the preceding exact sequence is the
precise map from the order-48 ternary-code symmetry to the three
\(D_4\) discriminant classes.

\begin{theorem}[A faithful action on the actual index-72 lattice]
\label{esn:thm:gl2-full-lattice-lift}
Fix the four identifications of the \(A_5\) coordinate lattices used in
the glue presentation, and the discriminant marking
\(10\mapsto v,01\mapsto s,11\mapsto c\) of \(D_4\).
The group \(G\cong GL_2(\mathbb F_3)\) has a faithful action by integral
isometries of the index-72 overlattice.  Write
\(\varphi_i:A\longrightarrow A(F_{\eta_i})\) for the restriction of
the fixed free-coordinate isometry \(e_j\mapsto e_{f_i289^j}\).
For \(X_i\in A(F_{\eta_i})\), the \(i\)-th output component is exactly
\(\varepsilon_i\varphi_i\varphi_{p(i)}^{-1}X_{p(i)}\).
In the identified \(A\) coordinates this is the action
\[
 L_g(X_1,X_2,X_3,X_4;d)
 =\bigl(\varepsilon_1X_{p(1)},\ldots,
        \varepsilon_4X_{p(4)};\delta(B)d\bigr),
\]
where \(\delta:GL_2(\mathbb F_2)\to O(D_4)\) is the section given by
\[
 \delta(B_R)=R_D=\operatorname{diag}(-1,1,1,1),\qquad
 \delta(B_H)=H_D=\frac12
 \begin{pmatrix}
 1&1&1&1\\1&1&-1&-1\\1&-1&1&-1\\1&-1&-1&1
 \end{pmatrix}.
\]
Under \(\Theta\), the two displayed full triality lifts are exactly
\(A_R\) and \(A_H\).
\end{theorem}

\begin{proof}
The four identified \(A_5\) summands carry the same Euclidean metric,
so the displayed signed permutation is a root-lattice isometry.  The
matrices \(R_D,H_D\) preserve \(D_4\): \(R_D\) negates one integral
coordinate and preserves parity of the sum; for \(x\in D_4\), each
coordinate of \(H_Dx\) is integral because its numerator has the same
parity as \(\sum x_i\), and the coordinate sum of \(H_Dx\) is \(2x_1\).
Both matrices are orthogonal involutions by multiplication.
Furthermore \(R_D=s_{e_1}\), \(H_D=s_c\), with
\(c=(-1,1,1,1)/2\), and \((e_1,c)=-1/2\).
On the plane with basis \(e_1,c\), their product has matrix
\(\left(\begin{smallmatrix}0&-1\\1&-1\end{smallmatrix}\right)\)
and has cube the identity; both reflections fix the orthogonal
complement pointwise.  The discriminant actions are the two distinct
transpositions \(B_R,B_H\).  The same word reduction as in the preceding
proof therefore shows that \(\langle R_D,H_D\rangle\) has exactly six
elements and that its discriminant action is an isomorphism onto
\(GL_2(\mathbb F_2)\).  Its inverse is the homomorphism \(\delta\).

Thus \(g\mapsto L_g\) is a homomorphism of root-lattice isometries.
It extends to the dual lattice because an isometry preserves the
integer-pairing condition defining the dual.  Its discriminant action is
the specified \((\varepsilon,p,B)\), which preserves \(\mathcal H\).
Consequently it preserves the inverse image
\[
 N=\{x\in R_{\mathrm{ES}}^\vee:
                  x+R_{\mathrm{ES}}\in\mathcal H\}.
\]
It is therefore an integral isometry of \(N\).
If \(L_g\) is the identity on \(N\), it is the identity on the full
real span, because \(N\) contains the full-rank lattice
\(R_{\mathrm{ES}}\).  It is therefore also the identity on
\(R_{\mathrm{ES}}^\vee/R_{\mathrm{ES}}\), in particular on the ternary
code.  Hence \(\Theta(g)=I\), and injectivity of \(\Theta\) gives \(g=1\).
This proves faithfulness.
The exact formulas for \(A_R,A_H\) were established above, so their
lifts are the stated full triality transformations.
\end{proof}

This constructs an explicit finite subgroup of \(O(N)\), including its
action on all five root summands and on their glue, and its exact
\(D_4\)-triality quotient.  The group and its projective quotient agree
with the \(A_5^4D_4\) row of \cite[Table~3, p.~21]{esn:CDH2014};
the isomorphisms here have been proved by their actual code actions.

For clarity, the two full triality elements, denoted
\(\widehat R=L_{\Theta^{-1}(A_R)}\) and
\(\widehat H=L_{\Theta^{-1}(A_H)}\), act in the fixed coordinates by
\[
\begin{split}
 \widehat R(X_1,X_2,X_3,X_4;d)
     &=(X_2,X_1,X_3,-X_4;R_Dd),\\
 \widehat H(X_1,X_2,X_3,X_4;d)
     &=(X_1,X_3,X_2,-X_4;H_Dd).
\end{split}
\]
Their squares and the cube of their product are the identity on the
entire lattice, by the proved faithful group homomorphism.
The product on the first three A5 summands is the indicated
three-cycle; its two negations on the fourth summand cancel.

\subsection{The exact reduced-shell divisor bijection}
\label{esn:subsec:shells}

Let \(p\) be a positive prime with \(p\equiv1\pmod {24}\), and
let the positive integer shell \(R\) satisfy \(R\equiv3\pmod4\).
Retain
\begin{equation}
 a=a_R(p)=\frac{p+R}{4},\qquad
 N_{\mathrm{sh}}=pa,\qquad g=\gcd(R,N_{\mathrm{sh}}),\qquad
 R_0=\frac Rg,\quad N_0=\frac{N_{\mathrm{sh}}}{g}.
 \label{esn:eq:shell-data}
\end{equation}
The integer \(N_{\mathrm{sh}}\) is the source's scalar \(N=pa\);
the subscript distinguishes it from the lattice \(N\) without changing
either object.  In particular \(a,R_0,N_0\) are positive integers,
\(\gcd(R_0,N_0)=1\), and the prescribed first denominator equation is
\begin{equation}
 \frac4p-\frac1a=\frac{4a-p}{pa}
       =\frac{R}{N_{\mathrm{sh}}}=\frac{R_0}{N_0}
       =\frac1b+\frac1c.
 \label{esn:eq:fixed-shell-equation}
\end{equation}
Here \(b,c\) are positive integer denominators, and their order is
initially retained.  Define the finite set
\[
 \mathcal D_{R,p}=
 \{d\in\mathbb Z_{>0}:d\mid N_0^2,
                             \ d\equiv-N_0\pmod {R_0}\},
 \qquad A_R(p)=|\mathcal D_{R,p}|.
\]

\begin{theorem}[Certificates, ordered pairs, and complementary pairs]
\label{esn:thm:shell-bijection}
There is an exact bijection
\begin{equation}
 \begin{split}
 \mathcal D_{R,p}&\longrightarrow
  \{(b,c)\in\mathbb Z_{>0}^2:
                    4/p=1/a+1/b+1/c\},\\
 d&\longmapsto
 \left(\frac{N_0+d}{R_0},
       \frac{N_0+N_0^2/d}{R_0}\right),
 \end{split}
 \label{esn:eq:shell-bijection}
\end{equation}
with inverse \(d=R_0b-N_0\).  The involution \(d\mapsto N_0^2/d\)
corresponds exactly to \((b,c)\mapsto(c,b)\).  Consequently the
number of unordered pairs of remaining denominators is
\begin{equation}
 U_R(p)=\frac{A_R(p)+\delta_{R,p}}2,
 \qquad
 \delta_{R,p}=\begin{cases}1,&R_0\mid2,\\0,&R_0\nmid2.\end{cases}
 \label{esn:eq:unordered-count}
\end{equation}
The unique possible fixed certificate is \(d=N_0\), giving
\(b=c=2N_0/R_0\).
\end{theorem}

\begin{proof}
For a positive solution of \eqref{esn:eq:fixed-shell-equation},
\(R_0b/N_0=1+b/c>1\), so \(d=R_0b-N_0>0\); likewise
\(e=R_0c-N_0>0\).  Multiplying the fraction equation by
\(N_0bc\) gives \(R_0bc=N_0(b+c)\), and therefore
\[
 (R_0b-N_0)(R_0c-N_0)=N_0^2.
\]
Thus \(d\mid N_0^2\), \(d\equiv-N_0\pmod {R_0}\), and
\(e=N_0^2/d\).

Conversely, let \(d\in\mathcal D_{R,p}\) and \(e=N_0^2/d\).
For \(R_0>1\), the coprimality \(\gcd(R_0,N_0)=1\) makes
\(d\equiv-N_0\) a unit, and hence
\[
 e\equiv N_0^2(-N_0)^{-1}\equiv-N_0\pmod {R_0}.
\]
For \(R_0=1\), this congruence is vacuous and holds automatically.
Thus both expressions in \eqref{esn:eq:shell-bijection} are positive
integers.  Substitution into \((R_0b-N_0)(R_0c-N_0)=N_0^2\)
and expansion give the original fraction equation.  Each construction
recovers exactly \(d=R_0b-N_0\) and \(e=R_0c-N_0\), so they
are inverse, proving the bijection and the involution statement.

For a positive divisor, \(d=N_0^2/d\) holds exactly when
\(d=N_0\).  This divisor is a certificate exactly when
\(N_0\equiv-N_0\pmod {R_0}\), that is, \(R_0\mid2N_0\).
Coprimality makes this equivalent to \(R_0\mid2\).
Every other orbit contains two certificates, so counting the fixed
orbits and the two-element orbits gives
\(U_R(p)=\delta_{R,p}+(A_R(p)-\delta_{R,p})/2\), exactly
\eqref{esn:eq:unordered-count}.
\end{proof}

Factor the exact integer \(N_0\) as
\(N_0=\prod_{j=1}^{m}q_j^{e_j}\), with distinct positive primes
\(q_j\) and positive exponents \(e_j\), and define
\begin{equation}
 \begin{split}
 B(N_0)&=\prod_{j=1}^{m}\{-e_j,-e_j+1,\ldots,e_j\},\\
 \lambda_{R_0,N_0}:B(N_0)&\longrightarrow(\mathbb Z/R_0\mathbb Z)^\times,
 \qquad
 \lambda_{R_0,N_0}(\beta)=\prod_{j=1}^{m}q_j^{\beta_j}\pmod {R_0}.
 \end{split}
 \label{esn:eq:exponent-box}
\end{equation}
The group for modulus one is taken to be the one-element group, so that
the same notation also expresses the vacuous congruence in that case.
For \(R_0>1\), every \(q_j\) is a unit by coprimality, so all
negative powers in this formula are defined.

\begin{theorem}[The exact exponent-box map and the two raw channels]
\label{esn:thm:raw-channels}
The map
\begin{equation}
 \beta\longmapsto d=\prod_{j=1}^{m}q_j^{e_j+\beta_j}
 \label{esn:eq:exponent-to-divisor}
\end{equation}
bijects \(\lambda_{R_0,N_0}^{-1}(-1)\) with
\(\mathcal D_{R,p}\).  It retains each valuation and both original
integers \(R_0,N_0\).

In addition, suppose \(\gcd(R,pa)=1\) and \(p\nmid a\).
Write
\[
 a=\prod_{j=1}^{s}q_j^{e_j},\qquad
 B(a)=\prod_{j=1}^{s}\{-e_j,\ldots,e_j\},
\]
and define, for each unit
\(t\) modulo \(R\),
\[
 c_R(a;t)=\#\left\{\beta\in B(a):
                                 \prod_j q_j^{\beta_j}\equiv t\pmod R\right\}.
\]
Then the number of raw divisor certificates, equivalently ordered
denominator pairs, is
\begin{equation}
 A_R(p)=c_R(a;-1)+2c_R(a;-p^{-1}).
 \label{esn:eq:two-raw-channels}
\end{equation}
Its complementary-divisor quotient is the count in
\eqref{esn:eq:unordered-count}, with \(R_0=R\) under these additional
coprimality hypotheses.
\end{theorem}

\begin{proof}
A divisor of \(N_0^2\) has one and only one exponent vector
\(\alpha_j\in\{0,\ldots,2e_j\}\).  The exact change
\(\beta_j=\alpha_j-e_j\) gives the entire displayed box and
\(dN_0^{-1}\equiv\prod_jq_j^{\beta_j}\pmod {R_0}\).
Multiplication of \(d\equiv-N_0\) by the unit \(N_0^{-1}\)
gives precisely the target \(-1\); modulo one both statements are
vacuous.  This proves the first bijection and its inverse.

Under the additional hypotheses, \(R_0=R\), \(N_0=pa\), and
the exponent of \(p\) in \(pa\) is exactly one.  Every positive integer
divisor \(d\) of \(p^2a^2\) therefore has exactly one valuation-index
representation
\[
 d=p^{t+1}\prod_jq_j^{e_j+\beta_j},\qquad
 (t,\beta)\in\{-1,0,1\}\times B(a).
\]
After division by \(pa\), its unit-group residue is
\[
 d(pa)^{-1}\equiv p^t\prod_jq_j^{\beta_j}\pmod R.
\]
Distinct valuation-index pairs may have the same unit-group residue;
they remain distinct integer divisors in the count.
The three conditions for this to equal \(-1\) are
\[
\begin{array}{c|c}
 t&\prod_jq_j^{\beta_j}\pmod R\\\hline
 -1&-p\\0&-1\\1&-p^{-1}.
\end{array}
\]
The negation \(\beta\mapsto-\beta\) is a bijection of \(B(a)\)
and inverts the unit-group value, so it bijects the fibres at
\(-p\) and \(-p^{-1}\).  Summing the sizes of these three disjoint
raw exponent classes proves \eqref{esn:eq:two-raw-channels}.  The
complementary divisor sends \((t,\beta)\) to \((-t,-\beta)\),
so identifying it changes the raw count as proved in
Theorem~\ref{esn:thm:shell-bijection}; it does not leave the factor
two in \eqref{esn:eq:two-raw-channels} unchanged.
\end{proof}

This counting convention corrects the wording attached to the source's
two-channel formula: the displayed sum is a raw count, whereas a count
after identifying complementary divisors is \eqref{esn:eq:unordered-count}.
The exact source locator and formula are retained in
\cite{esn:Residual2026,esn:Provenance}.

\subsection{A fixed-shell count does not descend to the residue fibre}
\label{esn:subsec:non-descent}

\begin{theorem}[Two primes in the same residue with different shell counts]
\label{esn:thm:non-descent}
The two primes \(2521\) and \(3361\) both reduce to \(1\in S_{840}\),
but for the same shell \(R=3\) their raw and unordered counts are
\begin{equation}
 A_3(2521)=0,\qquad A_3(3361)=6,
 \qquad U_3(2521)=0,\quad U_3(3361)=3.
 \label{esn:eq:non-descent-counts}
\end{equation}
One actual denominator certificate in the second case is
\begin{equation}
 \frac4{3361}=\frac1{841}+\frac1{974690}+\frac1{28266010}.
 \label{esn:eq:explicit-denominator-example}
\end{equation}
Thus the raw count, its unordered quotient, and the indicator of
completion of this fixed shell each fail to factor through prime
reduction modulo 840.
\end{theorem}

\begin{proof}
The residue equations are the literal equalities
\(2521=3\cdot840+1\) and \(3361=4\cdot840+1\).
For an explicit finite primality verification, the primes at most 57,
in increasing order, are
\[
 \mathcal P=(2,3,5,7,11,13,17,19,23,29,31,37,41,43,47,53).
\]
The remainders against its first 15 entries for \(2521\), all 16
entries for \(3361\), and its first nine entries for \(631\), are
respectively
\[
\begin{split}
 &(1,1,1,1,2,12,5,13,14,27,10,5,20,27,30),\\
 &(1,1,1,1,6,7,12,17,3,26,13,31,40,7,24,22),\\
 &(1,1,1,1,4,7,2,4,10).
\end{split}
\]
All are nonzero, and the respective integer square roots are 50, 57,
and 25.  The prime 29 has remainders \(1,2,4\) against \(2,3,5\),
the primes at most its integer square root 5.
Every composite positive integer greater than one has a prime divisor
at most its square root: in a factorization into two factors greater
than one, the smaller factor is at most the square root and has a prime
divisor no larger than itself.  These nonzero remainder lists therefore
prove each stated primality.

For \(R=3\),
\[
 a_3(2521)=631,\qquad a_3(3361)=841=29^2.
\]
Both primes are 1 modulo three; also \(631\equiv1\pmod3\) and
\(841\equiv1\pmod3\).  Thus \(\gcd(3,pa)=1\) in both cases.
Both \(a\) values are smaller than their respective primes, so
\(p\nmid a\).  Every additional hypothesis of
Theorem~\ref{esn:thm:raw-channels} is now verified.
The two targets \(-1\) and \(-p^{-1}\) coincide modulo three.

For \(a=631\), all exponents \(-1,0,1\) give residue one, so
the target \(-1\) has count zero.  For \(a=29^2\), among the five
exponents \(-2,-1,0,1,2\), exactly \(-1,1\) give \(-1\) modulo
three because \(29\equiv-1\pmod3\).  Both target counts are
therefore 2 and the raw count is \(2+2\cdot2=6\).  Here
\(R_0=3\nmid2\), so both unordered counts are half the corresponding
raw counts.  This proves \eqref{esn:eq:non-descent-counts}.

For the second prime, retain
\[
 N_{\mathrm{sh}}=3361\cdot841=2826601,
 \qquad d=3361\cdot29=97469,
 \qquad \frac{N_{\mathrm{sh}}^2}{d}=3361\cdot29^3=81971429.
\]
Then
\[
 \frac{N_{\mathrm{sh}}+d}{3}=974690,
 \qquad
 \frac{N_{\mathrm{sh}}+N_{\mathrm{sh}}^2/d}{3}=28266010.
\]
The bijection of Theorem~\ref{esn:thm:shell-bijection} now proves
\eqref{esn:eq:explicit-denominator-example}.  The first prime has
no completion with this prescribed first denominator by its zero
certificate count.  Two arguments with the same residue and unequal
function values cannot be the pullback of a single residue function:
if \(f(p)=h(p\bmod840)\), equality of the residues forces equality
of the values.  Applying this observation to the three displayed
functions proves the claimed failures of factorization.
\end{proof}

\subsection{An actual residue cocycle and a retained-sample Fourier map}
\label{esn:subsec:cocycle-fourier}

For \(r\in S_{840}\), multiplication by \(r\) permutes each fibre.
Let \(\rho(r)\) act by
\begin{equation}
 \rho(r)e_s=e_{rs}\quad(s\in V_{840}),\qquad
 \rho(r)|_D=I,
 \qquad b(r)=e_r-e_1\in A(S_{840})\subset N.
 \label{esn:eq:residue-action-cocycle}
\end{equation}
The two free coordinates in \(b(r)\) belong to the principal fibre.

\begin{theorem}[The full-lattice residue action and its exact cocycle]
\label{esn:thm:residue-cocycle}
The permutations \(\rho(r)\) extend to a faithful homomorphism
\(\rho:S_{840}\to O(N)\).  The map \(b:S_{840}\to N\) satisfies
\begin{equation}
 b(rs)=b(r)+\rho(r)b(s)\qquad(r,s\in S_{840}),
 \qquad b(1)=0,\qquad (b(r),b(r))=2\quad(r\ne1).
 \label{esn:eq:cocycle-identity}
\end{equation}
Thus reduction of any prime into \(S_{840}\), followed by \(b\),
is an explicit lattice-valued map on that prime domain.
\end{theorem}

\begin{proof}
Multiplication permutes the orthonormal free coordinates, preserves
their augmentation kernels, and obeys \(\rho(r)\rho(s)=\rho(rs)\).
On the \(i\)-th fibre its dual discriminant generator is
\[
 \lambda_i=e_{f_i}-\frac16\sum_{t\in F_{\eta_i}}e_t.
\]
Its image differs from it by
\(e_{rf_i}-e_{f_i}\in A(F_{\eta_i})\).  Hence \(\rho(r)\)
acts trivially on the entire A5 discriminant group, since
\(\lambda_i\) generates it.  It fixes D and therefore the whole
discriminant group, in particular \(\mathcal H\).
It follows from the inverse-image definition that it preserves \(N\).
If \(\rho(r)\) were the identity, applying it to
\(e_1-e_{289}\) would give
\(e_r-e_{289r}=e_1-e_{289}\).  Since \(289\ne1\), the
positive coordinate on the left is exactly \(r\), so comparison of
free-coordinate coefficients forces \(r=1\).  This proves faithfulness.

Finally
\[
 b(r)+\rho(r)b(s)
   =(e_r-e_1)+(e_{rs}-e_r)=e_{rs}-e_1=b(rs).
\]
For \(r\ne1\), the two basis vectors in the difference are
orthogonal unit vectors, giving squared norm two.  The stated prime map
is the composition \(p\mapsto p\bmod840\mapsto b(p\bmod840)\),
whose domains have all been specified.
\end{proof}

Fix a positive shell \(R\equiv3\pmod4\), and let \(P\) be any
finite set of positive primes whose reduction lies in \(S_{840}\).
For every \(p\in P\), use its actual count \(A_R(p)\) defined by
the reduced divisor set, including the full data of
\eqref{esn:eq:shell-data}.  Define
\begin{equation}
 \begin{split}
 w_{R,P}(r)&=\sum_{\substack{p\in P\\p\equiv r\ (840)}}A_R(p),\\
 \widehat w_{R,P}(k)&=
       \sum_{r\in S_{840}}w_{R,P}(r)
                     \exp\left(-\frac{2\pi\mathrm i k\kappa(r)}6\right),
             \qquad k\in\mathbb Z/6\mathbb Z.
 \end{split}
 \label{esn:eq:sample-pushforward}
\end{equation}
This is the pushforward along prime reduction, followed by the finite
Fourier transform with the already fixed generator 289.

\begin{theorem}[Inversion and an integral augmentation-valued count map]
\label{esn:thm:sample-fourier}
The function \(w_{R,P}\) is recovered by
\begin{equation}
 w_{R,P}(r)=\frac16\sum_{k=0}^5\widehat w_{R,P}(k)
                      \exp\left(\frac{2\pi\mathrm i k\kappa(r)}6\right).
 \label{esn:eq:fourier-inverse}
\end{equation}
Let \(u=\sum_{r\in S_{840}}e_r\),
\(W_{R,P}=\sum_rw_{R,P}(r)e_r\), and
\(M_{R,P}=\sum_rw_{R,P}(r)=\widehat w_{R,P}(0)\).
Then
\begin{equation}
 Q_{R,P}=6W_{R,P}-M_{R,P}u\in A(S_{840})\subset N
 \label{esn:eq:integral-count-augmentation}
\end{equation}
is a well-defined integral lattice-valued count map.  Put
\(\zeta_6=e^{2\pi\mathrm i/6}\) and
\[
 f_k=\sum_{j=0}^5\zeta_6^{kj}e_{289^j}\qquad(0\le k\le5).
\]
The five vectors \(f_1,\ldots,f_5\) form a basis of
\(A(S_{840})\otimes\mathbb C\), with
\begin{equation}
 \rho(289)f_k=\zeta_6^{-k}f_k,\qquad
 Q_{R,P}=\sum_{k=1}^5\widehat w_{R,P}(k)f_k.
 \label{esn:eq:fourier-augmentation-eigenvectors}
\end{equation}
Both the prime sample \(P\) and the shell \(R\) are retained in the
domain of these maps.
\end{theorem}

\begin{proof}
For an integer \(m\), the sum
\(\sum_{k=0}^5\zeta_6^{km}\) equals 6 if \(6\mid m\).
Otherwise it is the geometric sum
\((1-\zeta_6^{6m})/(1-\zeta_6^m)=0\).
Substituting \eqref{esn:eq:sample-pushforward} into the right side
of \eqref{esn:eq:fourier-inverse} therefore gives
\[
 \sum_{s\in S_{840}}w_{R,P}(s)
       \frac16\sum_{k=0}^5
              \zeta_6^{k(\kappa(r)-\kappa(s))}=w_{R,P}(r),
\]
because \(\kappa\) is bijective.  This proves inversion.

The coefficient sum of \(6W_{R,P}-M_{R,P}u\) is
\(6M_{R,P}-6M_{R,P}=0\), and its coefficients are integers, so
it belongs to the exact augmentation lattice.  The same geometric sum
shows that \(f_k\) has coefficient sum zero for \(1\le k\le5\)
and \(f_0=u\).  Moreover the six \(f_k\) are independent: the
coefficient Fourier transform of \(f_k\) is six in coordinate \(k\)
and zero in all other coordinates, by that same identity.  Since the
augmentation space has dimension five, its five nonzero-frequency
vectors form a basis.  Replacing \(j\) by \(j+1\) in
\(\rho(289)f_k\) gives the multiplier \(\zeta_6^{-k}\), including
its displayed sign.

Inversion in each free coordinate gives
\(W_{R,P}=(1/6)\sum_{k=0}^5\widehat w_{R,P}(k)f_k\).
Multiplication by 6 and subtraction of
\(M_{R,P}u=\widehat w_{R,P}(0)f_0\) prove the second identity
in \eqref{esn:eq:fourier-augmentation-eigenvectors}.
All sums were taken over the specified finite set \(P\) and its
residue image.  The transform therefore recovers these aggregate counts;
it makes no identification between the separate counts of distinct
primes with the same residue.  Their possible inequality is proved in
Theorem~\ref{esn:thm:non-descent}.
\end{proof}

% Input after esn:thm:gl2-full-lattice-lift and before the shell subsection.
% This fragment introduces no packages, macros, or theorem environments.
% New labels have prefix esn:fa:.  Its residue corollary uses the formula
% defining rho, which is restated here so placement before that formula is safe.

\subsection{The entire lattice automorphism group and its Weyl quotient}
\label{esn:fa:subsec:full-automorphisms}

Cheng, Duncan and Harvey define the umbral group as the quotient of
the Niemeier lattice automorphism group by its root-reflection subgroup;
their Table~3 identifies this quotient as \(GL_2(3)\) for
\(A_5^4D_4\) \cite[Section~2.4 and Table~3]{esn:CDH2014}.
The argument here proves that identification in the present residue,
metric, discriminant and glue coordinates, and locates the arithmetic
six-cycle in the full group. It does not claim the abstract group as new.
The displayed complement is the particular lift already constructed here;
no equality of its representation with a simple-root-preserving section
is assumed.

Retain the exact lattice \(N\), the root lattice \(R_{\mathrm{ES}}\),
the four fibre identifications \(\varphi_i\), the discriminant marking,
the glue subgroup \(\mathcal H\), and the lift \(g\mapsto L_g\)
constructed above.  For a lattice \(K\) in a Euclidean space, \(O(K)\)
means the group of real linear isometries of its real span that map
\(K\) onto itself.  For the present root lattice define
\[
 W(R_{\mathrm{ES}})
 =\langle s_\alpha:\alpha\in R_{\mathrm{ES}},\ (\alpha,\alpha)=2\rangle,
 \qquad s_\alpha(x)=x-(x,\alpha)\alpha.
\]
All actions below retain the previously fixed six-coordinate order in
each \(A_5\) summand and the marking \(10\mapsto v\),
\(01\mapsto s\), \(11\mapsto c\) in \(D^\vee/D\).

\begin{theorem}[All isometries of the displayed \(A_5\) lattice]
\label{esn:fa:thm:a5-automorphisms}
Let \(A=\{x\in\mathbb Z^6:\sum_{j=0}^5x_j=0\}\) with its displayed
Euclidean metric, and for a permutation \(\sigma\) of
\(\{0,1,2,3,4,5\}\) let \(Q_\sigma e_j=e_{\sigma(j)}\).
Every element of \(O(A)\) has exactly one expression
\(\epsilon Q_\sigma|_{\mathbb R A}\), where
\(\epsilon\in\{1,-1\}\).  Thus
\[
 O(A)=\langle Q_\sigma|_{\mathbb R A}:\sigma\in S_6\rangle
                  \times\{I,-I\},
 \qquad W(A)=\langle Q_\sigma|_{\mathbb R A}:\sigma\in S_6\rangle
                  \cong S_6.
\]
On the marked discriminant group
\(D_A=(\mathbb Z/6)\lambda\),
\(\lambda=e_0-\frac16\sum_{j=0}^5e_j\), this isometry acts by
\(k\lambda+A\mapsto\epsilon k\lambda+A\).  The kernel of this
discriminant action is exactly \(W(A)\).
\end{theorem}

\begin{proof}
The norm-two vectors are precisely
\(\alpha_{ij}=e_i-e_j\) for \(i\ne j\).  For two distinct such roots,
\[
 (\alpha_{ij},\alpha_{kl})
       =\delta_{ik}-\delta_{il}-\delta_{jk}+\delta_{jl}.
\]
Their inner product is one exactly when they have the same first index
or the same second index.  Indeed equality of first indices, with
distinct second indices, makes the displayed expression one; equality
of second indices, with distinct first indices, does the same.  If both
positive Kronecker terms vanish the expression is nonpositive, and if
both are one then the roots are equal, the excluded case.

Put
\[
 C_i=\{\alpha_{ij}:j\ne i\},\qquad
 -C_i=\{\alpha_{ji}:j\ne i\}\qquad(0\le i\le5).
\]
These twelve sets are exactly the maximal sets of roots in which every
two distinct members have inner product one.  To prove this, a singleton
can be enlarged by keeping its first index and choosing a third index,
so a maximal such set contains a pair.  If the pair is
\(\alpha_{ij},\alpha_{ik}\), with \(j\ne k\), a further root
\(\alpha_{lm}\) has inner product one with both only if
\[
 (l=i\ \text{or}\ m=j)\quad\text{and}\quad
 (l=i\ \text{or}\ m=k).
\]
Since \(j\ne k\), these conditions force \(l=i\).  Thus the whole
set is contained in \(C_i\), and maximality makes it equal to \(C_i\).
For a pair with common second index the same argument, after negating
both roots and every possible further member, gives \(-C_i\).
Each \(C_i\) and \(-C_i\) has five roots and is maximal by this argument.

Form a graph whose vertices are these twelve maximal sets, with an edge
when the two sets have nonempty intersection.  Its intersections are
exactly
\[
 C_i\cap C_j=(-C_i)\cap(-C_j)=\varnothing\quad(i\ne j),
 \qquad
 C_i\cap(-C_j)=
 \begin{cases}\{\alpha_{ij}\},&i\ne j,\\
               \varnothing,&i=j.
 \end{cases}
\]
Consequently it is the bipartite graph on the two displayed six-element
families with every cross-edge except the six pairs \((C_i,-C_i)\).
It is connected: two distinct vertices on the same side have a common
neighbor indexed by any index different from both; vertices on opposite
sides with different indices are adjacent; and \(C_i\) and \(-C_i\)
are joined by the path
\(C_i,-C_j,C_k,-C_i\) for distinct \(i,j,k\).
A connected bipartite graph has only two choices for its ordered
bipartition, because the side of a fixed vertex determines the side of
each other vertex along any connecting path.  Hence an isometry of
\(A\), which permutes roots and preserves their inner products and set
intersections, either preserves both displayed families or exchanges them.

In the first case write
\(C_i\mapsto C_{\sigma(i)}\) and
\(-C_i\mapsto-C_{\tau(i)}\).
The unique opposite-family vertex not adjacent to \(C_i\) is \(-C_i\).
Preservation of that missing edge forces \(\tau(i)=\sigma(i)\)
for each \(i\).  The singleton intersection formula then gives
\(\alpha_{ij}\mapsto\alpha_{\sigma(i)\sigma(j)}\).
In the case exchanging the families, composition with \(-I\) reduces
to the first case, so the root action is
\(\alpha_{ij}\mapsto-\alpha_{\sigma(i)\sigma(j)}\).
The roots \(e_j-e_5\), \(0\le j\le4\), span \(\mathbb R A\),
so the root action determines the whole linear isometry.  These cases
give precisely the asserted expressions.

The permutation \(\sigma\) is determined by the permutation of the
six vertices \(C_i\), after the sign is determined by whether the two
families are preserved or exchanged.  Thus the expression is unique.
Every \(Q_\sigma\) and \(-I\) preserves \(A\), and \(-I\)
commutes with every linear map; uniqueness proves the direct product.
Reflection in \(e_i-e_j\) interchanges free coordinates \(i,j\)
and fixes the others, as follows on substituting this root in
\(x\mapsto x-(x,e_i-e_j)(e_i-e_j)\).
Every permutation is a product of transpositions: each cycle
\((i_1\ i_2\ \cdots\ i_m)\) is
\((i_1\ i_m)(i_1\ i_{m-1})\cdots(i_1\ i_2)\).
These reflections therefore generate precisely the displayed \(S_6\).
Finally
\[
 Q_\sigma\lambda-\lambda=e_{\sigma(0)}-e_0\in A.
\]
Thus the discriminant action is the stated sign.  Since \(\lambda+A\)
has exact order six, its negative is distinct from itself; the action
is trivial exactly when \(\epsilon=1\).
\end{proof}

\begin{theorem}[All isometries of the displayed \(D_4\) lattice]
\label{esn:fa:thm:d4-automorphisms}
For the exact lattice
\(D=\{x\in\mathbb Z^4:\sum_i x_i\equiv0\pmod2\}\),
the action on its marked discriminant group gives the split exact sequence
\[
 1\longrightarrow W_D\longrightarrow O(D)
   \xrightarrow{\ d_D\ }GL_2(\mathbb F_2)\longrightarrow1,
 \qquad
 O(D)=W_D\rtimes\delta(GL_2(\mathbb F_2)).
\]
Here \(W_D\) consists of all signed coordinate permutations with an even
number of minus signs, and \(\delta\) is exactly the section generated
by the earlier matrices \(R_D,H_D\).  In particular \(O(D)=W_F\)
in the root and weight coordinates of
Theorem~\ref{esn:thm:d4-f4-triality}.
\end{theorem}

\begin{proof}
Every isometry of \(D\) preserves its dual: if \(T\in O(D)\),
\(x\in D^\vee\), and \(a\in D\), then
\((Tx,a)=(x,T^{-1}a)\in\mathbb Z\).
It therefore induces an additive automorphism of
\(D^\vee/D\cong\mathbb F_2^2\), giving the homomorphism \(d_D\).
The three nonzero classes have the exact norm-one representative sets
\[
 \begin{aligned}
 \mathcal V&=\{\pm e_i:1\le i\le4\},\\
 \mathcal S_+&=\{\tfrac12\textstyle\sum_i\delta_i e_i:
                                      \prod_i\delta_i=1\},\\
 \mathcal S_-&=\{\tfrac12\textstyle\sum_i\delta_i e_i:
                                      \prod_i\delta_i=-1\},
 \qquad \delta_i\in\{1,-1\}.
 \end{aligned}
\]
in the respective marked classes \(v,s,c\).
For the integral nonzero class, norm one forces exactly one coordinate
\(\pm1\).  A half-integral vector has norm at least
\(4(1/2)^2=1\), with equality exactly when every coordinate is
\(\pm1/2\).  The difference of such a vector from
\(s=(1,1,1,1)/2\) has one coordinate \(-1\) per minus sign,
so it belongs to \(D\) exactly when their number is even.
This proves the three sets and their marking directly.

If \(d_D(T)=I\), then \(T\) preserves \(\mathcal V\).
The vectors \(Te_i\) are orthonormal members of \(\mathcal V\),
so \(Te_i=\delta_i e_{\sigma(i)}\) for a permutation \(\sigma\)
of four indices and signs \(\delta_i\).  Preservation of
\(\mathcal S_+\) forces \(Ts\in\mathcal S_+\), hence
\(\prod_i\delta_i=1\).  Thus the kernel consists only of even
signed permutations.

Conversely, reflection in \(e_i-e_j\) interchanges coordinates
\(i,j\), and reflection in \(e_i+e_j\) interchanges them and
negates both.  Their composition produces a sign change in exactly
those two coordinates.  Any even set of sign changes is a product of
such disjoint pairs, and transpositions generate all permutations.
The root reflections therefore generate every even signed permutation,
and each root reflection is itself such a matrix.  If
\(\alpha\in D\) has norm two and \(x\in D^\vee\), then
\(s_\alpha x-x=-(x,\alpha)\alpha\in D\).  Every root reflection
acts trivially on the discriminant group.  This proves
\(\ker d_D=W_D\) in both directions.

For completeness the section has the original exact matrices
\[
 R_D=\operatorname{diag}(-1,1,1,1),\qquad
 H_D=\frac12
 \begin{pmatrix}
 1&1&1&1\\1&1&-1&-1\\1&-1&1&-1\\1&-1&-1&1
 \end{pmatrix}.
\]
They preserve \(D\), are orthogonal involutions, and act on
\((v,s,c)\) by \((s\ c)\) and \((v\ s)\), respectively,
as calculated in Theorem~\ref{esn:thm:d4-f4-triality}.
Their binary matrices in the retained marking are
\[
 B_R=\begin{pmatrix}1&1\\0&1\end{pmatrix},\qquad
 B_H=\begin{pmatrix}0&1\\1&0\end{pmatrix}.
\]
The same earlier coordinate calculation gives
\(R_D^2=H_D^2=(R_DH_D)^3=I\); thus their subgroup has at most six
elements, by reducing words to \((R_DH_D)^k\) or
\((R_DH_D)^kR_D\), \(k=0,1,2\).
The induced two transpositions generate the six permutations of the
three nonzero binary vectors.  There are exactly six invertible binary
matrices, by the \(3\cdot2\) choices of their independent columns,
and their action on the three nonzero vectors is faithful.  Hence the
six words have distinct discriminant actions, and the restriction of
\(d_D\) to \(\langle R_D,H_D\rangle\) is an isomorphism onto
\(GL_2(\mathbb F_2)\).  Its inverse is the specified \(\delta\).
For every \(T\in O(D)\),
\(T\delta(d_D(T))^{-1}\in\ker d_D=W_D\).
The two factors intersect trivially because \(d_D\delta=I\).
This proves the asserted splitting and exhausts \(O(D)\).
The earlier theorem proved
\(W_F=\langle W_D,R_D,H_D\rangle\), so it also proves the last equality.
\end{proof}

\begin{theorem}[The exact full automorphism quotient of \(N\)]
\label{esn:fa:thm:full-quotient}
Every isometry of \(N\) has a unique factorization
\[
 T=wL_g,\qquad w\in W(R_{\mathrm{ES}}),\quad g\in G.
\]
The Weyl group is normal, and the induced discriminant action gives
the split exact sequence and the indicated isomorphisms
\begin{equation}
 \begin{split}
 1&\longrightarrow W(R_{\mathrm{ES}})
   \longrightarrow O(N)\xrightarrow{\ d_{\mathrm{ES}}\ }G
   \longrightarrow1,\\
 O(N)&=W(R_{\mathrm{ES}})\rtimes L(G),\\
 L(G)&\xrightarrow{\ \cong\ }O(N)/W(R_{\mathrm{ES}})
               \cong G\xrightarrow{\ \Theta\ }GL_2(\mathbb F_3),\\
 W(R_{\mathrm{ES}})&\cong S_6^4\times W_D,
 \qquad |O(N)|=48\cdot(6!)^4\cdot(2^3\cdot4!).
 \end{split}
 \label{esn:fa:eq:full-group}
\end{equation}
In identified coordinates, if
\(w=(Q_{\sigma_1},Q_{\sigma_2},Q_{\sigma_3},Q_{\sigma_4};w_D)\)
and \(g=(\varepsilon,p,B)\), the exact semidirect-product action is
\begin{equation}
 L_gwL_g^{-1}
   =\bigl(Q_{\sigma_{p(1)}},Q_{\sigma_{p(2)}},
          Q_{\sigma_{p(3)}},Q_{\sigma_{p(4)}};
          \delta(B)w_D\delta(B)^{-1}\bigr).
 \label{esn:fa:eq:weyl-conjugation}
\end{equation}
In the original fibre coordinates the \(i\)-th permutation factor
on the right means
\(\varphi_iQ_{\sigma_{p(i)}}\varphi_i^{-1}\).
\end{theorem}

\begin{proof}
Theorem~\ref{esn:thm:niemeier} proves that the norm-two vectors of \(N\)
are exactly the 144 roots of its five original orthogonal summands.
Join two distinct roots by an edge when their inner product is nonzero.
The five summands are exactly the connected components of this graph.
There are no edges between different summands by orthogonality.
Within an \(A_5\) summand, the simple roots
\(e_0-e_1,e_1-e_2,e_2-e_3,e_3-e_4,e_4-e_5\)
have consecutive inner product \(-1\), forming a connected graph,
and span its real space.  Every root not already among them has nonzero
inner product with at least one of them; otherwise it is orthogonal
to its entire real space and has norm zero, contrary to norm two.
Thus its whole root graph is connected.  In \(D\), the four roots
\[
 e_1-e_2,\quad e_2-e_3,\quad e_3-e_4,\quad e_3+e_4
\]
span \(\mathbb R D\), and the second has inner product \(-1\)
with each of the other three.  The same spanning argument proves
connectedness of the \(D_4\) root graph.

An element \(T\in O(N)\) permutes the norm-two vectors and preserves
these edges, hence permutes their connected components.  Their
cardinalities are \(30,30,30,30,24\), so \(T\) preserves the
\(D_4\) component and permutes the four \(A_5\) components.
The roots generate each original lattice summand: the five and four
displayed simple roots are their integral bases by
Theorem~\ref{esn:thm:root-discriminant}.
Consequently \(T\) preserves their integer span
\(R_{\mathrm{ES}}\).  The same component argument applies to every
element of \(O(R_{\mathrm{ES}})\), since its norm-two vectors are
the same 144 roots.  It follows that every such element has, in the
fixed identified coordinates, the form
\begin{equation}
 T(X_1,X_2,X_3,X_4;d)
    =\bigl(\varepsilon_iQ_{\sigma_i}X_{p(i)}\bigr)_{i=1}^4
                                  \mathbin{;}T_Dd,
 \label{esn:fa:eq:all-root-isometries}
\end{equation}
where \(p\in S_4\), \(\varepsilon_i\in\{1,-1\}\),
\(\sigma_i\in S_6\), and \(T_D\in O(D)\).
This follows by transporting each isometry from its input \(A_5\)
summand to its output summand using the fixed \(\varphi_i\), then
applying Theorem~\ref{esn:fa:thm:a5-automorphisms}.
All choices in \eqref{esn:fa:eq:all-root-isometries} are unique.
The convention is precisely the earlier convention: a vector supported
on component \(k\) is sent to component \(p^{-1}(k)\).

Let \(\mathcal A_{\mathrm{ES}}\) be the group of transformations
of the exact discriminant group
\((\mathbb Z/6)^4\oplus\mathbb F_2^2\) given by
\[
 (x;y)\longmapsto
        \bigl((\varepsilon_i x_{p(i)})_{i=1}^4;By\bigr),
 \qquad \varepsilon_i\in\{1,-1\},\quad p\in S_4,
                 \quad B\in GL_2(\mathbb F_2).
\]
Each tuple gives a distinct transformation: the image of the order-six
generator in each individual summand determines its output summand
and its sign, and the images of the two binary basis vectors determine
\(B\).  The discriminant action of
\eqref{esn:fa:eq:all-root-isometries}, calculated by the two preceding
theorems, is exactly this tuple with \(B=d_D(T_D)\).
This defines a homomorphism
\[
 d_{\mathrm{ES}}:O(R_{\mathrm{ES}})\longrightarrow
                                      \mathcal A_{\mathrm{ES}}.
\]
It is a homomorphism because extending an isometry to the dual and then
passing to the quotient commutes with composition.  In particular the
component permutation rule remains
\(P_pP_q=P_{q\circ p}\); there is no change of the earlier convention.

This homomorphism is onto.  Given any tuple
\(a=(\varepsilon,p,B)\in\mathcal A_{\mathrm{ES}}\), define
\[
 \widetilde L_a(X_1,X_2,X_3,X_4;d)
   =\bigl(\varepsilon_iX_{p(i)}\bigr)_{i=1}^4
                                     \mathbin{;}\delta(B)d.
\]
It is a root-lattice isometry and has discriminant action \(a\).
The map \(a\mapsto\widetilde L_a\) is a homomorphism: if
\(a=(\varepsilon,p,B)\) and \(a'=(\varepsilon',q,B')\),
its composite has \(i\)-th output
\(\varepsilon_i\varepsilon'_{p(i)}X_{q(p(i))}\)
and \(D\)-component \(\delta(B)\delta(B')d=\delta(BB')d\),
which are exactly the tuple product in the stated convention.

The kernel of \(d_{\mathrm{ES}}\) consists exactly of the tuples
in \eqref{esn:fa:eq:all-root-isometries} with
\(p=1\), every \(\varepsilon_i=1\), and \(T_D\in W_D\).
It is therefore the direct product \(S_6^4\times W_D\).
Every root reflection acts on its own orthogonal summand and fixes
the other four pointwise, and the preceding two theorems identify the
groups generated inside those summands.  Hence this direct product
is exactly \(W(R_{\mathrm{ES}})\).

Every root reflection also preserves \(N\).  Indeed \(\alpha\in N\)
and integrality of \(N\) give
\(s_\alpha x=x-(x,\alpha)\alpha\in N\) for every \(x\in N\).
The reflection is an involution, so inclusion is equality.
Thus \(W(R_{\mathrm{ES}})\subset O(N)\).
It is normal there, since for \(T\in O(N)\)
\[
 Ts_\alpha T^{-1}(x)
   =x-(x,T\alpha)T\alpha=s_{T\alpha}(x),
\]
and \(T\alpha\) is again a norm-two root.

It remains to impose the actual glue, which completely determines the
remaining isometries.  For \(T\in O(R_{\mathrm{ES}})\) and
\(x\in R_{\mathrm{ES}}^\vee\), the inverse-image definition gives
\[
 Tx\in N
 \quad\Longleftrightarrow\quad
 d_{\mathrm{ES}}(T)(x+R_{\mathrm{ES}})\in\mathcal H.
\]
Consequently \(T(N)=N\) if and only if
\(d_{\mathrm{ES}}(T)(\mathcal H)=\mathcal H\).
For the forward implication, represent each class of \(\mathcal H\)
by a vector of \(N\) and apply both \(T\) and \(T^{-1}\).
For the reverse implication, the displayed equivalence applies to
\(T\) and \(T^{-1}\) because the action on \(\mathcal H\) is
bijective.  The group preserving \(\mathcal H\) inside precisely
\(\mathcal A_{\mathrm{ES}}\) is the group \(G\) of
Theorem~\ref{esn:thm:glue-gl2}, proved there to be exactly
\(GL_2(\mathbb F_3)\) through the code map \(\Theta\).
We have therefore proved the equality of subgroups
\[
 O(N)=d_{\mathrm{ES}}^{-1}(G)
                          \quad\text{inside }O(R_{\mathrm{ES}}).
\]
The section \(\widetilde L\) restricts on \(G\) to the already
specified lift \(L\).  This proves surjectivity of the restricted
map \(d_{\mathrm{ES}}:O(N)\to G\), with exactly the asserted kernel.
For \(T\in O(N)\), set \(g=d_{\mathrm{ES}}(T)\).
Then \(w=TL_g^{-1}\) lies in the kernel and \(T=wL_g\).
If \(wL_g=w'L_{g'}\), applying \(d_{\mathrm{ES}}\) gives
\(g=g'\), then cancellation gives \(w=w'\).
In particular \(L(G)\cap W(R_{\mathrm{ES}})=\{I\}\).
Normality and this unique factorization prove the semidirect product
and the quotient isomorphism, including its surjectivity.
The four \(S_6\) factors have order \((6!)^4\);
\(W_D\) has \(2^3\) sign choices and \(4!\) permutation choices,
and \(|G|=48\).  Unique factorization proves the displayed order.

Finally put \(Z=L_g^{-1}X\).  From the \(i\)-th output equation
\(X_i=\varepsilon_iZ_{p(i)}\) follows
\(Z_{p(i)}=\varepsilon_iX_i\), with each vector read in the fixed
identified \(A\) coordinates.  Therefore the \(i\)-th output of
\(L_gwL_g^{-1}X\) is
\[
 \varepsilon_iQ_{\sigma_{p(i)}}Z_{p(i)}
   =\varepsilon_iQ_{\sigma_{p(i)}}\varepsilon_iX_i
   =Q_{\sigma_{p(i)}}X_i.
\]
The \(D\)-component is
\(\delta(B)w_D\delta(B)^{-1}d\).
It belongs to \(W_D\) by normality of \(\ker d_D\).
These are exactly the components of
\eqref{esn:fa:eq:weyl-conjugation}, proving the asserted action with
all signs and component indices retained.
\end{proof}

\begin{theorem}[The residue six-cycle in the full automorphism group]
\label{esn:fa:thm:residue-weyl}
For \(r\in S_{840}=\langle289\rangle\), let
\(\rho(r)e_s=e_{rs}\) on each of the four fibres and
\(\rho(r)|_D=I\), with the exact fibre action used in
\eqref{esn:eq:residue-action-cocycle}.
Then
\begin{equation}
 \rho(S_{840})\subset W(R_{\mathrm{ES}}),\qquad
 d_{\mathrm{ES}}(\rho(r))=1,\qquad
 L_g\rho(r)L_g^{-1}=\rho(r)
       \quad(g\in G,\ r\in S_{840}).
 \label{esn:fa:eq:residue-position}
\end{equation}
In particular the map
\[
 S_{840}\times G\longrightarrow O(N),\qquad(r,g)\longmapsto\rho(r)L_g
\]
is an injective homomorphism with image of order \(6\cdot48=288\),
isomorphic to \(C_6\times GL_2(\mathbb F_3)\).
\end{theorem}

\begin{proof}
Put \(C e_j=e_{j+1\bmod6}\) on the free six-coordinate space and
restrict it to \(\mathbb R A\).  If
\(r=289^m\), with \(m\in\{0,1,2,3,4,5\}\), the exact fibre
identification gives
\[
 \rho(r)e_{f_i289^j}=e_{f_i289^{j+m}},\qquad
 \rho(r)|_{A(F_{\eta_i})}=\varphi_iC^m\varphi_i^{-1}.
\]
Thus in all four identified summands \(\rho(r)\) is the same
coordinate permutation \(C^m\), and it is the identity on \(D\).
Theorem~\ref{esn:fa:thm:a5-automorphisms} puts it in the corresponding
product of Weyl groups, giving the first two assertions.
Each sign \(\varepsilon_i I\) commutes with \(C^m\), and the
same operator \(C^m\) occurs in every component.  Directly the
\(i\)-th outputs of \(L_g\rho(r)\) and \(\rho(r)L_g\) are
respectively
\(\varepsilon_iC^mX_{p(i)}\) and
\(C^m\varepsilon_iX_{p(i)}\), which agree; the \(D\)-outputs
also agree since \(\rho(r)|_D=I\).
This proves the conjugation identity without changing the chosen origins.

The six restrictions \(C^m|_{\mathbb R A}\) are distinct: if
\(C^m(e_0-e_1)=e_0-e_1\), comparison of the unique coordinate with
coefficient \(+1\) gives \(m=0\pmod6\).
Thus \(\rho\) is faithful.  Commutation makes the displayed product
map a homomorphism.  If \(\rho(r)L_g=I\), then
\(L_g=\rho(r)^{-1}\) lies in
\(L(G)\cap W(R_{\mathrm{ES}})=\{I\}\).
Faithfulness of both factors gives \(r=1\) and \(g=1\).
This proves injectivity and the exact image order.
\end{proof}

% Exact proof module, 6 September 2026. Input after es_niemeier_shared.tex.
% All original coordinates and Euclidean root lengths are retained.
\section{The exact Leech neighbor of the retained index-72 lattice}
\label{esnl:sec:neighbor}

Retain the lattice, metric, discriminant marking, and glue literally:
\[
 \begin{split}
 A&=\{(x_0,\ldots,x_5)\in\mathbb Z^6:\sum x_j=0\},\\
 D&=\{(d_1,\ldots,d_4)\in\mathbb Z^4:\sum d_i\equiv0\pmod2\},
 \qquad R=A^{\perp4}\perp D,\\
 C_2&=\langle b_0=(1111;00),b_1=(1100;10),b_2=(1010;01)\rangle_{\mathbb F_2},\\
 C_3&=\langle z_1=(1,1,1,0),z_2=(1,2,0,1)\rangle_{\mathbb F_3},\\
 H&=\{(3u+2z;y):(u;y)\in C_2,\ z\in C_3\},\\
 N&=\{x\in R^\vee:x+R\in H\}.
 \end{split}
\]
The coordinates in each $A$ remain the six ordered elements $f_i289^j$
of its original residue fibre. The $A$ discriminant generator is
$\lambda=(5,-1,-1,-1,-1,-1)/6$. The $D$ markings $10,01,11$
remain represented by $v_D=(1,0,0,0)$,
$s_D=(1,1,1,1)/2$, $c_D=(-1,1,1,1)/2$.
The preceding proof establishes that $N$ is even unimodular,
$[N:R]=72$, $\det R=5184$, and its roots are exactly those of
$A_5^4D_4$. Every root has squared norm two. No metric rescaling
occurs anywhere below.

\subsection{The Weyl vector and the exact neighbor}

Use the original simple roots $a_j=e_j-e_{j+1}$, $0\le j\le4$,
in each $A$ component and
\[
 d_1=e_1-e_2,\quad d_2=e_2-e_3,\quad
 d_3=e_3-e_4,\quad d_4=e_3+e_4
\]
in $D$. Their highest roots are
$\theta_A=e_0-e_5=\sum_ja_j$ and
$\theta_D=e_1+e_2=d_1+2d_2+d_3+d_4$.
Both highest-root heights are five, so the common Coxeter number
is six. Define
\[
 \rho_A=(5,3,1,-1,-3,-5)/2,\qquad
 \rho_D=(3,2,1,0),\qquad \rho=\rho_A^{\oplus4}\perp\rho_D.
 \label{esnl:eq:rho}
\]
Each simple root pairs to one with $\rho$.
Moreover $\rho_A-3\lambda=(0,2,1,0,-1,-2)\in A$ and
$\rho_D\in D$. Thus the exact glue class of $\rho$ is
$(3,3,3,3;00)$, the image of $b_0$, and $\rho\in N$.
Directly
\[
 \rho_A^2=35/2,\quad \rho_D^2=14,\quad
 \rho^2=4(35/2)+14=84=2\cdot6\cdot7.
\]
Consequently $(\rho,N)\subset\mathbb Z$.
Fix $\alpha=a_0$ in the first $A$ component and put
\[
 \ell_\rho:N\to\mathbb Z/6,\quad x\mapsto(\rho,x)\bmod6,\qquad
 K=\ker\ell_\rho,\quad R_0=R\cap K,\quad
 v=\rho/6-\alpha,\quad \Lambda=K+\mathbb Z v.
 \tag{NL1}
\]
The pairing $(\rho,\alpha)=1$ proves
$[N:K]=[R:R_0]=6$. For $x\in K$, the number
$(v,x)=(\rho,x)/6-(\alpha,x)$ is integral, and
\[
 v^2=84/36-2/6+2=4,\qquad
 6v=\rho-6\alpha\in K,\qquad(\rho,6v)=78.
\]
Hence $\Lambda$ is positive definite even integral of rank 24.
If $mv\in N$, its pairing after adding $m\alpha$ with $\alpha$
is $m/6\in\mathbb Z$, so $6\mid m$.
It follows that
\[
 N\cap\Lambda=K,\qquad [\Lambda:K]=6,\qquad
 \det K=36,\qquad \det\Lambda=36/6^2=1.
 \tag{NL2}
\]
Changing $\alpha$ to another root of $\rho$-pairing one changes
$v$ by an element of $K$, and so does not change $\Lambda$.
We shall use the equivalent exact congruence description
\[
 \Lambda=\{x+k\rho/6:x\in N,\ k\in\mathbb Z,\
                                (\rho,x)+k\equiv0\pmod6\}.
 \tag{NL3}
\]
Indeed write $x+k\rho/6=(x+k\alpha)+kv$; the first summand
is in $K$ exactly under this congruence. Changing $k$ by six
can be absorbed in $x$ using $\rho\in N$, and preserves the
congruence since $\rho^2=84$.

\subsection{A complete rootlessness proof}

\begin{theorem}[The Leech neighbor in the retained coordinates]
\label{esnl:thm:leech-neighbor}
The lattice $\Lambda=K+\mathbb Z(\rho/6-\alpha)$ is positive
definite even unimodular of rank 24 and has minimum squared norm
exactly four. It is the holy-construction Leech lattice of the
retained $A_5^4D_4$ glue, and
$N\cap\Lambda=K$ with $[N:K]=[\Lambda:K]=6$.
\end{theorem}

\begin{proof}
The rank, parity, unimodularity and indices were proved in
(NL1)--(NL2). The holy-presentation identification is proved in
(NL8)--(NL11) below. We now prove the required minimum directly.
For $k=0$ in (NL3), a possible root is an original root of $N$.
An $A$ root pairs with $\rho$ to one of $\pm1,\ldots,\pm5$.
For $D$, all roots $\pm e_i\pm e_j$ have nonzero $\rho_D$-pairing
of absolute value at most five. Thus none belongs to $K$.
Negation and the reduction of $k$ modulo six leave $k=1,2,3$.

Define the exact shifted class minima by
\[
 m_A(k,t)=\min_{x\in t\lambda+A}|x+k\rho_A/6|^2,\qquad
 m_D(k,y)=\min_{x+D=y}|x+k\rho_D/6|^2.
\]
Their values are
\[
 \begin{array}{c|rrrrrr}
 k\backslash t&0&1&2&3&4&5\\\hline
 1&35/72&35/72&35/72&35/72&35/72&35/72\\
 2&11/18&4/9&11/18&4/9&11/18&4/9\\
 3&3/8&17/24&17/24&3/8&17/24&17/24
 \end{array}
 \quad
 \begin{array}{c|rrrr}
 k\backslash y&00&10&01&11\\\hline
 1&7/18&7/18&7/18&7/18\\
 2&2/9&5/9&5/9&5/9\\
 3&1/2&1/2&1/2&1/2
 \end{array}.
 \tag{NL4}
\]
We prove the table and its needed equality information.
The shifted $A$ coordinates are precisely
\[
 \left(n_j-\frac t6+\frac{k(5-2j)}{12}\right)_{j=0}^5,
               \qquad n_j\in\mathbb Z,\quad\sum n_j=t.
 \tag{NL5}
\]
At $k=1$, coordinatewise nearest representatives are the six
numbers $\pm1/12,\pm3/12,\pm5/12$ in an order depending on $t$.
They sum to zero and their squares sum to $35/72$.
At $k=2$ and even $t$, the nearest coordinates consist of
two copies of each sign of $1/6$ and one copy of each sign of
$1/2$; the half-integer ties admit exactly the needed zero-sum
choice and the squared sum is $11/18$.
At odd $t$, there are two each of $0,1/3,-1/3$, giving $4/9$.
For $t=1,3,5$ their actual ordered shifted vectors are
\[
 (-1,1,0,-1,1,0)/3,\quad
 (1,0,-1,1,0,-1)/3,\quad
 (0,-1,1,0,-1,1)/3.
 \tag{NL6}
\]
Subtracting $\rho_A/3$ gives original squared norms
$17/6,3/2,17/6$, respectively.
At $k=3,t=0,3$, the nearest coordinates alternate $\pm1/4$,
giving $3/8$ with zero sum.
At each other $t$, the coordinatewise nearest representatives
have three absolute values $1/12$ and three $5/12$, square sum
$13/24$, and coordinate sum $1$ or $-1$.
Meeting the zero-sum constraint requires an integer change.
Every nonzero coordinate change increases the squared sum by
at least $(7/12)^2-(5/12)^2=1/6$.
Changing one suitable $5/12$ coordinate attains this increase
and restores the zero sum. Thus the constrained minimum is
$17/24$. At $t=0,3$, the original minimizing vectors are
\[
 (-1,-1,0,0,1,1),\qquad (-3,-1,-1,1,1,3)/2,
 \tag{NL7}
\]
of squared norms $4,11/2$.

For $D$ and $k=1$ the shift is $(1/2,1/3,1/6,0)$.
The nearest integer coordinate distances are
$(1/2,1/3,1/6,0)$; the first-coordinate tie meets either integer
parity. The nearest half-integer distances are
$(0,1/6,1/3,1/2)$; the last-coordinate tie meets either
half-integer class. Both squared sums are $7/18$.
At $k=2$ the shift is $(1,2/3,1/3,0)$.
The unique nearest integer vector before shifting is
$(-1,-1,0,0)\in D$, giving $2/9$.
Changing parity costs at least $1/3$, attained by changing the
second or third coordinate, giving $5/9$.
The half-integer coordinate distances are
$(1/2,1/6,1/6,1/2)$, giving $5/9$; ties meet either marking.
At $k=3$ the shift is $(3/2,1,1/2,0)$.
Integer-class minimum distances are $(1/2,0,1/2,0)$,
and half-integer-class distances $(0,1/2,0,1/2)$.
Ties meet every marking and all square sums are $1/2$.
In class $00$, the two original minimizing vectors are
\[
 (-1,-1,0,0),\qquad(-2,-1,-1,0),
\]
of squared norms $2,6$. In class $10$, they are
\[
 (-1,-1,-1,0),\qquad(-2,-1,0,0),
\]
of squared norms $3,5$.
In either half-integer class they have the shapes
$(-3/2,-1/2,-1/2,\pm1/2)$ and
$(-3/2,-3/2,-1/2,\pm1/2)$.
Exactly one last-coordinate sign in each shape meets each
original marking; their squared norms are $3,5$.
This proves every entry and the relevant equality cases in (NL4).

For $k=1$, every full shifted vector has squared norm at least
$4(35/72)+7/18=7/3>2$.
For $k=2$, the lower bound from (NL4) is two.
Equality forces all $t_i$ odd and $y=00$.
In the actual glue this forces $u=1111$.
If $z=0$, all $t_i=3$ and (NL6) gives
$x^2=4(3/2)+2=8$.
If $z\ne0$, its weight is three, so exactly one $t_i=3$
and three $t_i\in\{1,5\}$; then
$x^2=3/2+3(17/6)+2=12$.
The norm-two equation is
\[
 2=x^2+\tfrac23(\rho,x)+\tfrac{28}{3},
 \qquad(\rho,x)=-11-\tfrac32x^2.
\]
Its pairing is $-23$ or $-29$, always $1$ modulo six.
Condition (NL3) requires residue $4$, so all nine equality
candidates, one per ternary word, are excluded.

For $k=3$, the lower bound two requires $z=0$:
every nonzero ternary word raises the minimum by
$3(17/24-3/8)=1$.
If $u=0$, the possible original squared norms are $18,22$.
If $u=1111$, they are $24,28$.
If $u$ has weight two, its $D$ class is nonzero and they are
$2(4)+2(11/2)+3=22$ or $2(4)+2(11/2)+5=24$.
These are all sixteen equality candidates.
But
\[
 2=x^2+(\rho,x)+21,\qquad(\rho,x)=-19-x^2,
\]
and (NL3) would require $x^2\equiv2\pmod6$.
None of $18,22,24,28$ has that residue.
Thus $\Lambda$ has no roots. Together with (NL2), this proves
that it is positive definite even unimodular of rank 24,
with minimum squared norm exactly four, attained by $v$.
\end{proof}

\subsection{The exact holy-construction presentation}

For $0\le t\le5$ set
\[
 \omega_t=(\underbrace{1,\ldots,1}_{t},
             \underbrace{0,\ldots,0}_{6-t})-\tfrac t6(1,\ldots,1).
\]
Then $\omega_t+A=t\lambda+A$ and
$(\rho_A,\omega_t)=t(6-t)/2=3\omega_t^2$.
For the $D$ markings set
\[
 \omega_{00}=0,\quad\omega_{10}=e_1,\quad
 \omega_{01}=(1,1,1,1)/2,\quad
 \omega_{11}=(1,1,1,-1)/2.
\]
The last differs from the original $c_D$ by $(1,0,0,-1)\in D$,
so retains its exact marking.
For $a=(t_1,t_2,t_3,t_4;y)\in H$ define
\[
 u_a=(\omega_{t_1},\omega_{t_2},\omega_{t_3},\omega_{t_4};\omega_y),
 \qquad g_a=u_a-\rho/6.
 \tag{NL8}
\]
These give all 72 representatives of $N/R$. Each nonzero $D$
representative has norm one and pairing three; hence
\[
 (\rho,u_a)=3u_a^2\in6\mathbb Z,\quad
 u_a\in K,\quad g_a^2=u_a^2-(\rho,u_a)/3+84/36=7/3.
 \tag{NL9}
\]
Let $\mathcal F$ be the 29 vectors consisting, in each original
component, of the negative original simple roots and the positive
highest root. Their pairings with $\rho$ are $-1$ or $5$.
There are positive integral zero relations of coefficient sum six:
\[
 \theta_A+\sum_{j=0}^4(-a_j)=0,\qquad
 \theta_D+(-d_1)+2(-d_2)+(-d_3)+(-d_4)=0.
 \tag{NL10}
\]
Every component of $g_a$ pairs to $1/6$ with the corresponding
vertices except one tip, where the pairing is $-5/6$.
For $A$ this follows because $\omega_t$ has simple-root pairing
one only at $a_{t-1}$ when $t\ne0$, and has
$(\omega_t,\theta_A)=1$ for $t\ne0$.
For $D$ the three nonzero representatives have simple-root
pairing one only at $d_1,d_4,d_3$, respectively; all have
highest-root pairing one. The zero representative points to
the highest root. These computations fix all orientations in
the holy generator diagram.

The following coefficient presentations are exact equalities:
\[
 \begin{split}
 N&=\left\{\sum_fa_ff+\sum_ab_ag_a:
                    a_f,b_a\in\mathbb Z,\ \sum_ab_a=0\right\},\\
 \Lambda&=\left\{\sum_fa_ff+\sum_ab_ag_a:
            a_f,b_a\in\mathbb Z,\ \sum_fa_f+\sum_ab_a=0\right\}.
 \end{split}
 \tag{NL11}
\]
For the first, the $f$ span $R$ and $g_a-g_0=u_a$ generate its
given extension $N$; the zero coefficient sum cancels $\rho/6$.
For the second, write $A=\sum a_f$, $B=\sum b_a$,
$r=\sum a_ff$, and $x=r+\sum b_au_a$.
The vector is $x-B\rho/6$, and by (NL9)
$(\rho,x)\equiv-A\equiv B$, exactly (NL3) with $k=-B$.
Conversely, for any $x+k\rho/6$ in (NL3), write
$x=r+\sum b_au_a$. The coefficient of $u_0=0$ can be set so
$B=-k$. Write $r$ with the negative simple-root vertices.
The congruence gives $A+B\equiv0\pmod6$.
Adding a suitable multiple of a zero relation (NL10) makes
$A+B=0$ exactly. Thus both directions of (NL11) are proved.

These are exactly the holy coefficient rules of Conway--Sloane
\cite[pp.~277--280]{esn:ConwaySloane1982}.
Borcherds gives the uniform proof of the holy-construction
theorem \cite[Section~7, especially the coefficient rules
(i)--(ii) and Lemma~7.3]{esn:Borcherds1985}.
Thus (NL11) is an explicit realization of that Leech
construction from the actual 72 marked glue words; all its
metric, determinant, integrality and rootlessness assertions
have also been independently proved here.

\subsection{The indices, discriminant morphism, and cyclic order twelve}

Reduction modulo $R$ gives an isomorphism $K/R_0\cong H$:
its kernel is $R_0$, and (NL9) supplies representatives in $K$
for every $H$ class. Consequently
\[
 [K:R_0]=72,\quad [N:R_0]=[\Lambda:R_0]=432,\quad
 \det R_0=6^2\cdot5184=186624.
 \tag{NL12}
\]
Thus the original and new inclusions, with every index, are
\[
 R_0\subset R\subset N,\quad [R:R_0]=6,\ [N:R]=72;
 \qquad
 R_0\subset K\subset\Lambda,\quad[K:R_0]=72,\ [\Lambda:K]=6.
\]
In particular $R\cap\Lambda=R_0$, an exact common sublattice.

\begin{theorem}[The exact order-twelve quotient generator]
\label{esnl:thm:cyclic-twelve}
The quotient of $\Lambda$ by the original height-zero root
sublattice $R_0$ is
$\Lambda/R_0\cong\mathbb Z/12\oplus(\mathbb Z/6)^2$.
The class of $v=\rho/6-\alpha$ has exact order twelve.
\end{theorem}

\begin{proof}
Let $\bar b_i$ denote the image in $H$ of the binary generator
$b_i$, namely its first four coordinates multiplied by three
and its last two coordinates retained.
Since $6v=\rho-6\alpha$ and $6\alpha\in R_0$,
\[
 6[v]=\bar b_0\quad\text{in }\Lambda/R_0.
 \tag{NL13}
\]
This is a nonzero order-two element. The classes
$\bar b_0,\bar b_1,\bar b_2,2z_1,2z_2$ form the displayed
independent binary and ternary generators of $H$.
Eliminate $\bar b_0$ using (NL13). The resulting surjective
presentation is
\[
 \Lambda/R_0\cong
 \mathbb Z/12\oplus(\mathbb Z/2)^2\oplus(\mathbb Z/3)^2
 \cong\mathbb Z/12\oplus(\mathbb Z/6)^2.
 \tag{NL14}
\]
There are no further relations: this presentation has order
$12\cdot4\cdot9=432$, equal to (NL12).
Thus $[v]$ has exact order twelve. This particular occurrence
of twelve is proved from the Coxeter denominator six and the
actual order-two Weyl-vector glue class.
\end{proof}

Set $M=N+\Lambda=N+\mathbb Z(\rho/6)$.
Then
\[
 M=K^\vee,\qquad K^\vee/K=
    \langle[\alpha],[v]\rangle\cong(\mathbb Z/6)^2.
 \tag{NL15}
\]
Indeed $M\subset K^\vee$, its index over $N$ is six by pairing
with $\alpha$, and $[K^\vee:N]=[N:K]=6$.
The cyclic images of $N/K$ and $\Lambda/K$ have order six and
trivial intersection, proving the stated generators.
Their exact discriminant data are
\[
 q_K([\alpha])=q_K([v])=0,\quad
 b_K([\alpha],[v])=1/6\pmod{\mathbb Z},\quad
 q_K(a[\alpha]+b[v])=ab/3\pmod{2\mathbb Z}.
 \tag{NL16}
\]
These follow from $\alpha^2=2$, $v^2=4$ and
$(\alpha,v)=1/6-2=-11/6$.
The two original lattices are therefore the inverse images
of the two indicated order-six isotropic subgroups in this
specified discriminant form.

\subsection{An exact Lorentzian quotient and symmetry transport}

Take $\mathcal I=N\perp U$ in coordinates $(x,m,n)$ with
\[
 \langle(x,m,n),(x',m',n')\rangle
 =(x,x')-mn'-m'n .
\]
Thus $U$ has Gram matrix
$\left(\begin{smallmatrix}0&-1\\-1&0\end{smallmatrix}\right)$.
The vectors $z=(0,0,1)$ and $w=(\rho,6,7)$ are primitive
isotropic: $w^2=84-2\cdot6\cdot7=0$, and six and seven are
coprime. There are integral quotient isometries
\[
 \begin{split}
 z^\perp/\mathbb Zz&\longrightarrow N,\quad[(x,0,n)]\mapsto x,\\
 w^\perp/\mathbb Zw&\longrightarrow\Lambda,\quad
                                  [(x,m,n)]\mapsto x-m\rho/6.
 \end{split}
 \tag{NL17}
\]
For the second map, orthogonality means $(\rho,x)=7m+6n$,
giving (NL3) with $k=-m$. Conversely take $m=-k$ and
$n=((\rho,x)+7k)/6$ for a vector in (NL3).
If the image vanishes, $x=m\rho/6$; pairing with $\alpha$
forces $m=6j$, and then $n=7j$, so the kernel is exactly
$\mathbb Zw$. Finally
\[
 |x-m\rho/6|^2
 =x^2-\tfrac m3(\rho,x)+\tfrac73m^2=x^2-2mn.
\]
Polarization gives the full pairing, proving the isometry.

For every $T\in O(N)$ the exact transport laws are
\[
 T(K_\rho)=K_{T\rho},\qquad
 T(\Lambda_\rho)=\Lambda_{T\rho},\qquad
 (x,m,n)\mapsto(Tx,m,n):w_\rho\mapsto w_{T\rho}.
 \tag{NL18}
\]
They follow by substitution in (NL3).
In particular the stabilizer of $\rho$ acts on the fixed
$\Lambda_\rho$. General isometries act on this specified
family of neighbors; no assertion that an arbitrary chosen
automorphism complement fixes this Weyl vector is required.

\subsection{The exact $F_4$ Coxeter action of order twelve}

The order twelve in (NL14) is a lattice quotient invariant of
this explicitly marked construction. There is also the ambient
$F_4$ Coxeter action on the same original $D_4$ real root space.
Use the precise $F_4$ simple roots
\[
 \beta_1=e_2-e_3,\quad\beta_2=e_3-e_4,\quad
 \beta_3=e_4,\quad\beta_4=(e_1-e_2-e_3-e_4)/2,
 \qquad C=s_{\beta_1}s_{\beta_2}s_{\beta_3}s_{\beta_4},
\]
where $s_\beta(x)=x-2(x,\beta)\beta/\beta^2$ and the rightmost
factor acts first. The first two roots have squared norm two,
and the last two squared norm one. No rescaling of the $D_4$
roots has occurred.

\begin{theorem}[Coxeter twelve and its marked triality quotient]
\label{esnl:thm:f4-coxeter-twelve}
In the retained four-coordinate metric,
\[
 C=\frac12
 \begin{pmatrix}
 1&1&1&1\\-1&1&1&-1\\1&1&-1&-1\\1&-1&1&-1
 \end{pmatrix},\qquad
 C^3=
 \begin{pmatrix}
 0&1&0&0\\-1&0&0&0\\0&0&0&-1\\0&0&1&0
 \end{pmatrix}.
 \tag{NL19}
\]
It satisfies $C^4-C^2+I=0$, $C^6=-I$ and has exact order
twelve. Its marked discriminant action is
$v_D+D\mapsto c_D+D\mapsto s_D+D\mapsto v_D+D$.
The restriction of the original triality quotient is the exact
sequence
\[
 1\longrightarrow\langle C^3\rangle\cong\mathbb Z/4
 \longrightarrow\langle C\rangle\cong\mathbb Z/12
 \longrightarrow\langle(v_D\,c_D\,s_D)\rangle\cong\mathbb Z/3
 \longrightarrow1,
 \tag{NL20}
\]
with $\langle C^3\rangle\subset W(D_4)$.
\end{theorem}

\begin{proof}
The first reflection exchanges coordinates two and three; the
second exchanges coordinates three and four; the third changes
the sign of coordinate four; the last is
$I-2\beta_4\beta_4^T$. Their product gives the first matrix in
(NL19). Multiplying it three times gives the second, whose
square is $-I$. Multiplying the first matrix also gives
$C^4-C^2+I=0$, and hence $C^{12}=I$.
Its first column gives
$Cv_D=c_D+(1,-1,0,0)$, with the displayed correction in $D$.
Directly $Cs_D=v_D$ and $Cc_D=(0,0,-1,-1)+s_D$.
The displayed integer correction lies in $D$ and proves the
stated three-cycle with the original $v_D,s_D,c_D$ markings.
The first two reflections preserve $D$ as root reflections.
The last two preserve $D$ because a sign change preserves
integer sum parity, and reflection in $\beta_4$ sends an
integer even-sum vector to an integer vector of even sum.
The inverse reflections have the same property, so these
maps are isometries of $D$.
Consequently their induced action on $D^\vee/D$ is defined
and has the exact order-three image calculated above.
The matrix $C^3$ is an even signed permutation of exact order
four, and thus lies in the already identified group $W(D_4)$.
If $C^m=I$, its discriminant image forces $3\mid m$,
and the exact order of $C^3$ then forces $4\mid m/3$.
Thus $C$ has exact order twelve and the discriminant kernel
on its cyclic subgroup is exactly $\langle C^3\rangle$.
This proves every map and group order in (NL20).
\end{proof}

In particular $C^2$ still has nontrivial order-three
discriminant action, while $C^3$ is in the long-root Weyl
group. The Coxeter denominator six in (NL1), the cyclic order
twelve in (NL14), and the $F_4$ Coxeter twelve in (NL20)
therefore have their stated, distinct domains and explicit
relationships through the original $D_4$ marking.

% Staged only. No canonical file has been edited.
% Requires amsmath, amssymb, amsthm, and the existing esn theorem labels.
% All new labels have namespace n24:.

\subsection{An exact cyclic-triality isometry of the entire Niemeier space}
\label{n24:sec:cyclic-isometry}

Retain the ordered fibres
\(f_i289^j\), \(1\leq i\leq4\), \(0\leq j\leq5\),
the Euclidean augmentation coordinates of \(A(F_{\eta_i})\), the
ordered orthonormal coordinates of \(D\subset\mathbb R^4\), and the
72-element glue \(\mathcal H\) of
\eqref{esn:eq:glue-and-overlattice}.  Thus the actual real span is
\[
 V_N=N\otimes_{\mathbb Z}\mathbb R
 =\left\{(x_{i,j};d_k)\in\mathbb R^{24}\oplus\mathbb R^4:
               \sum_{j=0}^5x_{i,j}=0\ (1\leq i\leq4)\right\},
\]
with squared norm \(\sum_{i,j}x_{i,j}^2+\sum_kd_k^2\).
The four displayed linear constraints, rather than an implicit change of
metric, give its dimension 24.  Choose an ordered orthonormal octonion
basis \((u_0,\ldots,u_7)\) with \(u_0=1\).  Every octonion product
below is the product of the chosen real octonion algebra, without a
change of multiplication convention.

For \(p\in\{0,1\}\) and \(j\in\{0,1,2\}\), set
\begin{align}
 m_{i,p}&=\frac13\sum_{j=0}^2x_{i,p+2j},&
 t_{i,p,j}&=x_{i,p+2j}-m_{i,p},\label{n24:eq:triplets}\\
 \alpha_i&=\frac{\sum_{j=0}^2x_{i,2j}
                    -\sum_{j=0}^2x_{i,1+2j}}{\sqrt6},&
 (b_0,\ldots,b_7)&=(\alpha_1,\alpha_2,\alpha_3,\alpha_4,
                           d_1,d_2,d_3,d_4).
 \label{n24:eq:fixed-coordinates}
\end{align}
In particular \(m_{i,0}=\alpha_i/\sqrt6\),
\(m_{i,1}=-\alpha_i/\sqrt6\), and
\(\sum_jt_{i,p,j}=0\).  The ordered map is
\begin{equation}
 \begin{split}
 \mathcal L:V_N&\longrightarrow\mathbb O^3,\\
 (x;d)&\longmapsto(q_1,q_2,q_3),\\
 q_{j+1}&=\sum_{i=1}^4\sum_{p=0}^1
       \left(t_{i,p,j}+\frac{b_{2(i-1)+p}}{\sqrt3}\right)
                                          u_{2(i-1)+p}.
 \end{split}
 \label{n24:eq:full-isometry}
\end{equation}

\begin{theorem}[Actual rank-24 map, inverse, pairing, and residue action]
\label{n24:thm:full-isometry}
The map \(\mathcal L\) is a real linear isometry onto
\(\mathbb O^3\) with its product norm.  Write
\(q_{j+1}=\sum_{a=0}^7q_{j,a}u_a\), and define
\begin{equation}
 b_a=\frac1{\sqrt3}\sum_{j=0}^2q_{j,a},\qquad
 t_{i,p,j}=q_{j,2(i-1)+p}
                    -\frac13\sum_{h=0}^2q_{h,2(i-1)+p}.
 \label{n24:eq:inverse-first}
\end{equation}
Its inverse is exactly
\begin{equation}
 x_{i,p+2j}=t_{i,p,j}+
                 \frac{(-1)^p b_{i-1}}{\sqrt6},
 \qquad d_k=b_{k+3}.
 \label{n24:eq:inverse-second}
\end{equation}
If \(r=\rho(289^2)=\rho(361)\), with the residue action
of Theorem~\ref{esn:fa:thm:residue-weyl}, and
\(c(q_1,q_2,q_3)=(q_3,q_1,q_2)\), then
\begin{equation}
 \mathcal Lr=c\mathcal L,\qquad r^3=c^3=I,
 \qquad rN=N.
 \label{n24:eq:cyclic-intertwining}
\end{equation}
Thus \(\mathcal L|_N\) embeds the original even unimodular lattice,
with every pairing and glue class retained, into the actual three
octonion sectors and intertwines an actual order-three lattice
isometry with cyclic triality.
\end{theorem}

\begin{proof}
The two means sum to zero because each six-coordinate fibre has total
sum zero.  The equations for \(m_{i,p}\) follow from
\(\alpha_i=6m_{i,0}/\sqrt6\).  Orthogonality to constant triplets gives
\[
 \sum_{j=0}^5x_{i,j}^2
 =\sum_{p=0}^1\sum_{j=0}^2t_{i,p,j}^2
                    +3m_{i,0}^2+3m_{i,1}^2
 =\sum_{p,j}t_{i,p,j}^2+\alpha_i^2.
\]
For each \(a=2(i-1)+p\), the same zero-sum identity gives
\[
 \sum_{j=0}^2\left(t_{i,p,j}+\frac{b_a}{\sqrt3}\right)^2
                      =\sum_{j=0}^2t_{i,p,j}^2+b_a^2.
\]
Summing over all eight values of \(a\) proves
\(\sum_{h=1}^3|q_h|^2=\sum_{i,j}x_{i,j}^2+\sum_kd_k^2\).
Polarization proves equality of all pairings.  Equations
\eqref{n24:eq:inverse-first} recover the means and residuals of each
output triplet.  Equation~\eqref{n24:eq:inverse-second} then recovers
the original means, the original fibre coordinates, and \(d\).
Conversely those formulas applied to arbitrary \(q\) yield
\(\sum_jx_{i,j}=0\) and return the prescribed \(q\).  This proves
bijectivity without a dimension-only argument.

The residue action sends the coordinate vector \(e_{f_i289^h}\)
to \(e_{f_i289^{h+2}}\), so the coefficient in output position
\(p+2j\) is the old coefficient in position \(p+2(j-1)\), with
\(j\) read modulo three.  Both means, all \(\alpha_i\), and all
\(d_k\) are fixed.  Hence the three output octonions become
\((q_3,q_1,q_2)\), proving the intertwining formula.  Each coordinate
permutation acts trivially on \(A^\vee/A\): for
\(\lambda=e_0-\frac16\sum_he_h\), its difference from \(\lambda\)
is an augmentation root.  The \(D\)-coordinates are fixed.  Thus
\(r\) acts trivially on the full discriminant group, preserves
\(\mathcal H\), and preserves its exact inverse image \(N\).
Its order is three, as seen on \(e_0-e_1\) in any fibre.
\end{proof}

\paragraph{The exact octonion-coordinate glue test.}
This embedding has the following fully explicit image, denoted
\(\Lambda_{\mathbb O}=\mathcal L(N)\).  Given \(q\), first recover
\((x;d)\) by \eqref{n24:eq:inverse-first}--\eqref{n24:eq:inverse-second}.
For each fibre require
\(x_{i,h}-x_{i,k}\in\mathbb Z\) for every \(h,k\), and put
\[
 k_i=-6x_{i,0}\pmod6\in\mathbb Z/6\mathbb Z.
\]
This is defined because
\(6x_{i,0}=\sum_h(x_{i,0}-x_{i,h})\in\mathbb Z\).
Require also that the four \(d_k\) are all integral or all
half-integral.  Let \(y(d)\in\mathbb F_2^2\) be the unique class with
\[
 d+D=0+D,\ v+D,\ s+D,\ c+D
 \quad\hbox{marked respectively by }00,10,01,11.
\]
Then the final and exact test is
\begin{equation}
 (k_1,k_2,k_3,k_4;y(d))
   =\bigl(3u+2(a+b,a+2b,a,b);\tau(u)\bigr)
 \label{n24:eq:image-glue}
\end{equation}
for \(u\in\mathbb F_2^4\) of even weight,
\((a,b)\in\mathbb F_3^2\), with the first four coordinates read
modulo six, and \(\tau(u)=(u_2+u_4,u_3+u_4)\).
Indeed the integral-difference conditions are precisely membership in
\(A^\vee\), and the common fractional part equals
\(-k_i/6\pmod{\mathbb Z}\), which is the fractional part of
\(k_i\lambda\).  The half-integral test is precisely membership in
\(D^\vee\), and \eqref{n24:eq:image-glue} is the original
\(\mathcal H\).  This proves both directions of the image
description.  No new lattice or rescaled glue is substituted.

\paragraph{The determinant and the factor two in the Albert metric.}
The matrix convention, also specifying this construction independently
of its host document, is
\[
 Q_{\mathbb O}(\lambda;q)=
 \begin{pmatrix}
 \lambda&q_3&\overline q_2\\
 \overline q_3&\lambda&q_1\\
 q_2&\overline q_1&\lambda
 \end{pmatrix}\in\mathfrak h_3(\mathbb O),\qquad
 X\circ Y=(XY+YX)/2.
\]
Put
\[
 E(q)=Q_{\mathbb O}(0;q_1,q_2,q_3),\qquad
 \mathfrak t(q_1,q_2,q_3)=\operatorname{Re}((q_1q_2)q_3),
 \qquad \mathfrak t_N(X)=\mathfrak t(\mathcal LX).
\]
The original ordered Albert determinant becomes the exact polynomial
\begin{equation}
 \det_{\mathbb O}\bigl(\lambda I+E(\mathcal LX)\bigr)
   =\lambda^3-\lambda(X,X)+2\mathfrak t_N(X).
 \label{n24:eq:niemeier-albert-cubic}
\end{equation}
Moreover
\begin{equation}
 \operatorname{tr}\bigl(E(\mathcal LX)\circ E(\mathcal LY)\bigr)
                  =2(X,Y).
 \label{n24:eq:albert-metric-two}
\end{equation}
For completeness, direct binary multiplication gives
$\operatorname{tr}E=0$, $\operatorname{tr}(E\circ E)=2\sum_i|q_i|^2$,
and $\operatorname{tr}(E\circ(E\circ E))=6\mathfrak t(q)$.
For the last equality, the three cyclic off-diagonal coordinates of
$E\circ E$ are $\overline{q_2q_3},\overline{q_3q_1},
\overline{q_1q_2}$; their trace pairings with the corresponding
$q_i$ each give $2\mathfrak t(q)$.
The Jordan trace formula
$\det X=\operatorname{tr}(X\circ(X\circ X))/3
-\operatorname{tr}X\operatorname{tr}(X\circ X)/2
+(\operatorname{tr}X)^3/6$
therefore gives the first identity after expanding $X=\lambda I+E$
and substituting the proved norm equality. For the second,
each off-diagonal octonion
occurs in both transposed matrix entries, giving exactly
\(2\sum_i\operatorname{Re}(q_i\overline{p_i})\).
Real associativity and cyclicity of octonion real parts give
\(\mathfrak t(q_3,q_1,q_2)=\mathfrak t(q_1,q_2,q_3)\):
\[
 \operatorname{Re}((q_3q_1)q_2)
 =\operatorname{Re}(q_3(q_1q_2))
 =\operatorname{Re}((q_1q_2)q_3).
\]
Consequently \(\mathfrak t_N(rX)=\mathfrak t_N(X)\), and the
whole determinant in \eqref{n24:eq:niemeier-albert-cubic} is
\(r\)-invariant with \(\lambda\) fixed.  Equivalently, for the real
permutation matrix \(P e_i=e_{i+1}\), indices modulo three,
\[
 P\bigl(\lambda I+E(q)\bigr)P^T
             =\lambda I+E(q_3,q_1,q_2).
\]
This is also a Jordan automorphism: multiplication by a real
permutation matrix merely reindexes the two-factor matrix products,
so it preserves \((AB+BA)/2\).  It therefore carries the positive
cone and its determinant boundary to themselves.  The cyclic bridge
is simultaneously linear, isometric on \(\mathbb O^3\), integral
on the stated lattice image, and compatible with the actual Albert
cubic and positive cone.

\paragraph{Dependence and the four-dimensional lattice summand.}
The map uses precisely the four specified origins \(f_i\), the
generator \(289\), the even/odd triplet order, the displayed order
of \(b\), the ordered \(D\)-coordinates, and the chosen orthonormal
octonion basis.  Changing one of these changes the displayed map;
the construction makes no claim of independence from these data.
Its restriction to the actual \(D\)-summand is
\begin{equation}
 d\longmapsto\frac1{\sqrt3}(z(d),z(d),z(d)),\qquad
 z(d)=d_1u_4+d_2u_5+d_3u_6+d_4u_7.
 \label{n24:eq:d-diagonal-embedding}
\end{equation}
This is an isometry on \(D\), and on \(D^\vee\) in its real span.
In particular its 24 shortest nonzero dual vectors become 24
specified points in this four-dimensional diagonal subspace.
The separate Cartan-weight map \(d\mapsto\sum_kd_k\varepsilon_k\)
is related to it by the exact isometry
\(\sum_kd_k\varepsilon_k\mapsto(z(d),z(d),z(d))/\sqrt3\).

\subsection{The Leech neighbour in the same three octonion coordinates}
\label{n24:sec:leech-coordinates}

The preceding map also transports the entire neighbour construction,
not only the real span of $N$. Put $N_{\mathbb O}=\mathcal L(N)$ and
retain the actual $\rho,\alpha,K,\Lambda$ of
Theorem~\ref{esnl:thm:leech-neighbor}. In the already fixed basis let
\[
 U=\sum_{a=0}^7u_a,\qquad
 V=\frac{u_0+u_1+u_2+u_3}{\sqrt2}
       +\sqrt3u_4+\frac{2u_5+u_6}{\sqrt3}.
\]
Substitution in the original isometry gives the explicit vectors
\begin{equation}
 \begin{aligned}
 \varrho=\mathcal L\rho&=(V+2U,V,V-2U),\\
 \delta=\mathcal L\alpha
 &=\left(\frac{(2+\sqrt2)u_0-2u_1}{3},
          \frac{(\sqrt2-1)u_0+u_1}{3},
          \frac{(\sqrt2-1)u_0+u_1}{3}\right).
 \end{aligned}
 \label{n24:eq:weyl-in-octonions}
\end{equation}
Indeed each even triplet of $\rho_A$ has mean $1/2$, each odd
triplet has mean $-1/2$, and both residual triplets are $(2,0,-2)$.
The four mean-difference coordinates are $\sqrt6/2$, while the
four $D$ coordinates are $(3,2,1,0)$. For $\alpha$, the two residual
triplets in the first fibre are $(2,-1,-1)/3$ and $(-2,1,1)/3$;
its only nonzero mean-difference coordinate is $2/\sqrt6$.
These facts prove every coefficient in
\eqref{n24:eq:weyl-in-octonions}. In particular
$|\varrho|^2=84$, $|\delta|^2=2$, and
$\langle\varrho,\delta\rangle=1$ in the product octonion metric.

Define
\begin{equation}
 \begin{aligned}
 K_{\mathbb O}&=\{q\in N_{\mathbb O}:
                          \langle\varrho,q\rangle\equiv0\pmod6\},\\
 \Lambda_{\mathrm{Leech},\mathbb O}
       &=K_{\mathbb O}+\mathbb Z(\varrho/6-\delta).
 \end{aligned}
 \label{n24:eq:leech-in-octonions}
\end{equation}
The isometry and its inverse prove, in both directions,
$\mathcal L(K)=K_{\mathbb O}$ and
$\mathcal L(\Lambda)=\Lambda_{\mathrm{Leech},\mathbb O}$.
Thus \eqref{n24:eq:leech-in-octonions} is a full coordinate
description of the Leech lattice inside the very same three sectors.
The glue test above, followed by its displayed height congruence,
is an exact membership test for $K_{\mathbb O}$; adjoining the six
cosets $k(\varrho/6-\delta)$, $0\le k<6$, gives all of the Leech
lattice. The common-sublattice indices, discriminant form, and
order-twelve quotient are transported without alteration.

The cyclic action also acts on the marked family of neighbours.
Retain $r=\rho(361)$ and its corresponding map
$c(q_1,q_2,q_3)=(q_3,q_1,q_2)$.
For any transported Weyl vector $\rho'=T\rho$, $T\in O(N)$,
the neighbour is defined using $T\alpha$, so that
$\langle\rho',T\alpha\rangle=1$. Then
\[
 c\,\mathcal L(K_{\rho'})=\mathcal L(K_{r\rho'}),\qquad
 c\,\mathcal L(\Lambda_{\rho'})=\mathcal L(\Lambda_{r\rho'}).
\]
This follows directly from $\mathcal Lr=c\mathcal L$ and
(NL18). The inverse maps use $r^{-1}$ and $c^{-1}$.
It specifies exactly which marked neighbour is the target; it does
not silently assume that the chosen $\rho$ is fixed by $r$.
Finally the common cubic evaluates on every one of these lattice
points as
\[
 \det Q_{\mathbb O}(\lambda;\mathcal LX)
    =\lambda^3-\lambda(X,X)
                   +2\operatorname{Re}(((\mathcal LX)_1
                                         (\mathcal LX)_2)(\mathcal LX)_3).
\]
This identity holds for all $X\in N_{\mathbb R}$, hence also
on the explicitly embedded $N$, $K$ and $\Lambda$.

\subsection{Exactly which root isometries preserve this glue}
\label{n24:sec:glue-symmetry}

Retain the coordinate convention \((P_px)_i=x_{p(i)}\).
An arbitrary root-lattice isometry has the original form
\[
 T(X_1,X_2,X_3,X_4;d)
   =\bigl((\varepsilon_iQ_{\sigma_i}X_{p(i)})_{i=1}^4;\ T_Dd\bigr),
\]
where \(\varepsilon_i\in\{\pm1\}\), \(\sigma_i\in S_6\),
\(p\in S_4\), and \(T_D\in O(D)\).  If \(B=d_D(T_D)\),
the necessary and sufficient condition for \(T(N)=N\) is
\begin{equation}
 B\tau(u)=\tau(P_pu)\ (u\in E),\qquad
 \bigl(\varepsilon_i\ell_{p(i)}\bigr)_{i=1}^4
                  =\bigl(\ell_iA\bigr)_{i=1}^4
 \quad\hbox{for some }A\in GL_2(\mathbb F_3),
 \label{n24:eq:all-glue-preservers}
\end{equation}
where
\(\ell_1=(1,1),\ell_2=(1,2),\ell_3=(1,0),\ell_4=(0,1)\).
The \(A\), and hence \(B=B_p\), are unique.  Indeed the
two-primary part of the glue is the graph of \(\tau\), and signs
are trivial modulo two.  Its three-primary part is the image of the
injective map \(j(a,b)=(a+b,a+2b,a,b)\), so preserving it is
equivalent to \(M_\varepsilon P_pj=jA\).  These are exactly
the two conditions in \eqref{n24:eq:all-glue-preservers}; the
primary decomposition proves their sufficiency as well as necessity.
Since the four \(\ell_i\) represent the four distinct projective
row lines over \(\mathbb F_3\), each \(A\) gives exactly one
signed permutation by this equation.  This recovers the precise
48-element glue quotient, not merely its action on three labels.

The canonical full-automorphism theorem proves, with these same
coordinates and section, that
\[
 O(N)=\bigl(S_6^4\times W(D_4)\bigr)\rtimes L(G),
 \qquad G\cong GL_2(\mathbb F_3).
\]
Conjugation by the actual map \(\mathcal L\) now gives the entire
orthogonal stabilizer of the embedded lattice, with explicit matrices
obtained from \eqref{n24:eq:full-isometry} and its inverse:
\begin{equation}
 O(\Lambda_{\mathbb O})
       =\mathcal L O(N)\mathcal L^{-1},\qquad
 \Gamma_G(g)=\mathcal LL_g\mathcal L^{-1},\qquad
 \mathcal LL_g=\Gamma_G(g)\mathcal L.
 \label{n24:eq:full-transported-action}
\end{equation}
Both inclusions in the first equality follow by conjugating an
isometry preserving the specified lattice.  This is a faithful
action on the whole 24-dimensional space. Each \(\Gamma_G(g)\)
commutes with \(c\),
because the canonical equation
\(L_g\rho(289^2)=\rho(289^2)L_g\) conjugates to that equation.

For completeness, the previously displayed \(D_4\)-triality
section and geometric frame permutations have a precise comparison
even though that particular proposed linear conjugacy fails.  On
\(V_N\), each of \(\widehat R,\widehat H\) has trace
\(5-5+2=2\): the transposed \(A_5\) pair contributes zero,
the other two components contribute \(5,-5\), and the
\(D_4\) reflection contributes 2.  Their three-cycle has trace
\(5+1=6\).  A geometric Albert-frame transposition acts on
\(\mathbb O^3\) by swapping two sectors with conjugation and by
conjugating the fixed sector; its trace is
\(1-7=-6\).  Its three-cycle has trace zero.  Therefore an
invertible map cannot conjugate these two particular \(S_3\)
representations, since trace is invariant under conjugation.
The complete real character decompositions are, respectively,
\[
 V_N|_{\langle\widehat R,\widehat H\rangle}
       \cong7\mathbf1\oplus5\operatorname{sgn}\oplus6\operatorname{Std},
 \qquad
 \mathbb O^3|_{S_3^{\rm frame}}
       \cong\mathbf1\oplus7\operatorname{sgn}\oplus8\operatorname{Std}.
\]
These multiplicities follow by taking inner products with the
characters \((1,1,1),(1,-1,1),(2,0,-1)\), on classes of sizes
\(1,3,2\).  The exact whole-space relation for the original section
is \eqref{n24:eq:full-transported-action}; the exact
determinant-compatible cyclic relation is
\eqref{n24:eq:cyclic-intertwining}.  In the latter relation the
residue element has trivial discriminant action while its octonion
action cycles the sectors.  Those are two explicitly calculated
actions of the same operator on two explicitly different constructions;
no identification between their quotient homomorphisms is assumed.

Finally, the connected group \(\operatorname{Spin}(8)\) acts on
the ambient triality module, and its subgroup preserving
\(\Lambda_{\mathbb O}\) is exactly its intersection with
\eqref{n24:eq:full-transported-action}.  It is finite: the columns
of a lattice isometry in a fixed lattice basis have prescribed
bounded lengths, so each column has finitely many possible lattice
values.  A connected group preserving a discrete full lattice
pointwise as a set acts trivially, since each orbit of a lattice
vector is connected and discrete and the lattice spans.  The
nontrivial connected triality action therefore does not preserve
the whole lattice.  The proved cyclic subgroup \(\langle c\rangle\)
does preserve both the lattice and its Albert cubic.  Equations
\eqref{n24:eq:all-glue-preservers} and
\eqref{n24:eq:full-transported-action} specify the full glue
stabilizer; they do not assert that every element of that finite
orthogonal group also preserves the cubic.

% Proof-bearing staged module. Namespace nca:. No canonical edits.
\section{The arithmetic of the retained Albert cubic on both lattices}
\label{nca:sec:cubic-arithmetic}

\subsection{Original data and the exact polynomial}

Retain the ordered coordinates \(X=(x_{i,j};d_k)\),
\(1\leq i\leq4\), \(0\leq j\leq5\), \(1\leq k\leq4\), with
\(\sum_{j=0}^5x_{i,j}=0\), and their original Euclidean metric.
Let
\[
 A_5=\{x\in\mathbb Z^6:\sum_jx_j=0\},\qquad
 D_4=\{d\in\mathbb Z^4:\sum_kd_k\equiv0\pmod2\},\qquad
 R=A_5^{\perp4}\perp D_4.
\]
Write \(\lambda=(5,-1,-1,-1,-1,-1)/6\),
\(v_D=(1,0,0,0)\), and \(s_D=(1,1,1,1)/2\).
The original discriminant glue is generated by
\[
 \begin{split}
 &(3,3,3,3;00),\quad(3,3,0,0;10),\quad(3,0,3,0;01),\\
 &(2,2,2,0;00),\quad(2,4,0,2;00)
 \end{split}
\]
in \((\mathbb Z/6)^4\oplus\mathbb F_2^2\), with the same marking
\(10=v_D+D_4\), \(01=s_D+D_4\). Its order is 72: the first
three generators have independent two-primary reductions, and the
last two independent three-primary reductions.  The lattice \(N\)
is its inverse image in \(R^\vee\).

Retain
\[
 \rho_A=(5,3,1,-1,-3,-5)/2,\qquad
 \rho_D=(3,2,1,0),\qquad \rho=\rho_A^{\oplus4}\perp\rho_D,
 \qquad \alpha=(e_0-e_1,0,0,0;0),
\]
and the exact neighbour
\[
 K=\{X\in N:(\rho,X)\equiv0\pmod6\},\qquad
 v=\rho/6-\alpha,\qquad \Lambda=K+\mathbb Zv.
\]
Here \(\rho\in N\), \(\rho^2=84\), and \((\rho,\alpha)=1\),
as follows immediately from these coordinates and the first glue
generator.  The name \(\Lambda\) always denotes this neighbour.

Fix an ordered orthonormal octonion basis \(u_0=1,u_1,\ldots,u_7\).
For the first part of the calculation no multiplication table is
imposed beyond the actual multiplication of these basis elements.
Define, in the original coordinates,
\begin{align}
 s_i&=\sum_{j=0}^5(-1)^jx_{i,j},\label{nca:eq:original-s}\\
 A_{2(i-1)+p}
  &=x_{i,p}-\frac{x_{i,p}+x_{i,p+2}+x_{i,p+4}}3,
       &&p=0,1,\label{nca:eq:original-A}\\
 B_{2(i-1)+p}
  &=x_{i,p+2}-\frac{x_{i,p}+x_{i,p+2}+x_{i,p+4}}3,
       &&p=0,1.\label{nca:eq:original-B}
\end{align}
Write \(A=A_0+\mathbf a\), \(B=B_0+\mathbf b\), where
\(\mathbf a=\sum_{h=1}^7A_hu_h\) and
\(\mathbf b=\sum_{h=1}^7B_hu_h\).  The common mean octonion is
\begin{equation}
 M=\frac{\sqrt2}{6}\sum_{a=0}^3s_{a+1}u_a
       +\frac{\sqrt3}{3}\sum_{k=1}^4d_ku_{k+3}.
 \label{nca:eq:mean}
\end{equation}
The retained isometry therefore has precisely the ordered form
\[
 \mathcal LX=(A+M,B+M,-A-B+M).
\]
This is an identity with its original three triplets: the residuals
sum to zero, the first four common mean coefficients are
\(s_i/\sqrt{18}=\sqrt2s_i/6\), and the last four are
\(d_k/\sqrt3=\sqrt3d_k/3\).
For completeness the original squared norm and the output norm are both
\[
 |A|^2+|B|^2+|A+B|^2+\frac16\sum_i s_i^2+\sum_kd_k^2.
\]
In the original fibres this follows by separating the means
\(s_i/6,-s_i/6\) from their zero-sum triplets; in the output it
follows from \(A+B+(-A-B)=0\) and
\(3|M|^2=\sum_i s_i^2/6+\sum_kd_k^2\). The inverse recovers
\(M=(q_1+q_2+q_3)/3\), then \(s_i=3\sqrt2M_{i-1}\),
\(d_k=\sqrt3M_{k+3}\), and \(A=q_1-M,B=q_2-M\).
Adding \((-1)^ps_i/6\) to the \(p\)-th residual triplet in the
\(i\)-th fibre recovers each original \(x_{i,p+2j}\). Thus the
map is onto and preserves the original pairing by polarization.

For imaginary octonions put
\[
 \mathbf a\times\mathbf b=\frac{\mathbf a\mathbf b-
                                    \mathbf b\mathbf a}{2},\qquad
 \phi(\mathbf a,\mathbf b,\mathbf c)
                     =\langle\mathbf a\times\mathbf b,\mathbf c\rangle.
\]
Let \(\phi_{ijk}=\phi(u_i,u_j,u_k)\).  Thus \(\phi\) is the
actual alternating octonion three-form in the chosen frame, including
all signs. Define
\begin{align}
 D&=A_0^2+A_0B_0+B_0^2-\langle\mathbf a,\mathbf a\rangle
             -\langle\mathbf a,\mathbf b\rangle
             -\langle\mathbf b,\mathbf b\rangle,\label{nca:eq:D}\\
 Z&=(2A_0+B_0)\mathbf a+(A_0+2B_0)\mathbf b
                               -3\mathbf a\times\mathbf b
       =\sum_{h=1}^7Z_hu_h.\label{nca:eq:Z}
\end{align}
Equivalently, the component in the final formula is the exact
quadratic polynomial
\[
 Z_h=(2A_0+B_0)A_h+(A_0+2B_0)B_h
                  -3\sum_{i,j=1}^7\phi_{ijh}A_iB_j.
\]

\begin{theorem}[Exact pullback with every radical and coefficient]
\label{nca:thm:polynomial}
The original cubic
\(F(X)=\operatorname{Re}(((\mathcal LX)_1(\mathcal LX)_2)
                                           (\mathcal LX)_3)\)
has the complete polynomial expression
\begin{equation}
 F(X)=P_0(X)+\sqrt2P_2(X)+\sqrt3P_3(X),
 \label{nca:eq:polynomial}
\end{equation}
where every variable is the displayed linear function of the original
\(x_{i,j},d_k\), and
\begin{align}
 P_0={}&-A_0B_0(A_0+B_0)+B_0\langle\mathbf a,\mathbf a\rangle
              +A_0\langle\mathbf b,\mathbf b\rangle
              +2(A_0+B_0)\langle\mathbf a,\mathbf b\rangle,
                \label{nca:eq:P0}\\
 P_2={}&\frac{-s_1D+s_2Z_1+s_3Z_2+s_4Z_3}{6}
       +\frac{s_1(4s_1^2-3\sum_{i=1}^4s_i^2)}{108}
       -\frac{s_1\sum_{k=1}^4d_k^2}{6},\label{nca:eq:P2}\\
 P_3={}&\frac{d_1Z_4+d_2Z_5+d_3Z_6+d_4Z_7}{3}.
                  \label{nca:eq:P3}
\end{align}
These identities hold throughout the original real vector space, in
particular on \(N\), \(K\), and \(\Lambda\). They retain the ordered
Albert determinant as
\[
 \det Q_{\mathbb O}(\lambda;\mathcal LX)
                =\lambda^3-\lambda(X,X)+2F(X).
\]
\end{theorem}

\begin{proof}
For \(a=a_0+\mathbf a\), \(b=b_0+\mathbf b\),
\(c=c_0+\mathbf c\), binary octonion multiplication gives
\[
 \operatorname{Re}((ab)c)
  =a_0b_0c_0-a_0\langle\mathbf b,\mathbf c\rangle
   -b_0\langle\mathbf a,\mathbf c\rangle
   -c_0\langle\mathbf a,\mathbf b\rangle
   -\phi(\mathbf a,\mathbf b,\mathbf c).
\]
Call the first four, symmetric, terms \(S(a,b,c)\).  Substitute
\((a,b,c)=(A+M,B+M,-A-B+M)\).  Terms with exactly two copies
of \(M\) in \(S\) sum to \(S(M,M,A+B-A-B)=0\). Terms
with exactly one copy sum to
\(-S(M,A,A)-S(M,A,B)-S(M,B,B)\).  The term with no copy is
\(S(A,B,-A-B)=P_0\), by multiplying its four displayed terms.
The alternating term with no copy vanishes; terms with two or
three copies vanish. Its three terms with one copy all equal
\(\phi(\mathbf a,\mathbf b,\operatorname{Im}M)\), since
\[
 \phi(\mathbf a,\operatorname{Im}M,-\mathbf a-\mathbf b)
 =\phi(\mathbf a,\mathbf b,\operatorname{Im}M),\quad
 \phi(\operatorname{Im}M,\mathbf b,-\mathbf a-\mathbf b)
 =\phi(\mathbf a,\mathbf b,\operatorname{Im}M).
\]
Thus all one-copy terms together are
\(-M_0D+\langle\operatorname{Im}M,Z\rangle\).
Finally the term with three copies is
\[
 \operatorname{Re}(M^3)=4M_0^3-3M_0|M|^2,
 \qquad |M|^2=\frac{\sum_i s_i^2}{18}
                              +\frac{\sum_kd_k^2}{3},
 \qquad M_0=\frac{\sqrt2s_1}{6}.
\]
Substitution proves \eqref{nca:eq:P0}--\eqref{nca:eq:P3} with
their displayed factors. In particular it proves the absence of a
\(\sqrt6\) term without discarding any mixed term. The determinant
formula uses the original isometry and its unchanged squared norm.
\end{proof}

\subsection{Finite generators of both original lattices}

Order the 24 original simple roots as
\[
 g_{5(i-1)+j+1}=a_{i,j}=e_{i,j}-e_{i,j+1}
 \quad(1\leq i\leq4,\ 0\leq j\leq4),
\]
followed by
\[
 g_{21}=e_1-e_2,\quad g_{22}=e_2-e_3,\quad
 g_{23}=e_3-e_4,\quad g_{24}=e_3+e_4
\]
in the original \(D_4\) coordinates.  All unspecified components
are zero.  Adjoin the five actual glue lifts
\begin{align*}
 g_{25}&=(3\lambda,3\lambda,3\lambda,3\lambda;0),&
 g_{26}&=(3\lambda,3\lambda,0,0;v_D),\\
 g_{27}&=(3\lambda,0,3\lambda,0;s_D),&
 g_{28}&=(2\lambda,2\lambda,2\lambda,0;0),\\
 g_{29}&=(2\lambda,4\lambda,0,2\lambda;0).
\end{align*}
Then
\begin{equation}
 N=\sum_{i=1}^{29}\mathbb Zg_i,\qquad
 H_i=(\rho,g_i)=
 \begin{cases}
 1,&1\leq i\leq24,\\
 30,18,18,15,20,&i=25,26,27,28,29\text{ respectively}.
 \end{cases}
 \label{nca:eq:N-generators}
\end{equation}
Indeed the first 24 vectors span \(R\), the last five generate
the original glue, and subtraction of the corresponding glue lift
leaves every element of \(N\) in \(R\).  The height calculation
uses \((\rho_A,\lambda)=5/2\),
\((\rho_D,v_D)=(\rho_D,s_D)=3\), and the simple-root heights one.

Define the following 30, fully specified neighbour generators:
\begin{equation}
 h_1=6\alpha,\qquad h_i=g_i-H_i\alpha\quad(2\leq i\leq29),
 \qquad h_{30}=v=\rho/6-\alpha.
 \label{nca:eq:Lambda-generators}
\end{equation}
Then \(K=\sum_{i=1}^{29}\mathbb Zh_i\) and
\(\Lambda=\sum_{i=1}^{30}\mathbb Zh_i\).
For if \(X=\sum_i n_ig_i\in K\), then
\[
 X=\sum_{i=2}^{29}n_i(g_i-H_i\alpha)
                         +\left(\sum_{i=1}^{29}n_iH_i\right)\alpha,
\]
and the last coefficient is divisible by six. Conversely every
displayed generator before \(h_{30}\) has height divisible by six.
Adjoining \(h_{30}\) is exactly the original neighbour definition.

\subsection{A complete finite calculation of the cubic value modules}

For a lattice \(L\) and a homogeneous real cubic \(F\), define
its additive value module by
\[
 \mathcal V_F(L)=\sum_{X\in L}\mathbb ZF(X)\subset\mathbb R.
\]
This definition retains every value and takes its generated additive
group. It does not claim that every element of the group is one
individual value of \(F\).

\begin{theorem}[Finite cubic-value generators, including negative integers]
\label{nca:thm:finite-values}
If \(L=\sum_{i=1}^m\mathbb Zb_i\), allowing dependent generators,
then \(\mathcal V_F(L)\) is generated by the following finite list:
\begin{align}
 A_i={}&F(b_i),&&1\leq i\leq m,\label{nca:eq:single}\\
 D^+_{ij}={}&F(b_i+b_j)-F(b_i)-F(b_j),&&i<j,\label{nca:eq:Dplus}\\
 D^-_{ij}={}&F(b_i-b_j)-F(b_i)+F(b_j),&&i<j,\label{nca:eq:Dminus}\\
 D_{ijk}={}&F(b_i+b_j+b_k)-F(b_i+b_j)-F(b_i+b_k)\notag\\
          &{}-F(b_j+b_k)+F(b_i)+F(b_j)+F(b_k),&&i<j<k.
 \label{nca:eq:Dthree}
\end{align}
Consequently \eqref{nca:eq:N-generators} gives exactly 4495
specified generators of \(\mathcal V_F(N)\), and
\eqref{nca:eq:Lambda-generators} gives exactly 4960 specified
generators of \(\mathcal V_F(\Lambda)\), for every fixed
orthonormal octonion frame and its actual \(\phi_{ijk}\).
\end{theorem}

\begin{proof}
Let \(B_F\) be the symmetric trilinear polarization, with
\(B_F(X,X,X)=F(X)\); explicitly its value is
\(1/6\) times the seven-term expression in
\eqref{nca:eq:Dthree}, with its three inputs replacing the three
generators. For every integer tuple \(n\), expansion gives
\[
 F\left(\sum_i n_ib_i\right)
 =\sum_i n_i^3A_i
 +\sum_{i<j}\bigl(n_i^2n_j C_{ij}+n_in_j^2 E_{ij}\bigr)
 +\sum_{i<j<k}n_in_jn_kD_{ijk},
\]
where \(C_{ij}=3B_F(b_i,b_i,b_j)\) and
\(E_{ij}=3B_F(b_i,b_j,b_j)\).  The two pair generators are
\(D^+_{ij}=C_{ij}+E_{ij}\) and
\(D^-_{ij}=-C_{ij}+E_{ij}\). Therefore their contribution is
\[
 \frac{n_in_j(n_i+n_j)}2D^+_{ij}
       +\frac{n_in_j(n_j-n_i)}2D^-_{ij}.
\]
Both coefficients are integers, for all positive, zero, and negative
integers \(n_i,n_j\): if either integer is even the claim is immediate,
and if both are odd their sum and difference are even. This proves
that the finite list generates all values. Conversely each entry
of the list is an integer linear combination of actual values on
\(L\). The proof does not use linear independence of the generators.
Finally \(29+2\binom{29}{2}+\binom{29}{3}=4495\) and
\(30+2\binom{30}{2}+\binom{30}{3}=4960\).
\end{proof}

For numerical arithmetic now specify the additional octonion frame
choice completely.  Use \(\mathbb O=\mathbb H\oplus\mathbb H\ell\)
in ordered coordinates \((a,b)\), with
\begin{equation}
 \begin{gathered}
 (a,b)(c,d)=(ac-\overline d\,b,\ da+b\overline c),\\
 u_0=(1,0),\quad u_1=(i,0),\quad u_2=(j,0),\quad u_3=(k,0),\\
 u_4=(0,1),\quad u_5=(0,i),\quad u_6=(0,j),\quad u_7=(0,k).
 \end{gathered}
 \label{nca:eq:Cayley-frame}
\end{equation}
Here \(ij=k\), \(jk=i\), \(ki=j\), and the reversed products
are negative.  The positive oriented three-form triples are exactly
\begin{equation}
 (123),\ (145),\ (176),\ (246),\ (257),\ (347),\ (365).
 \label{nca:eq:seven-triples}
\end{equation}
Every \(\phi_{ijk}\) is obtained from these by alternation, or is
zero.  This is a specified frame for the retained, previously
frame-dependent isometry, without a change to any lattice coordinates
or to the Albert matrix convention.

\begin{theorem}[Exact value modules and their quotient of order twelve]
\label{nca:thm:value-modules}
For the frame \eqref{nca:eq:Cayley-frame}, the exact additive modules are
\begin{align}
 \mathcal V_F(N)
    &=\frac1{27}\mathbb Z\ \oplus\
       \frac{\sqrt2}{54}\mathbb Z\ \oplus\
       \frac{\sqrt3}{18}\mathbb Z,\label{nca:eq:N-values}\\
 \mathcal V_F(\Lambda)
    &=\frac1{27}\mathbb Z\ \oplus\
       \frac{\sqrt2}{216}\mathbb Z\ \oplus\
       \frac{\sqrt3}{54}\mathbb Z.\label{nca:eq:Lambda-values}
\end{align}
In particular the inclusion of these two additive subgroups of
\(\mathbb R\) has the exact quotient
\begin{equation}
 \mathcal V_F(\Lambda)/\mathcal V_F(N)
                  \cong\mathbb Z/4\oplus\mathbb Z/3
                  \cong\mathbb Z/12.
 \label{nca:eq:value-quotient}
\end{equation}
This is an inclusion of cubic value modules; it does not assert the
false lattice inclusion \(N\subset\Lambda\).
\end{theorem}

\begin{proof}
The polynomial coefficients in
\eqref{nca:eq:P0}--\eqref{nca:eq:P3} are rational on both lattices
for this exact multiplication table.  The numbers
\(1,\sqrt2,\sqrt3\) are linearly independent over \(\mathbb Q\).
For example, isolating one radical in a rational relation and
squaring either gives a rational multiple of another irrational
radical, or a rational square equal to two, three, or \(3/2\),
which is impossible by parity of prime exponents. Thus a value has
a unique rational coefficient triple \((P_0,P_2,P_3)\).

Apply the complete finite list in
Theorem~\ref{nca:thm:finite-values} to the displayed original
generators. Exact substitution gives the following least common
multiples of coefficient denominators. The denominators in each
entry are ordered as \((P_0,P_2,P_3)\), and a zero coefficient has
denominator one:
\begin{equation}
\begin{array}{c|c|c|c|c}
 &A_i&D^+_{ij}&D^-_{ij}&D_{ijk}\\\hline
 N&(27,27,1)&(9,18,18)&(9,18,18)&(9,18,18)\\
 \Lambda&(27,216,27)&(27,108,54)&(27,108,54)&(9,36,54)
\end{array}.
\label{nca:eq:denominator-certificate}
\end{equation}
This table is a finite rational calculation, with no bound on
lattice points or heuristic sampling.  The complete arithmetic
certificate is the accompanying \texttt{verify\_cubic.py}: its
\texttt{residual\_data} and \texttt{cubic} functions are exactly
\eqref{nca:eq:original-s}--\eqref{nca:eq:P3}; it forms every
\(A_i,D^+_{ij},D^-_{ij},D_{ijk}\) for the 29 or 30 explicitly
defined vectors, using rational arithmetic.  The \texttt{cross}
function executes the three cyclic positive and three negative
entries for each of the seven triples
\eqref{nca:eq:seven-triples}.  Hence the full finite calculation
is specified by these equations, the finite index bounds in
Theorem~\ref{nca:thm:finite-values}, and integer arithmetic.
Its 4495 and 4960 entries verify the upper inclusions in
\eqref{nca:eq:N-values}--\eqref{nca:eq:Lambda-values}.

Here are short certificates for the reverse inclusions. Coordinates
in the next two tables are relative to the respective three
proposed generators on the right sides of
\eqref{nca:eq:N-values} and \eqref{nca:eq:Lambda-values}.
In the first table all pair symbols use the \(g_i\):
\begin{equation}
\begin{array}{c|r r r}
 F(g_1)&-4&4&0\\
 F(g_2)&-5&-4&0\\
 F(g_7)&0&18&0\\
 F(g_{28})&-80&52&0\\
 D^+_{21,28}&0&-36&16\\
 D^+_{2,27}&-27&27&3\\
 D^+_{2,25}&-81&45&0
\end{array}.
\label{nca:eq:N-lower-certificate}
\end{equation}
The determinants of the three columns obtained by transposing
rows \((1,4,5)\), \((2,3,6)\), and \((2,7,6)\) are,
respectively, \(1792,-270,-1647\). Their greatest common divisor
is one. In the second table all pair symbols use the \(h_i\):
\begin{equation}
\begin{array}{c|r r r}
 F(h_2)&14&-272&0\\
 F(h_4)&5&-128&0\\
 F(h_{28})&13780&-74672&0\\
 F(h_{30})&14&-63&-2\\
 D^+_{2,27}&7425&-60624&-9
\end{array}.
\label{nca:eq:Lambda-lower-certificate}
\end{equation}
Rows \((1,3,4)\), \((1,4,5)\), and \((2,4,5)\) give
determinants \(-5405504,2315394,1281267\), whose greatest
common divisor is also one.

For clarity, the elementary integer argument using these
determinants does not require a normal-form theorem.  Let
\(W\subset\mathbb Z^3\) be the group generated by one displayed
table.  For any three of its columns forming a matrix \(C\),
the adjugate identity \(C\operatorname{adj}(C)=\det(C)I\)
implies \(\det(C)\mathbb Z^3\subset W\).  Bezout's identity
for the three displayed determinants, whose gcd is one, therefore
gives \(\mathbb Z^3\subset W\). The opposite inclusion is
immediate from the integer entries. Every table entry is an
integer combination of actual values by its definition, proving
the two reverse inclusions.  Finally the coordinates of the
inclusion of the first module in the second are
\((a,b,c)\mapsto(a,4b,3c)\). Its cokernel is the displayed
product of cyclic groups, with order twelve. This also proves
the asserted exact quotient map.
\end{proof}

\subsection{Frame dependence and exact symmetry transport}

The finite-list theorem is valid for the actual coefficients of
every orthonormal frame. The numerical modules
\eqref{nca:eq:N-values}--\eqref{nca:eq:Lambda-values} refer to
the frame \eqref{nca:eq:Cayley-frame}.  This dependence has an
exact comparison map. If \(u_a'=Tu_a\) for a real isometry
\(T\) fixing \(1\), then
\[
 \mathcal L_{u'}=(T,T,T)\mathcal L_u.
\]
The symmetric part \(S\) of the trilinear tensor is unchanged,
because \(T\) preserves the scalar coordinate and inner product.
Writing \(\mathbf m=\operatorname{Im}M\), its exact change is
\begin{equation}
 F_{u'}(X)-F_u(X)
  =-3\bigl(\phi(T\mathbf a,T\mathbf b,T\mathbf m)
                  -\phi(\mathbf a,\mathbf b,\mathbf m)\bigr).
 \label{nca:eq:frame-transport}
\end{equation}
This follows from the complete expansion in the polynomial proof.
If \(T\) is an octonion automorphism it preserves \(\phi\), so
the cubic and its value modules are unchanged.

The dependence on a general orthonormal frame is substantive. Take
the actual lattice vector with the only nonzero fibre coordinates
\[
 x_{1,1}=1,\ x_{1,5}=-1,\quad
 x_{2,2}=1,\ x_{2,4}=-1,\qquad d=(1,1,0,0).
\]
It lies in \(R\subset N\), has \(A=u_1\), \(B=u_2\), and
\(M=(u_4+u_5)/\sqrt3\). In the Cayley frame its value is zero.
Replace only
\[
 u'_3=(2u_3-\sqrt5u_4)/3,\qquad
 u'_4=(\sqrt5u_3+2u_4)/3.
\]
This is an exact orthonormal frame change; the same original
lattice vector now has value \(-\sqrt{15}/3\), because
\(\phi(u_1,u_2,u'_4)=\sqrt5/3\) and
\(\phi(u_1,u_2,u_5)=0\). This value is outside
\(\mathbb Q+\sqrt2\mathbb Q+\sqrt3\mathbb Q\). Indeed the
independent sign changes of \(\sqrt2,\sqrt3\) are automorphisms
of \(\mathbb Q(\sqrt2,\sqrt3)\). If an element \(z\) in that
field had \(z^2=15\), each sign change would send it to
\(z\) or \(-z\). In the basis
\(1,\sqrt2,\sqrt3,\sqrt6\) this makes \(z\) a rational
multiple of a single basis vector. Its square would force one
of \(15,15/2,5,5/2\) to be a rational square, impossible by
parity of its prime exponents.
Thus the numerical specialization is not silently extended to
an unspecified orthonormal frame.

Return to the explicit Cayley frame for the following numerical
witnesses. Let \(\sigma=\rho(289)\) act simultaneously by
\(e_{i,j}\mapsto e_{i,j+1\bmod6}\), and trivially on \(D_4\).
The exact isometry gives
\(\mathcal L\sigma^2=c\mathcal L\) for
\(c(q_1,q_2,q_3)=(q_3,q_1,q_2)\). Cyclic real associativity
therefore proves \(F(\sigma^{2k}X)=F(X)\) for every real \(X\).
These maps preserve \(N\) by their original discriminant action.
The stabilizer of this cubic within the six-element residue group
is exactly its three-element subgroup: direct substitution at
\(\alpha=g_1\) gives
\begin{equation}
 F(\alpha)=\frac{-4+2\sqrt2}{27},\qquad
 F(\sigma\alpha)=\frac{-5-2\sqrt2}{27},
 \label{nca:eq:residue-nonpreserver}
\end{equation}
which differ. All odd powers have the same effect on the cubic
as \(\sigma\), by its proved invariance under even powers.

The original reflection exchanging \(e_{1,0},e_{1,1}\) sends
\(\alpha\) to \(-\alpha\), hence changes its nonzero cubic
value to its negative. The original full glue-preserving lifts
\[
 \begin{aligned}
 \widehat R(X_1,X_2,X_3,X_4;d)
       &=(X_2,X_1,X_3,-X_4;R_Dd),\\
 \widehat H(X_1,X_2,X_3,X_4;d)
       &=(X_1,X_3,X_2,-X_4;H_Dd).
 \end{aligned}
\]
also have explicit witnesses:
\[
 F(\widehat R\alpha)=0,\qquad
 F(g_7)=\frac{\sqrt2}{3},\qquad F(\widehat Hg_7)=0.
\]
Their original matrices remain
\[
 R_D=\operatorname{diag}(-1,1,1,1),\qquad
 H_D=\frac12\begin{pmatrix}1&1&1&1\\1&1&-1&-1\\
                           1&-1&1&-1\\1&-1&-1&1\end{pmatrix}.
\]
These maps preserve the original lattice, as can be checked on its
displayed generators.  Both \(R_D\) and \(H_D\) preserve \(D_4\):
the former changes one sign, while the latter has integral coordinates
on every even-sum integer vector and has output sum twice its first
input coordinate. They are orthogonal involutions. Their marked
discriminant actions fix \(v_D\) and exchange \(s_D,c_D\), and
exchange \(v_D,s_D\) while fixing \(c_D\), respectively.
On the binary glue generators, \(\widehat R\) fixes the first
two and sends the third to their second-plus-third sum;
\(\widehat H\) fixes the first and exchanges the second and third.
On the ternary code generators \(z_1=(1,1,1,0)\),
\(z_2=(1,2,0,1)\), the two signed component permutations act as
\((z_1,z_2)\mapsto(z_1,2z_2)\) and
\((z_1,z_2)\mapsto(z_1,2z_1+2z_2)\), respectively.
Thus both preserve the full original glue and its inverse image.
Every within-fibre coordinate permutation also preserves that inverse
image, since its action on \(\lambda+A_5\) is trivial:
\(Q\lambda-\lambda=Qe_0-e_0\in A_5\).
These are concrete failures to preserve this cubic, without
discarding their actual lattice actions or glue.

Negation preserves each lattice \(N\) and \(\Lambda\), and its
exact polynomial law is \(F(-X)=-F(X)\). In the neighbour the
nonzero witness is its displayed norm-four generator:
\begin{equation}
 F(\rho)=15\sqrt2+16\sqrt3,\qquad
 F(v)=\frac{14}{27}-\frac{7\sqrt2}{24}-\frac{\sqrt3}{27}.
 \label{nca:eq:special-values}
\end{equation}
Both equalities follow by substituting the original coordinates
of \(\rho\) and \(v\) into the proved polynomial.

Every original lattice isometry \(T\in O(N)\), whether or not
it preserves the cubic, has the exact replacement transport
\begin{equation}
 \Gamma_T=\mathcal LT\mathcal L^{-1},\qquad
 F^T(Y)=F(T^{-1}Y),\qquad F^T(TX)=F(X),
 \label{nca:eq:all-transport}
\end{equation}
and
\[
 T\Lambda_\rho=\Lambda_{T\rho},\qquad
 \mathcal V_{F^T}(\Lambda_{T\rho})=\mathcal V_F(\Lambda_\rho),
 \qquad \mathcal V_{F^T}(N)=\mathcal V_F(N).
\]
For the neighbour transport, write its exact elements as
\(x+k\rho/6\), where \(x\in N\) and
\((\rho,x)+k\equiv0\pmod6\); application of \(T\) gives
the identical congruence with \((T\rho,Tx)\). The inverse is
\(T^{-1}\). The value-module equalities then follow by the
bijections and the displayed identity \(F^T(TX)=F(X)\).
In particular for \(T=\sigma^{2k}\), one has \(F^T=F\), so
the same retained cubic has exactly the same value module on
every neighbour \(\Lambda_{\sigma^{2k}\rho}\) in this cyclic
orbit. This statement names its target neighbour and does not
assume the chosen \(\rho\) is fixed.

% Exact follow-on module. Input after the original Niemeier and Leech proofs.
% Original Euclidean coordinates, residue action, scales and glue are retained.
\section{The residue cubic action and the common integral Leech correspondence}
\label{esrf:sec:fixed-neighbor}

Retain the exact lattices $R=A^{\perp4}\perp D\subset N$ and the
binary and ternary codes from the index-72 construction.
The four $A$ coordinates remain the ordered six-element residue
fibres. To avoid confusing a vector with a representation, write
$\varrho$ for the Weyl vector
\[
 \varrho_A=(5,3,1,-1,-3,-5)/2,\qquad
 \varrho_D=(3,2,1,0),\qquad
 \varrho=\varrho_A^{\oplus4}\perp\varrho_D.
\]
Thus $\varrho\in N$ and $\varrho^2=84$.
The original residue representation is still denoted by $\rho$.
The specific element studied here is
\[
 r=\rho(361)=\rho(289)^2,\qquad
 r|_{A_i}=C^2,\quad Ce_j=e_{j+1\bmod6},\qquad r|_D=I.
 \tag{RF1}
\]
In particular $r^3=I$. This is the actual residue coordinate
permutation, with no conjugation or replacement.

For each $\sigma\in\Omega=O(N)\varrho$, choose an original root
$\alpha_\sigma$ with $(\sigma,\alpha_\sigma)=1$ and set
\[
 \begin{split}
 \ell_\sigma(x)&=(\sigma,x)\bmod6,\qquad
 K_\sigma=\ker(\ell_\sigma:N\to\mathbb Z/6),\\
 \Lambda_\sigma&=K_\sigma+
                      \mathbb Z(\sigma/6-\alpha_\sigma)\\
 &=\{x+k\sigma/6:x\in N,\ k\in\mathbb Z,\
                                    (\sigma,x)+k\equiv0\pmod6\}.
 \end{split}
 \tag{RF2}
\]
The preceding neighbor proof shows that the definition is independent
of $\alpha_\sigma$, every $\Lambda_\sigma$ is the specified rootless
even unimodular lattice, and
$N\cap\Lambda_\sigma=K_\sigma$ with both indices six.
The family considered in this section is precisely
$\mathscr F=\{\Lambda_\sigma:\sigma\in\Omega\}$.

\subsection{A necessary and sufficient equality criterion}

\begin{theorem}[Equality of the marked Weyl neighbors]
\label{esrf:thm:equality}
For $\sigma,\tau\in\Omega$, the following three statements are equivalent:
\[
 \Lambda_\sigma=\Lambda_\tau,\qquad
 K_\sigma=K_\tau,\qquad
 \tau-\epsilon\sigma\in6N
                 \text{ for some }\epsilon\in\{1,-1\}.
 \tag{RF3}
\]
\end{theorem}

\begin{proof}
Equality of the two neighbors gives equality of the kernels by
intersection with $N$.
Each $\ell_\sigma$ is onto because
$(\sigma,\alpha_\sigma)=1$.
If the kernels are equal, their induced isomorphisms
$N/K_\sigma\to\mathbb Z/6$ differ by an automorphism of
$\mathbb Z/6$. Its two units are $1,-1$.
Therefore $(\tau-\epsilon\sigma,N)\subset6\mathbb Z$.
Unimodularity $N^\vee=N$ gives
$(\tau-\epsilon\sigma)/6\in N$.
Conversely this congruence makes the two kernels equal.

It remains to show that this kernel criterion actually gives
equality of the chosen unimodular extensions.
The exact formula (RF2), after replacing $k$ by $-k$, gives
$\Lambda_{-\sigma}=\Lambda_\sigma$.
We may consequently suppose $\tau=\sigma+6t$ with $t\in N$.
Both Weyl vectors have squared norm 84, so
\[
 0=\tau^2-\sigma^2=12(\sigma,t)+36t^2,\qquad
                       (\sigma,t)=-3t^2\in6\mathbb Z.
 \tag{RF4}
\]
The last inclusion uses the actual evenness of $N$.
For $x+k\tau/6$ in (RF2), put $x'=x+kt\in N$.
Then $x+k\tau/6=x'+k\sigma/6$, and
\[
 (\sigma,x')+k
 \equiv(\sigma,x)+k
 \equiv(\tau,x)+k\equiv0\pmod6.
\]
This proves $\Lambda_\tau\subset\Lambda_\sigma$.
Interchanging $\sigma,\tau$ proves the reverse inclusion.
Thus no additional unverified choice of an extension is being
identified merely from equality of its common kernel.
\end{proof}

It follows in particular that the map
\[
 \Omega/{\sim}\ \longrightarrow\ \mathscr F,\qquad
 [\sigma]\longmapsto\Lambda_\sigma,\qquad
 \sigma\sim\tau\ \Longleftrightarrow\
                          \tau\equiv\pm\sigma\pmod{6N},
 \tag{RF5}
\]
is a bijection. It is $O(N)$-equivariant: the original formula
$T\Lambda_\sigma=\Lambda_{T\sigma}$ follows by substitution in
(RF2), and $T$ preserves the congruence relation since $T(N)=N$.

\subsection{The exact tetracode obstruction to a fixed neighbor}

\begin{theorem}[No fixed member of this specified family]
\label{esrf:thm:no-fixed-neighbor}
The actual residue element $r=\rho(361)$ fixes no lattice of
$\mathscr F$. Every orbit of its action on $\mathscr F$ has
exactly three members.
\end{theorem}

\begin{proof}
Suppose first that $r\sigma\equiv-\sigma\pmod{6N}$.
Applying $r$ successively, using $r(N)=N$ and $r^3=I$, gives
$\sigma\equiv-\sigma\pmod{6N}$.
Then $\sigma/3\in N$, contradicting its pairing
$1/3$ with the root $\alpha_\sigma\in N$.
Thus (RF3) permits fixedness only if
$t=(r\sigma-\sigma)/6\in N$.

Every isometry of $N$ preserves the original root lattice, permutes
its four $A_5$ components, and acts in each of them by a signed
coordinate permutation, as proved in the full automorphism theorem.
Negation permutes the six entries of $\varrho_A$.
Consequently each $A$ block $s=(s_0,\ldots,s_5)$ of every
$\sigma\in\Omega$ is a permutation of the six distinct numbers
$5/2,3/2,1/2,-1/2,-3/2,-5/2$.
In this block put $\Delta_j=s_{j-2\bmod6}-s_j$.
If the corresponding block $\Delta/6$ of $t$ belongs to $A^\vee$,
its coordinate differences are integral, so
\[
 \Delta_j\equiv c\pmod6\qquad\text{for every }j
 \tag{RF6}
\]
for one integer $c$ modulo six.
On either three-cycle of the index permutation, the three
$\Delta_j$ sum to zero. Hence $3c\equiv0\pmod6$ and
$c\in\{0,2,4\}$.
Each $\Delta_j$ is an integer between $-5$ and $5$.
If $c=0$, (RF6) forces all $\Delta_j=0$, contrary to the
distinctness of the entries on each three-cycle.
Therefore $c=2$ or $4$.

In the original discriminant marking, a vector of the $A$ class
$t_i\lambda+A$ has common fractional part $-t_i/6$.
The corresponding class of $\Delta/6$ is therefore
$t_i=-c\pmod6$, also in $\{2,4\}$.
The difference on the $D$ block is exactly zero because $r|_D=I$.
Thus membership of the full vector $(r\sigma-\sigma)/6$ in $N$
would require a glue word
\[
 (t_1,t_2,t_3,t_4;00)\in H,\qquad t_i\in\{2,4\}
                                                \text{ for every }i.
 \tag{RF7}
\]
In $H=\{(3u+2z;y)\}$ this forces $u=0$ and all four $z_i$
nonzero. But the exact ternary code is
\[
 C_3=\{(a+b,a+2b,a,b):a,b\in\mathbb F_3\}.
 \tag{RF8}
\]
If exactly one parameter is zero, exactly one coordinate is zero.
If both are nonzero, precisely one of $a+b,a+2b$ is zero.
Thus every nonzero word has weight three, and the zero word
has weight zero. No word in (RF7) belongs to $H$.
This contradicts $t\in N$.

Neither sign in (RF3) can occur. Since $r^3=I$, every orbit
therefore has cardinality three.
\end{proof}

The conclusion concerns exactly $\mathscr F$ and the action (RF1).
The integral correspondence for these three-member orbits is
constructed next; no assertion about other embeddings or
constructions of a Leech lattice follows from this obstruction.

\subsection{The canonical orbit and its common integral lattice}

Put $\sigma_j=r^j\varrho$, $\Lambda_j=\Lambda_{\sigma_j}$ and
$K_j=K_{\sigma_j}$, with $j$ read modulo three.
Keep the same original root
$\alpha=e_0-e_1$ in the first $A$ block.
Directly $(\sigma_j,\alpha)=1$ for all three $j$.
Define
\[
 v_j=\sigma_j/6-\alpha,\quad v=v_0,\qquad
 q_1=(\sigma_1-\sigma_0)/6,\quad
 q_2=(\sigma_2-\sigma_0)/6.
 \tag{RF9}
\]
The complete original coordinates are
\[
 \begin{split}
 q_1&=\bigl((-2,-2,1,1,1,1)/3\bigr)^{\oplus4}\perp0,\\
 q_2&=\bigl((-1,-1,-1,-1,2,2)/3\bigr)^{\oplus4}\perp0.
 \end{split}
 \tag{RF10}
\]
Their classes in $R^\vee/R$ are respectively
$(4,4,4,4;00)$ and $(2,2,2,2;00)$.
Neither is in $H$, while $3q_1,3q_2\in R$.
In particular each $q_i+N$ has exact order three.
Direct pairings give
\[
 q_1^2=q_2^2=16/3,\quad(q_1,q_2)=8/3,\quad
                    (\varrho,q_1)=(\varrho,q_2)=-16,
 \qquad q_1+q_2\in R.
 \tag{RF11}
\]

For $x\in N$, write its original glue class as $(3u+2z;y)$,
with $z=(a+b,a+2b,a,b)\in C_3$, and define the homomorphism
\[
 \beta:N\longrightarrow\mathbb F_3,\qquad\beta(x)=b=z_4.
 \tag{RF12}
\]
This map is defined through the original quotient $N/R$.
This definition uses the retained residue marking without changing it.
The sum of the four $z$ coordinates is also $b$ modulo three.
Write $\ell_j(x)=(\sigma_j,x)\bmod6$.
Then the exact three character formulas are
\[
 \ell_1(x)=\ell_0(x)+4\beta(x),\qquad
 \ell_2(x)=\ell_0(x)+2\beta(x)\quad\text{in }\mathbb Z/6.
 \tag{RF13}
\]
The maps $b\mapsto2b,4b$ from $\mathbb F_3$ to $\mathbb Z/6$
are well-defined homomorphisms.
To prove (RF13), the discriminant bilinear form gives
\[
 (q_1,x)\equiv\tfrac56\,4\sum t_i
       \equiv\tfrac13\sum t_i
       \equiv\tfrac23\beta(x)\pmod{\mathbb Z},
 \qquad t_i=3u_i+2z_i\pmod6.
\]
For $q_2$ the same calculation gives
$(q_2,x)\equiv\beta(x)/3\pmod{\mathbb Z}$.
Multiplication by six and (RF9) prove both formulas.

Using the original dominant representatives
$\omega_t=(1^t,0^{6-t})-\frac t6(1,\ldots,1)$, put
\[
 u_*=(\omega_2,\omega_4,0,\omega_2;0)\in N.
 \tag{RF14}
\]
Its ternary class is exactly the original generator
$z_2=(1,2,0,1)$.
It has $u_*^2=4$, $(\varrho,u_*)=12$ and $\beta(u_*)=1$.
The map
\[
 (\ell_0,\beta):N\longrightarrow\mathbb Z/6\oplus\mathbb F_3
 \tag{RF15}
\]
is onto: $\alpha$ maps to $(1,0)$ and $u_*$ to $(0,1)$.
Consequently
\[
 \begin{gathered}
 I=\ker(\ell_0,\beta)=K_0\cap K_1=K_1\cap K_2
     =K_2\cap K_0=\bigcap_jK_j,\\
 [N:I]=18,\qquad[K_j:I]=3,\qquad\det I=324.
 \end{gathered}
 \tag{RF16}
\]
The kernel equalities follow from (RF13), because either
$2b=0$ or $4b=0$ in $\mathbb Z/6$ forces $b=0$ in
$\mathbb F_3$.
The action $r$ permutes the $K_j$, so it preserves $I$.
Also $3q_i\in I$: these are in $R$, hence have $\beta=0$,
and $(\varrho,3q_i)=-48$ is divisible by six.

Define the full-rank lattice
\[
 J=I+\mathbb Z(3v).
 \tag{RF17}
\]
It is contained in $\Lambda_0$.
The element $6v=\varrho-6\alpha$ lies in $I$: its
$\ell_0$ value is zero, and its glue has zero ternary part.
The element $3v$ is not in $N$, because
$(3v,\alpha)=-11/2$ is not an integer.
Thus
\[
 N\cap J=I,\qquad[J:I]=2,\qquad
                    \det J=324/4=81,\qquad[\Lambda_0:J]=9.
 \tag{RF18}
\]
This lattice is positive definite, even, and rootless, because
it is a sublattice of $\Lambda_0$.

All exact root corrections in the cyclic action are retained.
Put $\kappa_j=\alpha-r^j\alpha\in R$.
The vectors $r\alpha=e_2-e_3$ and $r^2\alpha=e_4-e_5$
are original simple roots with $\varrho$-pairing one.
Hence $\kappa_j\in I$.
The precise formulas, rather than literal identification of
different generators, are
\[
 r^jv=v+q_j+\kappa_j\quad(j=1,2),\qquad
 r(3v)=3v+3q_1+3\kappa_1.
 \tag{RF19}
\]
Since $rI=I$, $3q_1\in I$ and $\kappa_1\in I$, this proves
$rJ\subset J$. Applying $r^3=I$ gives equality $rJ=J$.

\begin{theorem}[All three pairwise Leech intersections]
\label{esrf:thm:common-J}
The actual three-member orbit satisfies
\[
 \Lambda_0\cap\Lambda_1
 =\Lambda_1\cap\Lambda_2
 =\Lambda_2\cap\Lambda_0
 =\bigcap_j\Lambda_j=J,\qquad[\Lambda_j:J]=9.
 \tag{RF20}
\]
The element $r$ is an integral isometry of $J$, and induces
isometries $\Lambda_j\to\Lambda_{j+1}$ with this common restriction.
\end{theorem}

\begin{proof}
First $J\subset\Lambda_0\cap\Lambda_1$: $I$ lies in both,
and $3v=3v_1-3q_1$ with $3q_1\in I$.
Suppose $y\in\Lambda_0\cap\Lambda_1$.
Using the same $\alpha$, write
$y=x_0+k_0v_0=x_1+k_1v_1$, $x_j\in K_j$.
The fractional part of $(v_j,\alpha)$ is $1/6$ for both
$j$, so pairing with $\alpha\in N$ gives
$k_0\equiv k_1\pmod6$.
Absorbing multiples of $6v_j\in K_j$, take the same
$k\in\{0,\ldots,5\}$ in both expressions.
Then $x_0-x_1=kq_1\in N$.
The exact order three of $q_1+N$ forces $k=0$ or $3$.
For these values $kq_1\in I\subset K_1$, so
$x_0=x_1+kq_1\in K_0\cap K_1=I$.
Thus $y=x_0+kv\in J$, proving the first equality.
Apply $r$ and $r^2$, using $rJ=J$, to prove the other pairwise
equalities. Their common index is nine by (RF18) and the
isometries induced by $r$.
\end{proof}

For comparison with the original index-72 glue, let
$R_0=R\cap K_0$. Since $\beta|_R=0$, (RF13) gives
$R\cap K_j=R_0$ for every $j$, and $R\cap I=R_0$.
The reduction map $K_0/R_0\cong H$ identifies $I/R_0$
with the kernel of the original ternary coefficient $b$.
Thus
\[
 [I:R_0]=24,\qquad[J:R_0]=48,\qquad
 [K_0:R_0]=72,\qquad[\Lambda_j:R_0]=432.
 \tag{RF21}
\]
All these are inclusions of the displayed original lattices.

\subsection{The exact discriminant form and its cyclic three planes}

\begin{theorem}[The integral correspondence over $\mathbb F_3^4$]
\label{esrf:thm:discriminant-correspondence}
One has
\[
 J^\vee=\Lambda_i+\Lambda_j\quad(i\ne j),\qquad
 J^\vee/J\cong\mathbb F_3^4.
 \tag{RF22}
\]
In the ordered basis
\[
 e_1=u_*+J,\quad e_2=v+J,\quad
 e_3=ru_*+J,\quad e_4=rv+J,
 \tag{RF23}
\]
the discriminant quadratic form is
\[
 q_J(a,b,c,d)=\frac23(2ad+bc+bd)\pmod{2\mathbb Z},
 \tag{RF24}
\]
and three times its discriminant bilinear form has matrix
\[
 B=
 \begin{pmatrix}
 0&0&0&2\\0&0&1&1\\0&1&0&0\\2&1&0&0
 \end{pmatrix}\quad\text{over }\mathbb F_3.
 \tag{RF25}
\]
The three original Leech extensions are exactly the inverse
images in $J^\vee$ of the isotropic planes
\[
 \begin{split}
 P_0&=\{(a,b,0,0):a,b\in\mathbb F_3\},\\
 P_1&=\{(0,0,c,d):c,d\in\mathbb F_3\},\\
 P_2&=\{(a,b,a-b,b):a,b\in\mathbb F_3\}.
 \end{split}
 \tag{RF26}
\]
The induced action of the actual $r$ is the matrix
\[
 \mathcal R=
 \begin{pmatrix}
 0&0&2&1\\
 0&0&0&2\\
 1&0&2&2\\
 0&1&0&2
 \end{pmatrix},\qquad
 \mathcal R^3=I,\quad \mathcal R^TB\mathcal R=B,
 \quad\mathcal R P_j=P_{j+1}.
 \tag{RF27}
\]
All entries in (RF23)--(RF27) use the original glue generator
$z_2$, root $\alpha$, pairing, and residue permutation.
\end{theorem}

\begin{proof}
Since $J\subset\Lambda_i=\Lambda_i^\vee$, each $\Lambda_i$
lies in $J^\vee$.
By (RF20), the two subgroups $\Lambda_i/J$ and $\Lambda_j/J$
have order nine and intersect trivially for $i\ne j$.
Their sum has order 81, equal to
$|J^\vee/J|=\det J=81$, proving the first equality in (RF22).

The homomorphism $\beta:K_0/I\to\mathbb F_3$ is an isomorphism,
with generator $u_*+I$. Hence $\Lambda_0/J$ is generated by
$u_*+J,v+J$.
Both have order dividing three because $3u_*\in I$ and
$3v\in J$.
They are independent: if $a u_*+b v\in J$, reduce first modulo
$N$; the subgroup $J+N$ has the two cosets represented by
$0,3v$, whereas $v+N$ has order six.
For $a,b\in\{0,1,2\}$ this forces $b=0$.
Then membership of $a u_*$ in $I=N\cap J$ forces $a=0$.
Thus $\Lambda_0/J\cong\mathbb F_3^2$ with the first two
generators in (RF23).
Applying $r$ proves the corresponding statement for
$\Lambda_1/J$ and the last two generators.
Their trivial intersection and total order 81 prove that
(RF23) is indeed a basis of the full discriminant group.

Here are the full Euclidean Gram entries before any passage
to a quotient:
\[
 \operatorname{Gram}(u_*,v,ru_*,rv)=
 \begin{pmatrix}
 4&2&-2&-4/3\\
 2&4&-2/3&-2/3\\
 -2&-2/3&4&2\\
 -4/3&-2/3&2&4
 \end{pmatrix}.
 \tag{RF28}
\]
For an explicit verification, each of the three nonzero
$A$ blocks of $u_*$ pairs to $-2/3$ with its shifted block,
giving $(u_*,ru_*)=-2$.
One has
$(u_*,\alpha)=(u_*,r\alpha)=(ru_*,\alpha)=0$,
$(\varrho,u_*)=12$,
$(\sigma_1,u_*)=-8$, and
$(\sigma_2,u_*)=-4$, directly from (RF10) and (RF14).
These give respectively
$(u_*,v)=2$, $(u_*,rv)=-4/3$, and
$(v,ru_*)=-2/3$.
Also $(\varrho,\sigma_1)=-12$,
$(\varrho,r\alpha)=(\sigma_1,\alpha)=1$,
and $(\alpha,r\alpha)=0$. Thus
\[
 (v,rv)=-12/36-1/6-1/6=-2/3.
\]
The four diagonal entries equal four, and
$(ru_*,rv)=(u_*,v)=2$ by the exact isometry $r$.
This proves (RF28).
Reducing its entries modulo $\mathbb Z$ gives (RF25).
The diagonal and integral cross terms in the squared norm
are even. Reducing the remaining cross terms modulo
$2\mathbb Z$ gives exactly (RF24).
Changing any coordinate by three changes (RF24) by an even
integer, so it is well-defined on the asserted field basis.

It remains to calculate the action without identifying a
family permutation with a fixed Leech automorphism.
By definition $r e_1=e_3$, $r e_2=e_4$.
The vector $u_*+ru_*+r^2u_*$ lies in $I$: its $\beta$ value
is $1+1+1=0$, because $r$ acts trivially on $N/R$,
and its $\ell_0$ value is $12-4-8=0$.
Consequently
\[
 r^2u_*\equiv-u_*-ru_*\pmod J.
 \tag{RF29}
\]
By (RF19),
\[
 v+rv+r^2v
       =3v+q_1+q_2+\kappa_1+\kappa_2.
\]
The $\kappa_j$ belong to $I$, and $3v\in J$.
Furthermore $q_1+q_2-u_*+ru_*\in I$.
Indeed $q_1+q_2\in R$, so its $\beta$ value is zero;
subtracting $u_*$ and adding $ru_*$ keeps that value zero.
Its pairing with $\varrho$ is
$-32-12-4=-48$, divisible by six.
Thus
\[
 r^2v\equiv u_*-v-ru_*-rv\pmod J.
 \tag{RF30}
\]
The four images in (RF29)--(RF30) give the four columns of
$\mathcal R$ in (RF27).
The matrix has cube the identity and preserves $B$ either
by multiplication, or by the already proved facts $r^3=I$
and $rJ=J$ with its original isometric pairing.

The first two planes in (RF26) are respectively
$\Lambda_0/J$ and $\Lambda_1/J$ by their proved bases.
Their next image is the span of
$(-1,0,-1,0)$ and $(1,-1,-1,-1)$, from (RF29)--(RF30).
Solving their two linear parameters gives exactly
$(a,b,a-b,b)$, proving the stated $P_2=\Lambda_2/J$.
Each plane is isotropic: the first two vanish in (RF24),
and the third has
$2ab+b(a-b)+b^2=3ab=0$ in $\mathbb F_3$.
They have nine elements and pairwise intersection zero,
also following directly from (RF26).
Finally each $\Lambda_j$ is exactly the inverse image of
its quotient plane because it contains $J$.
This proves the claimed integral correspondence and all
its maps, indices, pairings, and cyclic actions.
\end{proof}

\subsection{The fixed fourth plane and its exact \texorpdfstring{$A_2^{12}$}{A2 to the twelfth power} extension}

The same finite discriminant form has the additional plane
\[
 P_\infty=\{(a,b,-a+b,-b):a,b\in\mathbb F_3\}.
 \tag{RF31}
\]
Put
\[
 p=u_*-ru_*,\qquad q=v+ru_*-rv,\qquad
                       M=J+\mathbb Zp+\mathbb Zq.
 \tag{RF32}
\]
These are literal vectors in the original Euclidean space.

\begin{theorem}[The fixed fourth extension]
\label{esrf:thm:fourth-extension}
The plane $P_\infty$ is isotropic and is fixed by the induced
residue action. Its inverse image $M$ is a positive definite
even unimodular lattice with exactly 72 roots, of root system
$A_2^{12}$ and minimum squared norm two.
The original residue element $r$ is an integral isometry of $M$.
It fixes four of these $A_2$ components pointwise and acts as
a coordinate three-cycle on each of the other eight.

The exact intersection and index relations are
\[
 \begin{gathered}
 M\cap\Lambda_j=J,\qquad [M:J]=[\Lambda_j:J]=9,\\
 M\cap N=L_0:=\{x\in N:\beta(x)=0,\
                                (\varrho,x)\equiv0\pmod2\},\\
 [M:L_0]=[N:L_0]=6.
 \end{gathered}
 \tag{RF33}
\]
If $R_M$ denotes the lattice generated by all its roots, then
\[
 R_M\cong A_2^{\perp12},\qquad
 \det R_M=3^{12}=531441,\qquad[M:R_M]=3^6=729.
 \tag{RF34}
\]
\end{theorem}

\begin{proof}
Substitution of (RF31) in (RF24) gives the numerator
$2a(-b)+b(-a+b)+b(-b)=-3ab=0$ in $\mathbb F_3$.
Thus $P_\infty$ is isotropic.
The matrix (RF27) maps its vector with parameters $(a,b)$ to
\[
 (a+b,b,-a,-b),
 \tag{RF35}
\]
which is again in $P_\infty$, now with parameters $(a+b,b)$.
The images of $p,q$ in (RF23) are
$(1,0,-1,0)$ and $(0,1,1,-1)$, a basis of $P_\infty$.
Consequently (RF32) is its exact inverse image under
$J^\vee\to J^\vee/J$, and $[M:J]=9$.
Isotropy gives even norms on $M$; polarization gives integral
pairings. Its determinant is $\det J/9^2=1$, and its
positive definiteness and rank are inherited from the unchanged
ambient real space. Equation (RF35), together with $rJ=J$,
proves $rM=M$.
Each of $P_0,P_1,P_2$ intersects $P_\infty$ in zero, as follows
by equating their coordinates in (RF26) and (RF31).
Taking inverse images proves $M\cap\Lambda_j=J$.

For its relation with $N$, (RF19) gives the exact expression
\[
 q=ru_*-q_1-\kappa_1 .
 \tag{RF36}
\]
Therefore
\[
 N+M=N+\mathbb Z(\varrho/2)+\mathbb Zq_1.
 \tag{RF37}
\]
The two indicated cosets modulo $N$ have orders two and three.
For the first, pair with $\alpha$; for the second, use (RF10).
Subgroups of coprime orders intersect trivially, so the quotient
in (RF37) has order six.
Also $p\in N$ has $\beta(p)=0$ and
$(\varrho,p)=12-(-4)=16$, hence image $(4,0)$ in (RF15).
Thus $p+I$ has order three.
The vector $3q$ lies in $J$ because $q+J$ has order three,
and lies in $N$ by (RF36) and $3q_1\in N$.
It is therefore in $N\cap J=I$.
We also have $6v\in I$.
Write an element of $M$ using generators $I,3v,p,q$ and reduce
the coefficient of $3v$ modulo two and of $q$ modulo three.
If this element is in $N$, their two coprime-order cosets in
(RF37) force both reduced coefficients to vanish.
Hence
\[
 M\cap N=I+\mathbb Zp.
 \tag{RF38}
\]
Its image in (RF15) is precisely $2(\mathbb Z/6)\oplus0$,
which proves the displayed formula for $L_0$ and
$[N:L_0]=6$. Since $\det N=\det M=1$, the equal covolumes
give $[M:L_0]=6$, proving (RF33).

We next prove that the following root count is exhaustive.
Every vector in $N+M$ has one of the six forms
\[
 x+a\varrho/2-bq_1,\qquad
                   x\in N,\ a\in\{0,1\},\ b\in\{0,1,2\}.
 \tag{RF39}
\]
When $a=0$, the vector is in $R^\vee$.
Every norm-two vector of this particular $R^\vee$ is already
a root of $R$.
Indeed, the $A$ class minima are
$0,5/6,4/3,3/2,4/3,5/6$, and the three nonzero $D$ classes
have minimum one. In a nonzero class of $R^\vee/R$ with
minimum at most two, the possible minima are exactly among
\[
 \{5/6,1,4/3,3/2,5/3,11/6\}.
 \tag{RF40}
\]
This exhausts them because at most two $A$ classes can be
nonzero under the bound; if two are nonzero, each has minimum
$5/6$; with a nonzero $D$ class there is at most one nonzero
$A$ class and its minimum must be $5/6$.
No number in (RF40) is congruent to two modulo $2\mathbb Z$.
All norms within a fixed discriminant class differ by an even
integer, so no such class has a norm-two vector.
This proves the claim about $R^\vee$.

All original roots have $\beta=0$, so (RF33) says that those
lying in $M$ are exactly the ones with even $\varrho$ height.
In each original $A$ block they are
$e_i-e_j$ with $i\ne j$ and $i\equiv j\pmod2$.
These give two orthogonal $A_2$ systems per block, on the
coordinate triples $\{0,2,4\}$ and $\{1,3,5\}$, and 48 roots
in total. In $D$ they are the eight roots
$\pm e_1\pm e_3,\ \pm e_2\pm e_4$.
They form four pairwise orthogonal $A_1$ systems.
This accounts for exactly 56 roots with $a=0$ in (RF39).

For $a=1$, absorb the $A$ dual vector $-bq_1$ into the
marked $A$ discriminant classes.
If the original class of $x$ is $(t_1,t_2,t_3,t_4;y)\in H$,
the shifted class in each $A$ is $t_i-4b$.
The exact half-Weyl-shift minima proved in the neighbor
calculation are
\[
 \min_{w\in t\lambda+A}|w+\varrho_A/2|^2
 =\begin{cases}3/8& t\equiv0\pmod3,\\17/24&t\not\equiv0\pmod3,\end{cases}
 \qquad
 \min_{w+D=y}|w+\varrho_D/2|^2=1/2.
 \tag{RF41}
\]
For completeness the first line follows by choosing centered
quarter coordinates when $t=0,3$; for other $t$, the nearest
twelfth coordinates have square sum $13/24$ and coordinate
sum $\pm1$, and restoring zero sum costs $1/6$.
The second line follows from the integer coordinate distances
$(1/2,0,1/2,0)$ and half-integer distances
$(0,1/2,0,1/2)$, with their ties meeting every $D$ marking.
Thus the total minimum in (RF39) with $a=1$ is at least two.
Equality requires $t_i-4b\equiv0\pmod3$ in every component.
Since $t_i\equiv2z_i\pmod3$, this says
$z_i=2b$ for all four coordinates.
For $b=1,2$ this is excluded by the exact code (RF8).
For $b=0$ it requires $z=0$.
The eight binary codewords and the two $D$ ties then give
exactly sixteen equality vectors $y=x+\varrho/2$.
Every one belongs to $M$: the norm equation gives
\[
 2=x^2+(\varrho,x)+21,\qquad
                         (\varrho,x)=-19-x^2,
 \tag{RF42}
\]
which is odd because $N$ is even.
Also $\beta(x)=0$ because $z=0$.
Therefore $x+3\alpha\in L_0$ and
$y=(x+3\alpha)+3v\in M$.
This proves exhaustively that $M$ has exactly $56+16=72$ roots.

Here are the full root coordinates and the resulting exact
orthogonal decomposition.
Let $\zeta_i$ be supported in the $i$th original $A$ block
with entries
\[
 \zeta=(1,-1,1,-1,1,-1)/4,\qquad \zeta^2=3/8.
\]
In the following table every $Z_y$ is an actual sum of those
four mutually orthogonal vectors, and every $\delta_y$ is
an original $D$ root:
\[
 \begin{array}{c|l|l}
 y&Z_y&\delta_y\\\hline
 00&\zeta_1+\zeta_2+\zeta_3+\zeta_4&e_1+e_3\\
 10&-\zeta_1-\zeta_2+\zeta_3+\zeta_4&e_1-e_3\\
 01&-\zeta_1+\zeta_2-\zeta_3+\zeta_4&e_2+e_4\\
 11&\zeta_1-\zeta_2-\zeta_3+\zeta_4&e_2-e_4
 \end{array}.
 \tag{RF43}
\]
The sixteen new roots are exactly
\[
 \{\epsilon Z_y+\eta\delta_y/2:
                   y\in\mathbb F_2^2,\ \epsilon,\eta\in\{1,-1\}\}.
 \tag{RF44}
\]
To verify the signs from the original glue, the $A$ half-shift
minimizer for binary coordinate $u_i=0$ is $\zeta_i$ and
for $u_i=1$ it is $-\zeta_i$.
For $y=00$ the two codewords are $0000,1111$; for $10$ they
are $1100,0011$; for $01$ they are $1010,0101$; and for $11$
they are $0110,1001$.
These give respectively the four opposite pairs $\pm Z_y$
in (RF43).
The two shifted $D$ minimizers in these exact original
markings are respectively
$\pm(e_1+e_3)/2,\pm(e_1-e_3)/2,
 \pm(e_2+e_4)/2,\pm(e_2-e_4)/2$.
This proves (RF44) with every sign and marking retained.

The four $Z_y$ are mutually orthogonal and have squared norm
$3/2$, as follows by dotting their four displayed sign rows.
The $\delta_y$ are mutually orthogonal of squared norm two.
The $Z_y$ are orthogonal to the entire original $D$ space,
and to the eight old $A_2$ systems because their coordinates
are constant on each parity triple.
For each $y$, the six roots
\[
 \{\pm\delta_y\}
       \ \cup\ \{\epsilon Z_y+\eta\delta_y/2:\epsilon,\eta=\pm1\}
 \tag{RF45}
\]
are an $A_2$ system: take simple roots
$\delta_y$ and $Z_y-\delta_y/2$, whose Gram matrix is
$\left(\begin{smallmatrix}2&-1\\-1&2\end{smallmatrix}\right)$;
their sum is $Z_y+\delta_y/2$.
The four such systems are mutually orthogonal and orthogonal
to the eight old systems. They exhaust the 72 proved roots.
Thus their root lattice is $A_2^{\perp12}$ of determinant
$3^{12}$, and unimodularity of $M$ gives the index $3^6$.
This proves (RF34) and the Coxeter number three of this
rooted Niemeier lattice directly from its twelve components.

Finally $r$ fixes each $\zeta_i$ because it permutes coordinates
within their parity triples; it fixes each $\delta_y$ because
it is the identity on $D$. Hence it fixes the four systems
(RF45) pointwise. On each of the eight old $A_2$ systems it
cyclically permutes its three coordinate positions, giving its
exact order-three Coxeter action.
This completes the asserted classification of the fixed
fourth extension and its precise relation to the three Leech
lattices through $J$.
\end{proof}

% New J/M extension; prerequisite definitions are in es_niemeier_cubic_arithmetic.tex.
\section{The cubic on the common lattice and the fourth Niemeier extension}
\label{jmv:sec:JM-extension}

Retain the original coordinates, Cayley multiplication, $F$, the generators $g_i$
and their heights $H_i$ from Section~\ref{nca:sec:cubic-arithmetic}, and the
exact common lattice and fourth extension of Section~\ref{esrf:sec:fixed-neighbor}.
The symbol $u$ below is the unchanged vector $u_*$ in (RF32).
\subsection{Integral generators of the common lattice and the fixed extension}

Define the unchanged residue action by
\[
 (rX)_{i,j}=x_{i,j-2\pmod6},\qquad (rX)_{D,k}=d_k.
\]
The two primary parts of the original glue have the exact descriptions
\begin{align*}
 H_2={}&\{(\varepsilon_1+\varepsilon_2+\varepsilon_3,
 \varepsilon_1+\varepsilon_2,\varepsilon_1+\varepsilon_3,
 \varepsilon_1;\varepsilon_2,\varepsilon_3):\varepsilon_i\in\mathbb F_2\},\\
 H_3={}&\{(t+s,t+2s,t,s;00):t,s\in\mathbb F_3\}.
\end{align*}
In the original \((\mathbb Z/6)^4\) coordinates this means
\(k=3w+2z\), with \(w\) the four binary entries of \(H_2\)
and \(z\) the four ternary entries of \(H_3\). The binary
\(D_4\) marking is retained. Chinese remainders recover \(w,z\)
uniquely, and the displayed parameterizations are injective.
Consequently
\[
 \beta:N\longrightarrow\mathbb F_3,\qquad \beta(X)=z_4=s
\]
is a well-defined additive homomorphism. The actual lifts
\eqref{nca:eq:N-generators} satisfy
\[
 \beta(g_i)=0\quad(1\leq i\leq28),\qquad \beta(g_{29})=1.
\]
Indeed roots have zero glue, the first three lifts are two-primary,
and the last two have \((t,s)=(1,0),(0,1)\), respectively.

Set, with no change of representatives,
\begin{align}
 I&=\{X\in N:(\rho,X)\equiv0\pmod6,\ \beta(X)=0\},
       &J&=I+\mathbb Z(3v),\label{jmv:eq:I-J}\\
 \omega_t&=(\underbrace{1,\ldots,1}_{t},0,\ldots,0)
                    -\frac t6(1,1,1,1,1,1),
       &u&=(\omega_2,\omega_4,0,\omega_2;0),\notag\\
 p&=u-r(u),&q&=v+r(u)-r(v),\notag\\
 M&=J+\mathbb Zp+\mathbb Zq.&&\label{jmv:eq:M}
\end{align}
Thus \(J\) and \(M\) are precisely the common-lattice and
residue-fixed extension presentations being evaluated here; the
following value-group proof needs no root count, lattice determinant,
or classification label for \(M\). For explicitness, the two extra
vectors in the original five blocks are
\begin{align*}
 p={}&((1,1,-1,-1,0,0),(1,1,0,0,-1,-1),0,
                         (1,1,-1,-1,0,0);0),\\
 q={}&\tfrac13((-2,4,4,-2,-2,-2),0,(2,2,-1,-1,-1,-1),
                         (1,1,1,1,-2,-2);0).
\end{align*}
The zero blocks have their original sizes six and four. These
coordinates follow by substituting \(v=\rho/6-\alpha\) into
\eqref{jmv:eq:M}; in particular the shifted \(\alpha\) term in
\(r(v)\) is retained. Also \(u\in N\),
\(\beta(u)=1\), \((\rho,u)=12\), and \((u,u)=4\):
its class is \((2,4,0,2;00)\),
\((\rho_A,\omega_t)=t(6-t)/2\), and
\((\omega_t,\omega_t)=t(6-t)/6\).

\begin{theorem}[Integral generating lists with both kernel conditions]
\label{jmv:thm:IJM-generators}
Define the following 32 ordered vectors:
\begin{align}
 b_1&=6\alpha,&
 b_i&=g_i-H_i\alpha\quad(2\leq i\leq28),\notag\\
 b_{29}&=3g_{29}-60\alpha,&
 b_{30}&=3v=\rho/2-3\alpha,\notag\\
 b_{31}&=p,&b_{32}&=q.\label{jmv:eq:IJM-generators}
\end{align}
Then the integral, not merely rational, equalities are
\[
 I=\sum_{i=1}^{29}\mathbb Zb_i,\qquad
 J=\sum_{i=1}^{30}\mathbb Zb_i,\qquad
 M=\sum_{i=1}^{32}\mathbb Zb_i.
\]
\end{theorem}
\begin{proof}
Every one of the first 29 vectors has \(\beta=0\) and height
divisible by six, by the exact heights and characters already
calculated. Conversely take any integral expression
\(X=\sum_{i=1}^{29}n_ig_i\) for \(X\in I\). Applying
\(\beta\) shows \(n_{29}=3a\) for an integer \(a\),
regardless of which expression was chosen. The exact equality
\[
 X=\sum_{i=2}^{28}n_i(g_i-H_i\alpha)
        +a(3g_{29}-60\alpha)
        +\left(n_1+\sum_{i=2}^{28}n_iH_i+60a\right)\alpha
\]
has last coefficient \((\rho,X)\), which is divisible by six.
It therefore expresses \(X\) as an integer combination of the
first 29 \(b_i\). The reverse containment was proved for each
generator. The last two equalities follow by adjoining exactly
the vectors in \eqref{jmv:eq:I-J} and \eqref{jmv:eq:M}.
No assertion of linear independence is used.
\end{proof}

\subsection{Complete finite certificates for the two new value groups}

\begin{theorem}[Equality of the common-lattice and fixed-extension value-generated groups]
\label{jmv:thm:JM-values}
For the retained Cayley multiplication and the original ordered
coordinates, the two additive cubic-value groups are exactly
\begin{equation}
 \mathcal V_F(J)=\mathcal V_F(M)
 =\frac1{27}\mathbb Z\ \oplus\frac{\sqrt2}{216}\mathbb Z
                         \ \oplus\frac{\sqrt3}{18}\mathbb Z.
 \label{jmv:eq:JM-values}
\end{equation}
The equality concerns generated additive groups, not equality
of individual-value sets \(F(J)\) and \(F(M)\).
\end{theorem}
\begin{proof}
Apply Theorem~\ref{nca:thm:finite-values} to the first 30 and
all 32 vectors in \eqref{jmv:eq:IJM-generators}. The resulting
list lengths are exactly
\[
 30+2\binom{30}{2}+\binom{30}{3}=4960,\qquad
 32+2\binom{32}{2}+\binom{32}{3}=5984.
\]
Write all pair and triple symbols in the next tables using these
\(b_i\), not the \(g_i\) or \(h_i\). Substitution in the full
polynomial gives the complete denominator certificate
\begin{equation}
\begin{array}{c|c|c|c|c}
 &A_i&D^+_{ij}&D^-_{ij}&D_{ijk}\\\hline
 J&(27,216,1)&(9,36,18)&(9,36,18)&(9,36,18)\\
 M&(27,216,1)&(9,36,18)&(9,36,18)&(9,36,18)
\end{array}.
\label{jmv:eq:JM-denominators}
\end{equation}
As before each triple is the least common multiple of all reduced
rational denominators in its full category, in the order
\((P_0,P_2,P_3)\); a zero has denominator one. The finite
arithmetic is completely specified by
\eqref{nca:eq:original-s}--\eqref{nca:eq:P3},
\eqref{jmv:eq:IJM-generators}, and all indices in
\eqref{nca:eq:single}--\eqref{nca:eq:Dthree}.
The accompanying standard-library checker
\texttt{check\_jm\_values.py}, using the included
\texttt{cubic\_engine.py}, executes every entry with exact
fractions and asserts this table. It reconstructs all 72 marked
glue classes, checks all 29 \(\beta(g_i)\), and checks both
kernel conditions on all the displayed \(I\) generators.
These finite computations establish the upper containment in
\eqref{jmv:eq:JM-values}, including all mixed polarization terms
involving \(p,q\), not just their individual cubic values.

Here is a short certificate for the reverse containment, using
only \(J\). Coordinates of each row are taken in the ordered
proposed value-group basis
\((1/27,\sqrt2/216,\sqrt3/18)\):
\begin{equation}
\begin{array}{c|r r r}
 F(b_2)&14&-272&0\\
 F(b_4)&5&-128&0\\
 F(b_{27})&23220&-111672&-90\\
 F(b_{28})&13780&-74672&0\\
 F(b_{30})&378&-1701&-18\\
 D^+_{2,27}&7425&-60624&-3\\
 D^+_{11,21}&36&48&-16
\end{array}.
\label{jmv:eq:JM-lower-certificate}
\end{equation}
Let \(C_1,C_2,C_3\) have, respectively, columns obtained by
transposing rows \((2,3,5)\), \((2,5,6)\), and \((1,4,7)\)
of this table. Direct three-by-three determinants are
\begin{align*}
 d_1&=-39859290,&d_2&=11531403,&d_3&=-43244032.
\end{align*}
The following explicit integer identity certifies their gcd one:
\begin{equation}
 (-6882196252)d_1+(-23788908901)d_2+(-907)d_3=1.
 \label{jmv:eq:JM-bezout}
\end{equation}
Put \((c_1,c_2,c_3)=(-6882196252,-23788908901,-907)\).
For each standard basis column \(e_t\in\mathbb Z^3\) the exact
integer reconstruction is
\begin{equation}
 e_t=\sum_{j=1}^3 C_j\bigl(c_j\operatorname{adj}(C_j)e_t\bigr).
 \label{jmv:eq:JM-axis-reconstruction}
\end{equation}
All columns of every \(C_j\) represent integer combinations of
actual cubic values on \(J\); their integer coefficient vectors
in \eqref{jmv:eq:JM-axis-reconstruction} therefore give each
of the three claimed generators as such a combination. This is
an explicit reverse-containment morphism, not an index inference.
The checker also expands every adjugate and verifies all three
axis reconstructions, saving their integer coefficients in
\texttt{EXACT\_RESULTS.json}. Thus \(\mathcal V_F(J)\) contains
the right-hand side of \eqref{jmv:eq:JM-values}. Since \(J\subset M\),
the same certificate proves its containment in \(\mathcal V_F(M)\).
The upper bounds already established finish both equalities.
\end{proof}

For additional exact checks on the unchanged representatives, the
three individual values are
\[
 F(3v)=14-\frac{63}{8}\sqrt2-\sqrt3,\qquad
 F(p)=-2,\qquad F(q)=\frac{64}{27}.
\]
They agree with the corresponding rows of the complete calculation;
they are not substituted for the mixed-term certificate.

\subsection{Exact inclusions, quotient maps, and the factor of the Albert determinant}

\begin{theorem}[The exact four-to-three chain of value groups]
\label{jmv:thm:chain}
There is a strict chain of additive subgroups of the original real
coefficient field
\[
 \mathcal V_F(N)\ \subsetneq\ \mathcal V_F(J)=\mathcal V_F(M)
                       \ \subsetneq\ \mathcal V_F(\Lambda).
\]
The inclusion matrices in the ordered bases already displayed are
\[
 \mathcal V_F(N)\longrightarrow\mathcal V_F(J)=\mathcal V_F(M):
       (a,b,c)\longmapsto(a,4b,c),
\]
\[
 \mathcal V_F(J)=\mathcal V_F(M)\longrightarrow\mathcal V_F(\Lambda):
       (a,b,c)\longmapsto(a,b,3c).
\]
Their exact quotient maps are
\begin{align*}
 \frac a{27}+\frac{b\sqrt2}{216}+\frac{c\sqrt3}{18}
       &\longmapsto b\pmod4,\\
 \frac a{27}+\frac{b\sqrt2}{216}+\frac{c\sqrt3}{54}
       &\longmapsto c\pmod3.
\end{align*}
In particular
\begin{align*}
 \mathcal V_F(J)/\mathcal V_F(N)&\cong\mathbb Z/4,&
 \mathcal V_F(M)/\mathcal V_F(N)&\cong\mathbb Z/4,\\
 \mathcal V_F(\Lambda)/\mathcal V_F(J)&\cong\mathbb Z/3,&
 \mathcal V_F(\Lambda)/\mathcal V_F(M)&\cong\mathbb Z/3,\\
 \mathcal V_F(M)/\mathcal V_F(J)&=0.&&
\end{align*}
\end{theorem}
\begin{proof}
The first basis comparison uses
\(\sqrt2/54=4(\sqrt2/216)\), with the other basis entries equal.
The second uses \(\sqrt3/18=3(\sqrt3/54)\), with its other
entries equal. Uniqueness of rational radical coefficients proves
that the displayed kernels are exactly the corresponding subgroups.
Surjectivity is witnessed by \(\sqrt2/216\) and \(\sqrt3/54\),
respectively, both actual elements of the proved additive groups.
Their orders modulo the smaller groups are exactly four and three:
the coefficient divisibility conditions are necessary and sufficient.
This also proves strictness. The last quotient is zero by the equality
of the two proved groups. In particular the composite inclusion is
\((a,b,c)\mapsto(a,4b,3c)\), recovering the full quotient
\(\mathbb Z/4\oplus\mathbb Z/3\), with explicit map
\[
 \frac a{27}+\frac{b\sqrt2}{216}+\frac{c\sqrt3}{54}
       \longmapsto (b\bmod4,c\bmod3).
\]
Its identification with \(\mathbb Z/12\) is
\((\bar b,\bar c)\mapsto9b+4c\pmod{12}\), whose inverse is
reduction modulo four and modulo three. Every assertion follows
from the proved cubic coefficient groups, not from an index of
the original lattices.
\end{proof}

The factor two in the original Albert determinant is unchanged:
\[
 \det Q_{\mathbb O}(\lambda;\mathcal LX)
       =\lambda^3-\lambda(X,X)+2F(X).
\]
Thus, if one instead asks for the additive group generated by the
actual cubic term \(2F\), it is exactly \(2\mathcal V_F(L)\).
For \(J,M\) this is
\[
 \frac2{27}\mathbb Z\ \oplus\frac{\sqrt2}{108}\mathbb Z
                           \ \oplus\frac{\sqrt3}{9}\mathbb Z.
\]
Multiplication by two is an injective additive map of the real
coefficient field and restricts to the exact isomorphism
$\mathcal V_F(L)\to\mathcal V_{2F}(L)=2\mathcal V_F(L)$,
$x\mapsto2x$, with inverse $y\mapsto y/2$.
It intertwines the same two inclusion matrices and
the same quotients; no factor has been absorbed into \(F\).

\subsection{Reproduction and scope of the exact certificate}

The accompanying evidence folder is
\begin{flushleft}\small
\texttt{supporting\_materials/workbench/evidence/}\\
\texttt{period\_lattice\_continuation\_20260906/jm\_cubic/}
\end{flushleft}
It contains both full source files \texttt{cubic\_engine.py} and
\texttt{check\_jm\_values.py}. Running the latter with Python's
standard library performs the following finite proof certificate:
all \(20399\) polarization entries for \(N,\Lambda,J,M\), all
displayed denominator and determinant assertions, the integer
Bezout and adjugate axis reconstructions, and the exact inclusion
coefficient comparisons. The output \texttt{EXACT\_RESULTS.json}
retains each selected column, determinant, Bezout coefficient,
and every expanded integer linear combination reconstructing a
value-group axis.

There is additionally an independent coefficient-identity check:
the original 24 simple roots form a real basis of the original
24-dimensional space. The direct Cayley--Dickson implementation
and the rational polynomial implementation are compared on the
24 basis vectors, their 552 ordered-by-index plus/minus pairs,
and their 2024 three-element sums, a total of 2600 exact inputs.
The finite polarization proof shows these evaluations determine
every coefficient of a homogeneous cubic in those 24 variables.
Consequently this is a full polynomial-identity certificate on
the real ambient space, not sampling of lattice points. The
field implementation retains all four rational coefficients in
\(\mathbb Q(\sqrt2,\sqrt3)\), and verifies that the \(\sqrt6\)
coefficient is identically zero. A further 65 direct checks cover
the displayed new generating lists and \(u,p,q\).

This computation proves the stated additive value groups and
their morphisms. It asserts neither equality of individual-value
sets nor that the cubic group reconstructs a lattice or classifies
all cubic-preserving lattice isometries.

% Complete proof module: Wilson construction and retained-coordinate comparison.
\section{A cyclic Wilson Leech lattice in the original triality space}
\label{wcl:sec:construction}

\subsection{The two retained metrics and the octonion convention}

Let $V=\mathbb O_1\oplus\mathbb O_2\oplus\mathbb O_3$ be the original
ordered octonion space. Retain its metric and cubic
\[
\begin{aligned}
 h_0(x,x)&=N(x_1)+N(x_2)+N(x_3),\\
 \tau(x_1,x_2,x_3)&=\operatorname{Re}((x_1x_2)x_3),\qquad
 d(Q_{\mathbb O}(0;x))=2\tau(x_1,x_2,x_3),
\end{aligned}
\]
where
\[
 Q_{\mathbb O}(\lambda;x)=
 \begin{pmatrix}\lambda&x_3&\bar x_2\\
 \bar x_3&\lambda&x_1\\x_2&\bar x_1&\lambda\end{pmatrix}.
\]
Wilson's norm on the same underlying ordered space is
\[
 q_{\rm W}(x)=\tfrac12 h_0(x,x).
\]
The map
\begin{equation}
 \iota:(V,q_{\rm W})\longrightarrow(V,h_0),\qquad
 \iota(x_1,x_2,x_3)=(x_1/\sqrt2,x_2/\sqrt2,x_3/\sqrt2)
 \label{wcl:eq:metric-isometry}
\end{equation}
is an explicit isometry. Both metrics remain specified. In particular,
for the original $\lambda$, which is not rescaled,
\begin{equation}
 d(Q_{\mathbb O}(\lambda;\iota x))
 =\lambda^3-\lambda q_{\rm W}(x)+\frac1{\sqrt2}\tau(x_1,x_2,x_3).
 \label{wcl:eq:cubic-scaling}
\end{equation}
Indeed the quadratic coordinates scale by $1/2$, while the trilinear
term scales by $(1/\sqrt2)^3=1/(2\sqrt2)$ and its original coefficient
is $2$. Thus the off-diagonal cubic after $\iota$ is
$1/(2\sqrt2)$ times its value before $\iota$.

For an explicit realization inside the same $\mathbb O$, choose a
Cayley--Dickson basis
$(1,\mathsf i,\mathsf j,\mathsf k,\ell,\mathsf i\ell,
\mathsf j\ell,\mathsf k\ell)$, with
\[
 (a+b\ell)(c+d\ell)=(ac-\bar d b)+(da+b\bar c)\ell,\qquad
 \mathsf i\mathsf j=\mathsf k .
\]
Wilson's ordered orthonormal basis is realized by
\begin{equation}
 (i_\infty,i_0,i_1,i_2,i_3,i_4,i_5,i_6)
 =(1,\mathsf i,\mathsf j,\ell,\mathsf k,
                  \mathsf j\ell,-\mathsf k\ell,\mathsf i\ell).
 \label{wcl:eq:actual-basis-map}
\end{equation}
The seven oriented Fano triples are $(t,t+1,t+3)$ modulo $7$:
$i_ti_{t+1}=i_{t+3}$, with cyclic products within each triple and
opposite signs on reversal. To verify \eqref{wcl:eq:actual-basis-map},
its seven corresponding products, in order $t=0,\ldots,6$, are
\[
 \mathsf i\mathsf j=\mathsf k,\quad
 \mathsf j\ell=\mathsf j\ell,\quad
 \ell\mathsf k=-\mathsf k\ell,\quad
 \mathsf k(\mathsf j\ell)=\mathsf i\ell,\quad
 (\mathsf j\ell)(-\mathsf k\ell)=\mathsf i,\quad
 (-\mathsf k\ell)(\mathsf i\ell)=\mathsf j,\quad
 (\mathsf i\ell)\mathsf i=\ell.
\]
Alternativity supplies cyclic products on each quaternionic line.
The seven lines cover all pairs of different imaginary basis units;
squares are $-1$ and $1$ is the unit. This proves the algebra
isomorphism on the full multiplication table, hence by bilinearity on
all octonions. It also preserves conjugation and the norm. No
conjugation of a coordinate of $Q_{\mathbb O}$ is performed.

\subsection{The explicit eight-dimensional lattice and two sublattices}

In the displayed Wilson basis write $e_0=i_\infty,e_{t+1}=i_t$ and put
\[
\begin{aligned}
 D_8&=\{(z_0,\ldots,z_7)\in\mathbb Z^8:\textstyle\sum z_i\equiv0\pmod2\},\\
 s&=\tfrac12(-1+i_0+i_1+i_2+i_3+i_4+i_5+i_6),\qquad
 L_{\rm W}=D_8+\mathbb Zs.
\end{aligned}
\]
This is Wilson's lattice $L$, renamed $L_{\rm W}$ to distinguish it
from the project's original rank-24 Niemeier lattice. Since $2s\in
D_8$ and $s\notin D_8$, it is exactly $D_8\cup(s+D_8)$.

\begin{lemma}\label{wcl:lem:e8}
The lattice $L_{\rm W}$ is even and unimodular, has minimum norm $2$,
and has the explicit basis
\[
 b_1=s,\quad b_j=e_{j-1}-e_j\ (2\le j\le7),\quad b_8=e_6+e_7.
\]
Its Gram matrix for the unchanged octonion norm is
\begin{equation}
 G_8=
 \begin{pmatrix}
 2&0&0&0&0&0&0&1\\
 0&2&-1&0&0&0&0&0\\
 0&-1&2&-1&0&0&0&0\\
 0&0&-1&2&-1&0&0&0\\
 0&0&0&-1&2&-1&0&0\\
 0&0&0&0&-1&2&-1&-1\\
 0&0&0&0&0&-1&2&0\\
 1&0&0&0&0&-1&0&2
 \end{pmatrix},\qquad \det G_8=1 .
 \label{wcl:eq:e8-gram}
\end{equation}
Its norm-two vectors are precisely the $112$ vectors $\pm e_i\pm e_j$
with $i<j$ and the $128$ vectors
$\frac12(\pm e_0\pm\cdots\pm e_7)$ with an odd number of minus signs.
\end{lemma}

\begin{proof}
For $z\in D_8$, $N(z)\equiv\sum z_i=0\pmod2$ and
$h(s,z)=\frac12\sum z_i-z_0\in\mathbb Z$.
As $N(s)=2$, every norm in $D_8+\mathbb Zs$ is even.
The covolume of $D_8$ is $2$, and its displayed extension has index
$2$, so $L_{\rm W}$ has covolume $1$. The integral bilinear form follows
by polarizing the even norm, and covolume one proves self-duality.
The columns $b_i$ all belong to $L_{\rm W}$. Their determinant has
absolute value one: expansion along the first coordinate gives
$1/2$ times the determinant $2$ of the indicated $D_7$ basis.
They therefore form a basis. Their scalar products give the displayed
matrix. Positivity and evenness give minimum at least $2$, and
$N(b_i)=2$ attains it. Integer vectors of norm two have exactly two
nonzero entries $\pm1$. Half-integer vectors have norm at least
$8/4=2$, with equality exactly when all eight entries are $\pm1/2$.
Their membership in $s+D_8$ is exactly the stated odd-minus parity.
\end{proof}

Define, with the bars in the positions printed in Wilson,
\[
 \mathsf M=L_{\rm W}\bar s,\qquad \mathsf N=L_{\rm W}s .
\]
All products are right products. Let $B$ be the basis matrix with
columns $b_i$. Right multiplication by $s$ in this basis is the
following integer matrix:
\begin{equation}
 S=B^{-1}R_s B=
 \begin{pmatrix}
 3&0&0&0&0&0&0&2\\
 -2&-1&0&1&-1&1&0&-2\\
 -4&-1&-1&1&0&0&1&-3\\
 -6&-1&-1&0&0&1&0&-4\\
 -8&0&-1&0&-1&1&1&-5\\
 -10&-1&0&0&-1&0&1&-5\\
 -5&0&-1&1&-1&0&0&-2\\
 -7&0&0&0&0&0&0&-4
 \end{pmatrix}.
 \label{wcl:eq:S}
\end{equation}
This is a direct multiplication-table calculation using
\eqref{wcl:eq:actual-basis-map}; it can be checked by multiplying each
column $b_j$ by the displayed $s$. Set $T=-I_8-S$. Since
$s+\bar s=-1$, $s\bar s=2$, and alternativity allows repeated $s$
and $\bar s= -1-s$, one has
\begin{equation}
 T=B^{-1}R_{\bar s}B,\qquad S+T=-I_8,\qquad ST=TS=2I_8,
 \qquad S^{\mathsf T}G_8S=T^{\mathsf T}G_8T=2G_8 .
 \label{wcl:eq:ST-identities}
\end{equation}

\begin{lemma}\label{wcl:lem:complementary}
Both $\mathsf M,\mathsf N$ are sublattices of $L_{\rm W}$ of index
$16$ and minimum norm $4$, and
\[
 2L_{\rm W}\subset\mathsf M,\mathsf N\subset L_{\rm W},\qquad
 \mathsf M+\mathsf N=L_{\rm W},\qquad
 \mathsf M\cap\mathsf N=2L_{\rm W}.
\]
Their images $\overline{\mathsf M},\overline{\mathsf N}$ in
$\mathcal E=L_{\rm W}/2L_{\rm W}$ are complementary
four-dimensional $\mathbb F_2$ subspaces.
\end{lemma}

\begin{proof}
The integer matrices $S,T$ give the inclusions in $L_{\rm W}$;
$ST=2I$ gives the inclusions of $2L_{\rm W}$.
Norm composition shows that each right multiplication scales all
squared lengths by $N(s)=N(\bar s)=2$; thus its determinant has
absolute value $(\sqrt2)^8=16$, proving the indices and minimum norms.
The identity $S+T=-I$ gives $\mathsf M+\mathsf N=L_{\rm W}$.
Modulo $2L_{\rm W}$, each image has size $2^8/16=2^4$.
Their sum is the eight-dimensional space $\mathcal E$, so their
intersection is zero. This proves the asserted intersection upstairs.
\end{proof}

\subsection{The binary glue code and its exact generators}

Wilson's original lattice is exactly
\begin{equation}
\begin{split}
 \Lambda_{\rm W}=\{(x,y,z)\in L_{\rm W}^3:\ &
 x+y,x+z,y+z\in\mathsf M,\quad x+y+z\in\mathsf N\}.
\end{split}
 \label{wcl:eq:Wilson-definition}
\end{equation}
It carries Wilson's norm $q_{\rm W}$, as printed. Define the lattice
inside the original triality metric by
\[
 \Lambda_{\rm cyc}=\iota(\Lambda_{\rm W})\subset(V,h_0).
\]
The simultaneous use of $\mathsf M=L_{\rm W}\bar s$ for pair sums
and $\mathsf N=L_{\rm W}s$ for the total sum is essential.

Write any class in $\mathcal E$ uniquely as $m+n$, with
$m\in\overline{\mathsf M}$ and $n\in\overline{\mathsf N}$.
Then the image of \eqref{wcl:eq:Wilson-definition} in $\mathcal E^3$
is the explicit binary code
\begin{equation}
 \mathcal C=
 \{(m_1+n,m_2+n,m_1+m_2+n):
       m_1,m_2\in\overline{\mathsf M},\
       n\in\overline{\mathsf N}\}.
 \label{wcl:eq:code}
\end{equation}
Indeed the pair conditions say that the three $\mathsf N$ components
are equal; the total-sum condition says that the sum of the three
$\mathsf M$ components is zero. This proves both directions of the
code description. Its dimension is $4+4+4=12$.

For a concrete lift, the first four columns of $S$ and of $T$ are
bases of their images modulo $2$. The determinant of the eight-column
matrix $(T_1,T_2,T_3,T_4,S_1,S_2,S_3,S_4)$ is $-7$, so these columns
are independent modulo $2$; each of the two images already has
dimension four. Put $m_j=b_j\bar s$, $n_j=b_js$ for $1\le j\le4$.
A complete set of $36$ integer generators of $\Lambda_{\rm W}$ is
\begin{equation}
\begin{gathered}
 (2b_i,0,0),\ (0,2b_i,0),\ (0,0,2b_i)\quad(1\le i\le8),\\
 (m_j,m_j,0),\ (m_j,0,m_j),\ (n_j,n_j,n_j)\quad(1\le j\le4).
\end{gathered}
 \label{wcl:eq:generators}
\end{equation}
These vectors are all in \eqref{wcl:eq:Wilson-definition}.
Their reductions span precisely \eqref{wcl:eq:code}; the first
$24$ generate its full kernel $(2L_{\rm W})^3$. This proves that
they generate the entire lattice, not only a sublattice.

\begin{theorem}\label{wcl:thm:Leech}
The explicitly generated lattice $\Lambda_{\rm W}$ with $q_{\rm W}$,
and therefore $\Lambda_{\rm cyc}$ with the unaltered $h_0$, is an
even unimodular rank-24 lattice of minimum norm $4$. In particular
it is rootless. This is the actual Wilson Leech construction, through
the explicit isometry \eqref{wcl:eq:metric-isometry}.
\end{theorem}

\begin{proof}
The code proves
$[L_{\rm W}^3:\Lambda_{\rm W}]=2^{24-12}=2^{12}$.
Since $L_{\rm W}$ has covolume one for its original octonion metric,
$\Lambda_{\rm W}$ has covolume $2^{12}$ for $h_0$.
Changing to the separately recorded metric $q_{\rm W}$ multiplies
volume in dimension $24$ by $(1/\sqrt2)^{24}=2^{-12}$,
so its covolume is exactly one.

For every triple of vectors in a real inner product space,
\[
 N(x)+N(y)+N(z)
 =N(x+y)+N(x+z)+N(y+z)-N(x+y+z).
\]
For a vector of $\Lambda_{\rm W}$ the four terms on the right
are norms in $\mathsf M$ or $\mathsf N$, so each is divisible by
$4$: these are twice norms in the even lattice $L_{\rm W}$.
It follows that $q_{\rm W}(x,y,z)$ is even. Polarization of these
even norms gives an integral bilinear form. An integral lattice of
covolume one equals its dual, proving unimodularity.

If a nonzero vector had $q_{\rm W}=2$, its three coordinate norms
would sum to $4$. Each nonzero coordinate in $L_{\rm W}$ has norm
at least $2$, so one coordinate is zero. The other two then lie in
$\mathsf M$, each of whose nonzero vectors has norm at least $4$.
At most one can be nonzero. That coordinate also lies in $\mathsf N$
by the total-sum condition, so it lies in
$\mathsf M\cap\mathsf N=2L_{\rm W}$ and has norm at least $8$,
a contradiction. Hence all nonzero norms are at least $4$.
The vectors $(2b_i,0,0)$ have $q_{\rm W}=4$, so the minimum is
exactly $4$. The isometry $\iota$ transfers every statement to the
original metric.
\end{proof}

\subsection{Cyclic symmetry, the fixed lattice, and exact three-primary glue}

Retain the canonical cycle $c(x,y,z)=(z,x,y)$. Denote the opposite
cycle used in Wilson coordinates by $c_{\rm W}(x,y,z)=(y,z,x)$.
Direct substitution gives
\[
 c_{\rm W}=c^{-1}=c^2,\qquad c_{\rm W}^2=c,\qquad
 \langle c_{\rm W}\rangle=\langle c\rangle.
\]
Both actions are retained with these directions; no action is renamed
by silently reversing it.
The defining three coordinate memberships, three pair conditions,
and total-sum condition are all unchanged by $c_{\rm W}$; hence both cycles preserve
$\Lambda_{\rm W}$ and $\Lambda_{\rm cyc}$. They preserve their metrics
and the original cubic:
\[
 \operatorname{Re}((xy)z)=\operatorname{Re}((yz)x).
\]
This follows from the octonion adjoint identities
$h(ab,d)=h(b,\bar a d)=h(a,d\bar b)$, which imply real associativity,
and $\operatorname{Re}(ab)=\operatorname{Re}(ba)$.
In the Albert algebra $c_{\rm W}$ is realized by conjugation with the
\emph{real} permutation matrix for $(1\,3\,2)$.
It sends the coordinates $(x_1,x_2,x_3)$ to $(x_2,x_3,x_1)$ with no
conjugations, and permutes the ordered diagonal frame accordingly.
The inverse real permutation matrix, for $(1\,2\,3)$, realizes the
canonical $c$ and sends $(x_1,x_2,x_3)$ to $(x_3,x_1,x_2)$.
Thus the fixed lattices of the two cycles are identical. For the
retained residue isometry $\mathcal Lr=c\mathcal L$, the opposite
action satisfies $\mathcal Lr^{-1}=c_{\rm W}\mathcal L$.

\begin{theorem}\label{wcl:thm:fixed-coinvariant}
Before applying $\iota$, the cyclic fixed lattice and its orthogonal
complement in $\Lambda_{\rm W}$ are exactly
\[
 F=\{(n,n,n):n\in\mathsf N\},\qquad
 C=\{(m_1,m_2,-m_1-m_2):m_1,m_2\in\mathsf M\}.
\]
Their ranks are $8$ and $16$. For the retained Wilson metric, and
therefore after $\iota$ for the original triality metric, their Gram
matrices in the displayed bases are
\begin{equation}
 G_F=3G_8,\qquad
 G_C=\begin{pmatrix}2G_8&G_8\\G_8&2G_8\end{pmatrix}.
 \label{wcl:eq:fixed-coinvariant-grams}
\end{equation}
The fixed-lattice basis is $(b_js,b_js,b_js)$, $1\le j\le8$;
the coinvariant bases are $(b_j\bar s,0,-b_j\bar s)$ and
$(0,b_j\bar s,-b_j\bar s)$. In particular the fixed lattice has
minimum norm $6$ and determinant $3^8$.
The coinvariant lattice has determinant $3^8$, and the map which
negates its second $L_{\rm W}$ parameter identifies its displayed
Gram matrix with the usual $A_2\otimes E_8$ Gram matrix.
\end{theorem}

\begin{proof}
A fixed triple has the form $(x,x,x)$. Its pair condition holds for
every $x\in L_{\rm W}$ because $2L_{\rm W}\subset\mathsf M$.
Its total condition is $3x\in\mathsf N$. As $L_{\rm W}/\mathsf N$
has exponent two, $3x\in\mathsf N$ is equivalent to $x\in\mathsf N$.
This proves the first description. Orthogonality to the whole
diagonal real eight-space says $x+y+z=0$. The pair conditions are
then exactly $x,y,z\in\mathsf M$, and the total condition is automatic.
This proves the description of $C$.

Right multiplication by either $s$ or $\bar s$ multiplies the
octonion metric by $2$. Thus
\[
 q_{\rm W}(as,as,as)=3N(a),
\]
and
\[
 q_{\rm W}(a\bar s,b\bar s,-(a+b)\bar s)
 =N(a)+N(b)+N(a+b)=2N(a)+2N(b)+2h(a,b).
\]
Polarization gives both full Gram matrices. Their determinants are
\[
 \det G_F=3^8\det G_8=3^8,\qquad
 \det G_C=\det\begin{pmatrix}2&1\\1&2\end{pmatrix}^{8}
 (\det G_8)^2=3^8.
\]
Minimum norm $6$ for $F$ follows from the exact
factor $3$ and the minimum $2$ in $L_{\rm W}$.
The sign map $b\mapsto-b$ changes both off-diagonal blocks from
$G_8$ to $-G_8$, proving the stated exact comparison to
$A_2\otimes E_8$.
\end{proof}

There is a complete integral comparison, including the glue:
\begin{equation}
 0\longrightarrow F\oplus C\longrightarrow\Lambda_{\rm W}
 \xrightarrow{\ \Sigma\ }\mathsf N/3\mathsf N
 \longrightarrow0,\qquad
 \Sigma(x,y,z)=x+y+z\pmod{3\mathsf N}.
 \label{wcl:eq:three-primary-exact}
\end{equation}
To prove surjectivity, for every $a\in L_{\rm W}$ the explicit vector
\[
 v(a)=(as,a,-a)
\]
lies in $\Lambda_{\rm W}$ and has sum $as$. Its pair sums are
$a(s+1)=-a\bar s\in\mathsf M$,
$a(s-1)=-a\bar s-2a\in\mathsf M$, and $0$;
its total lies in $\mathsf N$. If a lattice vector has total $3n$,
subtracting $(n,n,n)\in F$ leaves a vector of $C$. This proves the
kernel assertion. Thus
\[
 \Lambda_{\rm W}/(F\oplus C)\cong\mathsf N/3\mathsf N
 \cong L_{\rm W}/3L_{\rm W}\cong(\mathbb Z/3)^8,\qquad
 [\Lambda_{\rm W}:F\oplus C]=3^8.
\]
All maps commute with $c_{\rm W}$ and with its inverse $c$; both
cycles act trivially on the quotient.

The discriminant anti-isometry is also explicit. Use the quadratic
discriminant convention $q_D(u+D)=h_D(u,u)\pmod{2\mathbb Z}$
for an even lattice $D$ with its retained metric. The orthogonal
projections of the lift $v(a)$ are
\begin{align}
 f(a)&=\tfrac13(as,as,as),\nonumber\\
 \gamma(a)&=\left(\tfrac23as,\ a-\tfrac13as,\ -a-\tfrac13as\right).
 \label{wcl:eq:glue-projections}
\end{align}
They belong to $F^*$ and $C^*$, respectively: their pairings with
$F$ or $C$ are the pairings of the integral vector $v(a)$ with the
same lattice, because the other projection is orthogonal.
The classes $f(a)+F$ identify $L_{\rm W}/3L_{\rm W}$ with $F^*/F$.
If $\gamma(a)\in C$, then $f(a)=v(a)-\gamma(a)\in
\Lambda_{\rm W}\cap F_{\mathbb R}=F$, so $a\in3L_{\rm W}$.
Since $C^*/C$ has order $\det G_C=3^8$, the map
\[
 f(a)+F\longmapsto\gamma(a)+C
\]
is an isomorphism of the two discriminant groups. Its quadratic
sign is determined without an unspecified convention:
\[
 q_F(f(a))=\tfrac13N(a),\qquad
 q_C(\gamma(a))=q_{\rm W}(v(a))-\tfrac13N(a)
              =2N(a)-\tfrac13N(a).
\]
As $N(a)$ is even, their sum $2N(a)$ is zero modulo $2\mathbb Z$.
This is the required anti-isometry. Equations
\eqref{wcl:eq:three-primary-exact} and
\eqref{wcl:eq:glue-projections} specify the overlattice itself,
including its actual coset representatives.

\subsection{Finite lattice symmetries and exact cubic tests}

Independent sign changes of the three sectors preserve
$\Lambda_{\rm W}$. Changing the sign of $x$ changes a pair sum by
$-2x\in2L_{\rm W}\subset\mathsf M$ and changes the total by
$-2x\in2L_{\rm W}\subset\mathsf N$. The same argument applies to
each sector. All six sector permutations preserve the defining
conditions too. Under signs $(\epsilon_1,\epsilon_2,\epsilon_3)$
the cubic $\tau$ is multiplied by $\epsilon_1\epsilon_2\epsilon_3$;
hence exactly the even sign changes in this subgroup preserve it.
Cyclic permutations preserve $\tau$ as already proved.

The odd permutation $(x,y,z)\mapsto(y,x,z)$ need not preserve it,
even on the lattice itself. In Wilson's basis put
\[
 a=2(1+i_0),\qquad b=2(1+i_1),\qquad d=2(i_2+i_3).
\]
Every coordinate belongs to $2L_{\rm W}$, so $(a,b,d)\in
\Lambda_{\rm W}$. The real part of
$((1+i_0)(1+i_1))(i_2+i_3)$ is $-1$, because
$i_0i_1=i_3$ and the $i_3^2$ term is the only real contribution.
Reversing the first two factors changes this contribution to $+1$.
Consequently
\begin{equation}
 \tau(a,b,d)=-8,\qquad \tau(b,a,d)=8.
 \label{wcl:eq:transposition-cubic}
\end{equation}
After the specified scalar isometry the original determinant cubic
has values $-4\sqrt2$ and $4\sqrt2$.

The simultaneous Fano-index operations
\[
 A:i_t\longmapsto i_{t+1},\qquad
 B_0:i_t\longmapsto i_{2t},\qquad i_\infty\longmapsto i_\infty
\]
are octonion algebra automorphisms, as the defining oriented Fano
products show. They are orthogonal permutations preserving $D_8$ and
the odd-minus half coset, and they fix the displayed $s$ and $\bar s$.
They consequently preserve $\mathsf M,\mathsf N,\Lambda_{\rm W}$.
Applied simultaneously to all three sectors, they preserve $\tau$.
They satisfy $A^7=B_0^3=1$ and $B_0AB_0^{-1}=A^2$; their permutations
are the distinct maps $t\mapsto 2^j t+k$ with $j=0,1,2$ and
$k\in\mathbb F_7$. Thus they form an explicitly realized
$C_7\rtimes C_3$ of order $21$.

The even sign changes and the cycle $c_{\rm W}$ form
$(\mathbb Z/2)^2\rtimes C_3\cong A_4$. They commute with the
simultaneous Fano operations, and their intersection is trivial:
a simultaneous Fano operation fixes the real unit in every sector
and does not permute sectors. Hence an explicit finite group
\[
 A_4\times(C_7\rtimes C_3)
\]
of order $252$ preserves both the literal lattice and the original
cubic. This is a proved subgroup; no assertion that it is the full
cubic stabilizer is needed.

Wilson also gives the lattice symmetries
\[
 r_t(x,y,z)=(x,yi_t,zi_t)\qquad(t=0,\ldots,6).
\]
For completeness their lattice action can be checked solely from
the explicit finite matrices above. Let
$U_t=B^{-1}R_{i_t}B$. The signed permutation multiplication table
and \eqref{wcl:eq:S} give
\begin{equation}
 U_t\in GL_8(\mathbb Z),\qquad
 T^{-1}(U_t-I_8)\in M_8(\mathbb Z),\qquad
 \tfrac12(U_t-I_8)T\in M_8(\mathbb Z)
 \quad(0\le t\le6).
 \label{wcl:eq:unit-matrix-certificates}
\end{equation}
These are finite exact matrix assertions, not assumptions about
octonionic matrix associativity. The full multiplication rule and
the exact rational algorithm checking every entry for all seven
values are included in the accompanying \texttt{check\_wilson.py}.
Concretely, start with the $64$ products
$1i_t=i_t1=i_t$, $i_t^2=-1$, and the seven oriented Fano triples,
form the eight columns $R_{i_t}(e_j)$, conjugate by the explicitly
displayed $B$, and substitute the integer $T=-I_8-S$.
For each of the three matrices in \eqref{wcl:eq:unit-matrix-certificates}
the reduced denominators of every entry are $1$, for each
$t=0,1,2,3,4,5,6$. Also $\det U_t=1$.
This is a reproducible $7\cdot3\cdot64$-entry certificate.
The first condition is also immediate from the signed permutation:
right multiplication by $i_t$ flips four coordinate signs, preserving
the odd-minus parity of $L_{\rm W}$.

The second matrix assertion says
$L_{\rm W}(i_t-1)\subset\mathsf M$, and the third says
$\mathsf M(i_t-1)\subset2L_{\rm W}$. Thus $x+yi_t$ and $x+zi_t$
differ from the original pair sums by elements of $\mathsf M$.
The pair $(y+z)i_t$ remains in $\mathsf M$, and the transformed total
differs from $x+y+z$ by
$(y+z)(i_t-1)\in2L_{\rm W}\subset\mathsf N$.
All defining conditions are preserved. The inverse follows from
$r_t^2(x,y,z)=(x,-y,-z)$ and the proved sign symmetries.
Norm composition proves that $r_t$ is an isometry.

Nevertheless none of these seven $r_t$ preserves the original cubic.
For $a_t=2(1+i_t)$, the triple $(a_t,a_t,a_t)$ lies in
$(2L_{\rm W})^3\subset\Lambda_{\rm W}$. All products below occur in
the associative real subalgebra generated by $i_t$, and give
\[
 \tau(a_t,a_t,a_t)=8\operatorname{Re}((1+i_t)^3)=-16,
\]
whereas $(1+i_t)i_t=-1+i_t$ and
\[
 \tau(a_t,a_ti_t,a_ti_t)
 =8\operatorname{Re}((1+i_t)(-1+i_t)^2)=16.
\]
Therefore the original determinant cubic on the isometric lattice
$\Lambda_{\rm cyc}$ changes from $-8\sqrt2$ to $8\sqrt2$.
This supplies actual lattice vectors as counterexamples, not only
ambient real vectors outside the lattice.

\subsection{Exact comparison with coordinate lattices and the marked lane}

The construction has the explicit integral sandwich
\begin{equation}
 (\sqrt2L_{\rm W})^3
 \ \subset\ \Lambda_{\rm cyc}\
 \subset\ (L_{\rm W}/\sqrt2)^3,
 \qquad
 [\Lambda_{\rm cyc}:(\sqrt2L_{\rm W})^3]
 =[(L_{\rm W}/\sqrt2)^3:\Lambda_{\rm cyc}]=2^{12}.
 \label{wcl:eq:coordinate-sandwich}
\end{equation}
This is \eqref{wcl:eq:code} transported by the exact isometry
\eqref{wcl:eq:metric-isometry}. The inclusions are literal inclusions
in the fixed three octonion sectors, and each commutes with both
$c_{\rm W}$ and the canonical $c=c_{\rm W}^{-1}$.
Together with the explicit bases, the code, and the three-primary
glue \eqref{wcl:eq:glue-projections}, it determines actual
commensurability maps and cosets rather than only an abstract
isometry class.

Under \eqref{wcl:eq:actual-basis-map}, the single negative column
takes the Wilson odd-minus half roots to even-minus half roots in
the current ordered Cayley--Dickson basis. The element $s$ remains
exactly
\[
 s=\tfrac12(-1+\mathsf i+\mathsf j+\mathsf k+\ell
                     +\mathsf i\ell+\mathsf j\ell-\mathsf k\ell),
\]
with coordinate vector $(-1,1,1,1,1,1,1,-1)/2$.
The signed permutation has determinant $+1$ and all $64$ basis
products agree, as independently checked in the accompanying code.
Thus the lattice, its cyclic action, and the triality cubic live
in one specified multiplication convention.

\subsection{Literal intersections with both retained lattice embeddings}
\label{wcl:sec:literal-intersections}

Let $U=\operatorname{span}_{\mathbb R}(1,\mathsf i,\mathsf j,\mathsf k)$
be the first four coordinates in the \emph{current} ordered
Cayley--Dickson basis, and let
\[
 D_4=\{a=(a_0,a_1,a_2,a_3)\in\mathbb Z^4:
                          a_0+a_1+a_2+a_3\in2\mathbb Z\}\subset U.
\]
This quaternionic $D_4$ denotes the displayed four-coordinate lattice;
the original Niemeier $D$ summand below remains in its original place.

Let $J_{\rm W}$ denote the signed permutation matrix in
\eqref{wcl:eq:actual-basis-map}. If $U_0:\mathbb R^4\to\mathbb R^8$
is the coordinate inclusion in the current basis, then
\begin{equation}
 S^{-1}B^{-1}J_{\rm W}^{-1}U_0=
 \begin{pmatrix}
 1/2&-1/2&-1/2&-1/2\\
 -1/2&0&0&1/2\\
 -1&1/2&0&1/2\\
 -3/2&1&1/2&1/2\\
 -2&1&1&1/2\\
 -5/2&3/2&1&1\\
 -5/4&3/4&3/4&3/4\\
 -7/4&3/4&3/4&3/4
 \end{pmatrix}.
 \label{wcl:eq:quaternion-membership-matrix}
\end{equation}
All matrices on the left have already been explicitly specified.
Thus the equation is a direct rational multiplication certificate.

\begin{lemma}\label{wcl:lem:quaternion-intersection}
In this fixed coordinate frame,
\[
 \mathsf N\cap U=2D_4 .
\]
\end{lemma}
\begin{proof}
Since $\mathsf N\subset L_{\rm W}$, an element of $\mathsf N\cap U$
has integral first four coordinates: the half-integer coset of the
current $E_8$ description has no zero coordinate among its last four.
Write those integral coordinates as $k_0,k_1,k_2,k_3$. Membership in
$\mathsf N$ is exactly integrality of the matrix
\eqref{wcl:eq:quaternion-membership-matrix} applied to $k$.
Subtracting its last two rows gives $k_0/2$, so $k_0$ is even.
The second row then forces $k_3$ even, the third forces $k_1$ even,
and the fourth forces $k_2$ even. Write $k=2a$.
The first six rows are now integral. The seventh and eighth rows
are integral precisely when
$-5a_0+3a_1+3a_2+3a_3$ is even, equivalently
$a_0+a_1+a_2+a_3$ is even. Conversely these conditions make all eight
rows integral. This proves both inclusions with no unstated lattice
identification.
\end{proof}

We now retain the original lattices and their actual map $\mathcal L$.
The original coordinates are $X=(x_{i,j};d_k)$,
$1\le i\le4$, $0\le j\le5$, $1\le k\le4$, with
$\sum_jx_{i,j}=0$. The root lattice is
$R=A_5^{\perp4}\perp D$, with its unchanged Euclidean metric.
Its index-$72$ overlattice $N$ is the inverse image of the original
glue
\[
 (k_1,k_2,k_3,k_4;y)
   =\bigl(3u+2(a+b,a+2b,a,b);\,(u_2+u_4,u_3+u_4)\bigr),
\]
where $u\in\mathbb F_2^4$ has even weight and $a,b\in\mathbb F_3$;
the first four entries are modulo $6$. Here $k_i=-6x_{i,0}\pmod6$
and $y$ is the original $D^\vee/D$ marking, so $d=0$ has $y=00$.
Retain the original neighbour
\[
 \rho_A=(5,3,1,-1,-3,-5)/2,\quad
 \rho_D=(3,2,1,0),\quad \rho=\rho_A^{\oplus4}\perp\rho_D,
\]
\[
 K=\{X\in N:(\rho,X)\in6\mathbb Z\},\qquad
 \alpha=(e_0-e_1,0,0,0;0),\qquad
 \Lambda=K+\mathbb Z(\rho/6-\alpha).
\]
Thus $\rho\in N$, $\rho^2=84$, $(\rho,\alpha)=1$.
Every vector of $N$ or $\Lambda$ has rational original coordinates.

For exact use of the retained isometry, set
\[
 \sigma_i=\sum_{j=0}^5(-1)^jx_{i,j},\qquad
 t_{i,p,j}=x_{i,p+2j}-\frac13\sum_{h=0}^2x_{i,p+2h}.
\]
Write $u_0=1,u_1=\mathsf i,u_2=\mathsf j,u_3=\mathsf k$ and retain
the remaining current Cayley--Dickson basis vectors $u_4,\ldots,u_7$.
The already specified map has the exact coordinate form
\[
 (\mathcal LX)_{j+1,a}=
 \begin{cases}
 t_{i,p,j}+\sqrt2\,\sigma_{a+1}/6,&0\le a\le3,\\
 t_{i,p,j}+\sqrt3\,d_{a-3}/3,&4\le a\le7,
 \end{cases}
 \qquad a=2(i-1)+p.
\]
The placement of $\sigma_{a+1}$ is retained, including its difference
from the index $i$ of the residual. No coordinate is reassigned.

\begin{theorem}\label{wcl:thm:literal-intersection}
The literal intersections in the original metric space are exactly
\begin{equation}
 \Lambda_{\rm cyc}\cap\mathcal L(N)
 =\Lambda_{\rm cyc}\cap\mathcal L(\Lambda)
 =\mathcal I
 :=\{\sqrt2(a,a,a):a\in D_4\subset U\}.
 \label{wcl:eq:literal-common-lattice}
\end{equation}
In original Niemeier coordinates the inverse image of this common
lattice is
\[
 \mathcal L^{-1}(\mathcal I)
 =\{x_{i,j}=(-1)^j a_{i-1},\ d=0:\ a\in D_4\}\subset K.
\]
It has rank four, minimum norm $12$, and Gram matrix
\begin{equation}
 G_{\mathcal I}=6
 \begin{pmatrix}
 2&-1&0&0\\
 -1&2&-1&-1\\
 0&-1&2&0\\
 0&-1&0&2
 \end{pmatrix},
 \qquad\det G_{\mathcal I}=4\cdot6^4=5184.
 \label{wcl:eq:common-gram}
\end{equation}
An explicit basis is $\sqrt2(\delta_j,\delta_j,\delta_j)$,
where
$\delta_1=(1,-1,0,0)$, $\delta_2=(0,1,-1,0)$,
$\delta_3=(0,0,1,-1)$, $\delta_4=(0,0,1,1)$.
\end{theorem}

\begin{proof}
Each coordinate of $\Lambda_{\rm cyc}$ belongs to
$\mathbb Q/\sqrt2=\sqrt2\mathbb Q$ in the current frame.
For a rational original vector $X$, the first four coordinate
formulas above therefore force $t_{i,p,j}=0$ at the corresponding
positions, by linear independence of $1,\sqrt2$ over $\mathbb Q$.
The last four force both $t_{i,p,j}=0$ and $d_{a-3}=0$, by linear
independence of $1,\sqrt2,\sqrt3$. For the latter independence, a
relation with nonzero $\sqrt3$ coefficient would place $\sqrt3$ in
$\mathbb Q(\sqrt2)$; squaring $A+B\sqrt2$ forces $AB=0$ and then
would make $3$ or $3/2$ a rational square, a contradiction.
All residuals therefore vanish and $d=0$. The zero fibre sums give
\[
 x_{i,j}=(-1)^j\sigma_i/6.
\]

First let $X\in N$. Its $A_5^\vee$ integral-difference conditions
are exactly $k_i:=\sigma_i/3\in\mathbb Z$, yielding
$x_{i,j}=(-1)^jk_i/2$. Their discriminant labels are
$-3k_i=3k_i\pmod6$. Their three-primary parts are zero, so the
injectivity of $(a,b)\mapsto(a+b,a+2b,a,b)$ forces $a=b=0$.
Since $d=0$, the binary condition is
$u_2+u_4=u_3+u_4=0$ and $\sum u_i=0$. These equations give
$u_1=u_2=u_3=u_4$. Thus all four $k_i$ have the same parity.
Conversely these conditions satisfy every original glue requirement.
The image is the diagonal triple
\[
 \mathcal LX=(k/\sqrt2,k/\sqrt2,k/\sqrt2),\qquad k\in U.
\]
By the fixed-lattice description of Wilson's construction it
belongs to $\Lambda_{\rm cyc}$ exactly when $k\in\mathsf N$.
The previous lemma makes this equivalent to $k=2a$, $a\in D_4$,
which proves the first intersection equality.

For the second equality, a vector of $\Lambda$ can be written
$X=n+m\rho/6$ with $n\in N$, $m\in\mathbb Z$.
The already proved necessary condition $d_X=0$ gives
$d_n=-m(3,2,1,0)/6\in D^\vee$.
All coordinates of a vector of $D^\vee$ are integral or all
half-integral; the fourth coordinate is zero, so all are integral.
Its third coordinate forces $6\mid m$, and hence $X\in N$.
In fact $\Lambda\cap N=K$: in
$k+m(\rho/6-\alpha)\in N$ the pairing with the original root
$\alpha$ forces $6\mid m$, and
$(\rho,\rho-6\alpha)=84-6=78\in6\mathbb Z$.
Thus an intersection vector belongs to $K$.
For the diagonal-fibre vectors just found,
\[
 (\rho,X)=\frac32\sum_{i=1}^4 k_i.
\]
Consequently the condition $X\in K$ is $\sum k_i\equiv0\pmod4$.
Every $k=2a$ with $a\in D_4$ satisfies it, proving equality of
the two intersections and the displayed inverse image.

Finally
$h_0(\sqrt2(a,a,a),\sqrt2(b,b,b))=6h(a,b)$.
The four $\delta_j$ form a basis of $D_4$, since their determinant
has absolute value two and $D_4$ has index two in $\mathbb Z^4$.
Their scalar products give \eqref{wcl:eq:common-gram}.
The least nonzero norm in $D_4$ is two: a norm-one integer vector
has odd coordinate sum, while $(1,-1,0,0)$ has norm two.
This proves the exact minimum and determinant.
\end{proof}

The two slice indices are also literal. Define the common
four-dimensional real subspace
$\mathcal U=\{(u,u,u):u\in U\}$. The proof gives
\[
 \mathcal L(N)\cap\mathcal U
 =\{(k,k,k)/\sqrt2:k\in P\},\quad
 P=2\mathbb Z^4+\mathbb Z(1,1,1,1),
\]
\[
 \mathcal L(\Lambda)\cap\mathcal U
 =\{(k,k,k)/\sqrt2:k\in P_0\},\quad
 P_0=\{k\in P:\textstyle\sum k_i\equiv0\pmod4\}
     =2D_4+\mathbb Z(1,1,1,1).
\]
The first proof does not need Wilson membership to obtain these
slice equations. The lattice $P$ has basis
$2e_0,2e_1,2e_2,(1,1,1,1)$ and covolume $8$.
The lattice $2D_4$ has covolume $2^4\cdot2=32$, so
$[P:2D_4]=4$. The class of $(1,1,1,1)$ has exact order two
modulo $2D_4$, yielding $[P_0:2D_4]=2$. Therefore
\[
 [\mathcal L(N)\cap\mathcal U:\mathcal I]=4,\qquad
 [\mathcal L(\Lambda)\cap\mathcal U:\mathcal I]=2.
\]
The rank-four intersection, strictly below rank $24$, also proves
that these full lattices are not commensurable \emph{in the retained
embedding}: a common finite-index subgroup would have full rank.
Their exact existing relation is the common lattice with the
explicit embeddings, inverse map, metrics, and indices just proved.

\subsection{The cubic value group on the literal common lattice}
\label{wcl:sec:common-cubic}

Retain $F(y)=\tau(y_1,y_2,y_3)$, so the Albert off-diagonal
determinant is $2F$. For $y=\sqrt2(a,a,a)\in\mathcal I$ write
$a=a_0+\mathbf a$ in the same quaternionic $U$. Since
$\mathbf a^2=-N(\mathbf a)$ and $a_0$ is real, direct multiplication
inside this associative subalgebra gives
\[
 \operatorname{Re}(a^3)
 =a_0^3-3a_0N(\mathbf a)=4a_0^3-3a_0N(a).
\]
Every scalar and factor is therefore
\begin{equation}
 F(\sqrt2(a,a,a))
 =2\sqrt2\bigl(4a_0^3-3a_0N(a)\bigr).
 \label{wcl:eq:common-cubic}
\end{equation}
For $a\in D_4$, $N(a)$ is even, because
$N(a)\equiv\sum a_i=0\pmod2$. The integer in parentheses is
therefore even, so every value belongs to $4\sqrt2\mathbb Z$.
For the actual lattice vector $a=(1,1,0,0)$ it equals
$4-3\cdot2=-2$, giving $F=-4\sqrt2$.
Consequently the \emph{additive subgroup generated by all values}
is exactly
\[
 \bigl\langle F(\mathcal I)\bigr\rangle_{\mathbb Z}
       =4\sqrt2\mathbb Z,\qquad
 \bigl\langle d(Q_{\mathbb O}(0;\mathcal I))\bigr\rangle_{\mathbb Z}
       =8\sqrt2\mathbb Z.
\]
This assertion concerns the generated additive groups; it does not
claim that every element of either group is an individually attained
cubic value.
For an explicit distinction, put
$P(a)=a_0(a_0^2-3(a_1^2+a_2^2+a_3^2))$.
If $3\mid P(a)$, reduction modulo three forces $3\mid a_0$,
and substitution $a_0=3b$ gives
$P(a)=9b(3b^2-a_1^2-a_2^2-a_3^2)$.
Thus $3\mid P(a)$ implies $9\mid P(a)$.
The value $P(a)=6$ is impossible, so $F=12\sqrt2$ is not
attained, although it belongs to the proved additive value group.


\subsection{Primary source and exact reproducibility data}

Wilson's octonionic construction is the source of the pair and total
conditions \cite[Sections~2--4, pp.~2--6]{esn:Wilson2009}. His page~2
uses the odd-minus lattice $L$ and $R=\bar L$; page~3 gives pair sums
in $L\bar s$, the total in $Ls$, and the metric equal to one half of
the sum of octonion norms. The displayed basis map and
\eqref{wcl:eq:metric-isometry} retain these conjugations and the exact
scale in the present coordinates. The distinct coordinate change on
Wilson's page~6 is not substituted for this isometry.

The complete integer-arithmetic program and its output are supplied in the
following directory (the two displayed lines form one path):
\begin{quote}\small
\path{supporting_materials/workbench/evidence/}\\
\path{lattice_foundation_20260906/wilson/}
\end{quote}
The files are \path{check_wilson.py} and \path{EXACT_CHECKS.json}.
They record the generating-lattice index $4096$, Gram determinant $1$,
even diagonal, integral cyclic action, the fixed and coinvariant Gram
matrices, and the index $6561$ of their sum. The same program checks
all $64$ multiplication-table products for
\eqref{wcl:eq:actual-basis-map} and the stated exact cubic tests.
The proofs above establish rootlessness and the global descriptions;
these conclusions are not inferred from a finite search.


% Complete TeX conversion of THREE_LINE_IMAGE_PROOF.md, Sections 1--5.
% Retains the original F=tau from es_niemeier_wilson_cyclic.tex.
\section{The exact cubic image of the original common lattice}
\label{d4i:sec:complete-image}

Retain the coordinates $a_0,a_1,a_2,a_3$ and the original lattice
\[
 D_4=\{a\in\mathbb Z^4:a_0+a_1+a_2+a_3\in2\mathbb Z\}.
 \label{d4i:eq:original-domain}\tag{D4I1}
\]
The polynomial already used in \cref{wcl:sec:common-cubic} is
\[
 P(a)=a_0\bigl(a_0^2-3(a_1^2+a_2^2+a_3^2)\bigr).
 \label{d4i:eq:original-polynomial}\tag{D4I2}
\]
Here and throughout this section $F$ is the original triality cubic
$F=\tau$, with its ordered octonionic arguments as in
\cref{wcl:sec:construction,wcl:sec:common-cubic}. In particular, on the
literal common lattice of \cref{wcl:thm:literal-intersection},
\[
 \mathcal I=\{\sqrt2(a,a,a):a\in D_4\},\qquad
 F(\sqrt2(a,a,a))=2\sqrt2\,P(a).
 \label{d4i:eq:retained-cubic}\tag{D4I3}
\]
No coordinate permutation, scaling of $P$, or alteration of the lattice
is made. The witnesses below are actual points of this original lattice.
The image proof is elementary and requires no external arithmetic
existence theorem.

\subsection{Necessary restrictions in the original coordinates}
\label{d4i:sec:necessary}

For every integer $b$, one has $b^2\equiv b\pmod2$. Consequently,
for $a\in D_4$,
\[
 a_0^2+a_1^2+a_2^2+a_3^2
 \equiv a_0+a_1+a_2+a_3\equiv0\pmod2.
\]
Writing $s=a_1^2+a_2^2+a_3^2$, the factor $a_0^2-3s$ is congruent
to $a_0^2+s\equiv0\pmod2$. Thus $P(a)$ is even.

Modulo three the original polynomial satisfies $P(a)\equiv a_0^3$.
If $3\mid P(a)$, the field $\mathbb Z/3\mathbb Z$ has no nonzero
element with zero cube, so $3\mid a_0$. Substitution of $a_0=3b$
into the original formula gives
\[
 P(a)=3b(9b^2-3s)=9b(3b^2-s).
 \label{d4i:eq:three-divisibility}\tag{D4I4}
\]
Therefore
\[
 P(a)\in2\mathbb Z,\qquad
                 3\mid P(a)\ \Longrightarrow\ 9\mid P(a).
 \label{d4i:eq:necessary-restrictions}\tag{D4I5}
\]

\subsection{Three explicit affine lines and their exact values}
\label{d4i:sec:three-lines}

Define the common direction
\[
 v=(-3,1,1,1).
 \label{d4i:eq:common-direction}\tag{D4I6}
\]
Its coordinate sum is zero, and direct substitution gives
\[
                   P(v)=-3\bigl(9-3(1+1+1)\bigr)=0.
\]
For every $r\in\mathbb Z$ and $\varepsilon\in\{-1,1\}$, put
\[
 L_\varepsilon(r)=(-3r-\varepsilon,r,r,r+\varepsilon)
                =rv+(-\varepsilon,0,0,\varepsilon).
 \label{d4i:eq:two-lines}\tag{D4I7}
\]
Its coordinate sum is
$-3r-\varepsilon+r+r+r+\varepsilon=0$, so it belongs to $D_4$.
The original last-three-coordinate square sum and first-coordinate
square are respectively
\[
 s=r^2+r^2+(r+\varepsilon)^2=3r^2+2r\varepsilon+1,
 \qquad (-3r-\varepsilon)^2=9r^2+6r\varepsilon+1.
\]
Thus the factor in the original cubic is exactly
\[
 a_0^2-3s=(9r^2+6r\varepsilon+1)
                 -(9r^2+6r\varepsilon+3)=-2,
\]
and therefore
\[
 P(L_\varepsilon(r))=(-3r-\varepsilon)(-2)=6r+2\varepsilon.
 \label{d4i:eq:two-line-values}\tag{D4I8}
\]

For every $h\in\mathbb Z$, define the third line
\[
 L_0(h)=(-3h,h+1,h-1,h)=hv+(0,1,-1,0).
 \label{d4i:eq:third-line}\tag{D4I9}
\]
Its coordinate sum is zero, so it also belongs to $D_4$. Here
\[
 s=(h+1)^2+(h-1)^2+h^2=3h^2+2,
\]
and consequently
\[
 P(L_0(h))=(-3h)\bigl(9h^2-3(3h^2+2)\bigr)
                         =(-3h)(-6)=18h.
 \label{d4i:eq:third-line-values}\tag{D4I10}
\]
These identities hold for all signed integer parameters. In particular,
$L_0(0)=(0,1,-1,0)$ has cubic value zero, and the zero vector also
has cubic value zero.

\subsection{The global image and a proved reconstruction map}
\label{d4i:sec:global-image}

Set
\[
 S=\{n\in\mathbb Z:n\in2\mathbb Z\text{ and }
                                  (3\nmid n\text{ or }9\mid n)\}.
 \label{d4i:eq:image-domain}\tag{D4I11}
\]

\begin{theorem}[Complete cubic image and explicit section]
\label{d4i:thm:complete-image}
The original polynomial has the exact image
\[
 P(D_4)=18\mathbb Z\ \cup\ (6\mathbb Z+2)\ \cup\ (6\mathbb Z-2)=S.
 \label{d4i:eq:complete-image}\tag{D4I12}
\]
Equivalently,
\[
 P(D_4)=\{n\in\mathbb Z:
                 n\bmod18\in\{0,2,4,8,10,14,16\}\}.
 \label{d4i:eq:residue-image}\tag{D4I13}
\]
The following map is an actual section $\sigma:S\to D_4$:
\[
 \sigma(n)=
 \begin{cases}
 (-3h,h+1,h-1,h),&n=18h,\\
 (-3r-\varepsilon,r,r,r+\varepsilon),
        &n=2(3r+\varepsilon),\quad\varepsilon\in\{-1,1\}.
 \end{cases}
 \label{d4i:eq:global-section}\tag{D4I14}
\]
Its displayed parameters are uniquely determined in their respective
cases, and $P\circ\sigma=\operatorname{id}_S$.
\end{theorem}

\begin{proof}
Equation~\eqref{d4i:eq:necessary-restrictions} proves
$P(D_4)\subseteq S$. For the reverse inclusion let $n\in S$.
If $3\mid n$, then $9\mid n$, and evenness gives $n=18h$ for
a unique $h\in\mathbb Z$. Equation~\eqref{d4i:eq:third-line-values}
gives $P(L_0(h))=n$.

If $3\nmid n$, write $n=2q$. Then $3\nmid q$, so there is a
unique $\varepsilon\in\{-1,1\}$ with
$q\equiv\varepsilon\pmod3$. The number
$r=(q-\varepsilon)/3$ is consequently a unique integer.
Equation~\eqref{d4i:eq:two-line-values} gives
\[
                    P(L_\varepsilon(r))=6r+2\varepsilon=2q=n.
\]
The two displayed cases are disjoint and cover $S$, and both
witnesses lie in $D_4$ by \cref{d4i:sec:three-lines}. These
calculations prove the asserted section identity and both inclusions
in \eqref{d4i:eq:complete-image}. The first arithmetic progression
contributes residue zero modulo eighteen; the second contributes
$2,8,14$; and the third contributes $16,4,10$. This proves
\eqref{d4i:eq:residue-image}.
\end{proof}

The classification includes every positive, zero, and negative value.
The full excluded set is the union of the odd integers with
$18\mathbb Z+6$ and $18\mathbb Z+12$.
Let
\[
 A=\{a\in\mathbb Z^4:a_0+a_1+a_2+a_3=0\}
 \label{d4i:eq:sum-zero-slice}\tag{D4I15}
\]
be the literal sum-zero sublattice in the original coordinates. The
three witnesses lie in $A\subseteq D_4$, so the two inclusions
$S\subseteq P(A)\subseteq P(D_4)=S$ prove
\[
                              P(A)=P(D_4)=S.
 \label{d4i:eq:slice-image}\tag{D4I16}
\]
The image of the original cubic $F=\tau$ on the common lattice is
therefore exactly
\[
\begin{aligned}
 F(\mathcal I)&=2\sqrt2\,S\\
 &=4\sqrt2\bigl((9\mathbb Z)\cup(3\mathbb Z+1)
                                     \cup(3\mathbb Z-1)\bigr).
\end{aligned}
 \label{d4i:eq:original-cubic-image}\tag{D4I17}
\]
Because the original off-diagonal Albert determinant is $2F$, its
corresponding image is $4\sqrt2\,S$. The additive group generated
by that determinant image is $8\sqrt2\mathbb Z$, while the
additive group generated by $F(\mathcal I)$ is $4\sqrt2\mathbb Z$.
Indeed every $P$ value is even, and the first two lines at parameter
zero attain $P=2$ and $P=-2$. Thus the generated additive group of
the $P$ values is exactly $2\mathbb Z$; applying the two original
scalar factors proves both group assertions without identifying an
image set with its generated group.

\subsection{Exact local images with explicit local sections}
\label{d4i:sec:local-images}

For a prime $p$, let
\[
 D_{4,p}=\{a\in\mathbb Z_p^4:
                            a_0+a_1+a_2+a_3\in2\mathbb Z_p\}.
 \label{d4i:eq:local-domain}\tag{D4I18}
\]
This is precisely $D_4\otimes_{\mathbb Z}\mathbb Z_p$ in the original
coordinates. To verify the equality directly, use the integer basis
\[
\begin{gathered}
 b_0=(1,-1,0,0),\quad b_1=(0,1,-1,0),\\
 b_2=(0,0,1,-1),\quad b_3=(0,0,1,1).
\end{gathered}
 \label{d4i:eq:integral-basis}\tag{D4I19}
\]
For any $a$ in the displayed definition of $D_{4,p}$, the unique
coefficients of these basis vectors are
\[
\begin{aligned}
 c_0&=a_0,&c_1&=a_0+a_1,\\
 c_3&=(a_0+a_1+a_2+a_3)/2,&c_2&=c_3-a_3.
\end{aligned}
 \label{d4i:eq:local-basis-inverse}\tag{D4I20}
\]
All belong to $\mathbb Z_p$, and substitution gives the original
four coordinates. Conversely a $\mathbb Z_p$-linear combination
of the four vectors has coordinate sum $2c_3$. These two direct
calculations prove the tensor-product identification, including
at $p=2$. Taking integral coefficients gives the same basis
description of the original $D_4$.

The three line identities of \cref{d4i:sec:three-lines} hold over
$\mathbb Z_p$, since they are integer polynomial identities. Each
line stays in the sum-zero submodule of $D_{4,p}$.

\begin{theorem}[All local images and the exact local-to-global statement]
\label{d4i:thm:local-images}
For the unchanged polynomial $P$,
\[
\begin{aligned}
 P(D_{4,2})&=2\mathbb Z_2,\\
 P(D_{4,3})&=\mathbb Z_3^{\times}\cup9\mathbb Z_3,\\
 P(D_{4,p})&=\mathbb Z_p\qquad(p\ge5).
\end{aligned}
 \label{d4i:eq:local-images}\tag{D4I21}
\]
An integer $n$ belongs to every local image if and only if
$n\in S$, and hence if and only if $n\in P(D_4)$. The global
section is exactly \eqref{d4i:eq:global-section}; explicit local
sections are given in the proof.
\end{theorem}

\begin{proof}
For $p=2$, the parity calculation of \cref{d4i:sec:necessary}, reduced
modulo two, gives $P(D_{4,2})\subseteq2\mathbb Z_2$. Conversely,
given an arbitrary $n=2q\in2\mathbb Z_2$, put
\[
 r=(q-1)/3\in\mathbb Z_2.
\]
Division by three is valid since three is a unit in $\mathbb Z_2$.
Then $L_1(r)\in D_{4,2}$ and
\[
                         P(L_1(r))=6r+2=2q=n.
\]
Thus $n\mapsto L_1((n/2-1)/3)$ is an explicit section on
$2\mathbb Z_2$, and proves the first equality.

For $p=3$, reduction of the original polynomial modulo three and
the substitution $a_0=3b$ in
\eqref{d4i:eq:three-divisibility} prove that any value divisible by
three belongs to $9\mathbb Z_3$. Every unit
$n\in\mathbb Z_3^{\times}$ is congruent modulo three to exactly
one of $2\varepsilon$, with $\varepsilon\in\{-1,1\}$. Set
\[
 r=(n-2\varepsilon)/6\in\mathbb Z_3.
\]
This quotient is integral because the numerator is divisible by
three and two is a unit. The line $L_\varepsilon(r)$ then has
$P$ value $6r+2\varepsilon=n$.
Every $n\in9\mathbb Z_3$ can be written uniquely as $18h$ with
$h=n/18\in\mathbb Z_3$, again because two is a unit.
The third line has $P(L_0(h))=18h=n$. These two disjoint cases
give the explicit local section
\[
 n\longmapsto
 \begin{cases}
 L_\varepsilon((n-2\varepsilon)/6),
        &n\in\mathbb Z_3^{\times},\quad n\equiv2\varepsilon\pmod3,\\
 L_0(n/18),&n\in9\mathbb Z_3.
 \end{cases}
 \label{d4i:eq:three-adic-section}\tag{D4I22}
\]
They prove the second equality in \eqref{d4i:eq:local-images}.

For every prime $p\ge5$, six is a unit. Given any
$n\in\mathbb Z_p$, choose $r=(n-2)/6\in\mathbb Z_p$.
The point $L_1(r)$ belongs to $D_{4,p}$ and has
$P(L_1(r))=6r+2=n$. Thus $n\mapsto L_1((n-2)/6)$ is a
section on all of $\mathbb Z_p$, proving the third equality.

For integers $n$, the first local condition is precisely evenness;
the second is precisely $3\nmid n$ or $9\mid n$; and the other
primes impose no condition. Membership in every local image is
therefore exactly $n\in S$. The explicit integral section in
\cref{d4i:thm:complete-image} proves that those conditions suffice
for global integral representability. This proves the local-to-global
assertion without approximation or inference from a finite modular
calculation.
\end{proof}

\subsection{The exact factorization and quotient metric on the sum-zero slice}
\label{d4i:sec:factorization}

Define the integral map into the original sublattice $A$ of
\eqref{d4i:eq:sum-zero-slice} by
\[
 \Phi:\mathbb Z^3\longrightarrow A,\qquad
 \Phi(w,u,v)=(-3w-u-v,w+u,w+v,w).
 \label{d4i:eq:slice-map}\tag{D4I23}
\]
Its coordinate sum is zero. For $a\in A$, its proposed inverse is
\[
                 \Phi^{-1}(a)=(a_3,a_1-a_3,a_2-a_3).
 \label{d4i:eq:slice-inverse}\tag{D4I24}
\]
Substitution returns the original second, third, and fourth
coordinates immediately; the first becomes $-a_1-a_2-a_3=a_0$
by the defining equation of $A$. Substituting $\Phi$ into the
inverse returns $(w,u,v)$. These two equalities prove that $\Phi$
is an exact integral-module isomorphism.

For $a=(a_0,y,z,w)\in A$, put $T=y+z+w=-a_0$. Direct expansion
of the original polynomial gives
\[
\begin{aligned}
 P(a)&=-T\bigl(T^2-3(y^2+z^2+w^2)\bigr)\\
     &=2T(y^2+z^2+w^2-yz-yw-zw)\\
     &=T\bigl((y-z)^2+(y-w)^2+(z-w)^2\bigr).
\end{aligned}
 \label{d4i:eq:slice-expansion}\tag{D4I25}
\]
Under $\Phi$, one has $T=3w+u+v$, and the three differences are
$u-v,u,v$. Consequently
\[
                 P(\Phi(w,u,v))=2(3w+u+v)(u^2-uv+v^2).
 \label{d4i:eq:factorization}\tag{D4I26}
\]
This identity explains the three original lines without changing
their coordinates. On $L_\varepsilon(r)$, the quotient differences
are $(u,v)=(-\varepsilon,-\varepsilon)$, so
$u^2-uv+v^2=1$; on $L_0(h)$, they are $(u,v)=(1,-1)$, so
$u^2-uv+v^2=3$. The original first coordinates remain
$-3r-\varepsilon$ and $-3h$, respectively.

The exact quotient map is
\[
 \chi:A\longrightarrow\mathbb Z^2,\qquad
                  \chi(a)=(a_1-a_3,a_2-a_3).
 \label{d4i:eq:quotient-map}\tag{D4I27}
\]
It is onto because $\chi\Phi(0,u,v)=(u,v)$. An element lies in
its kernel exactly when its last three coordinates all equal a
common integer $w$; its first coordinate is then $-3w$.
Thus $\ker\chi=\mathbb Z(-3,1,1,1)$. The section
$(u,v)\mapsto\Phi(0,u,v)$ proves the split exact sequence
\[
 0\longrightarrow\mathbb Z(-3,1,1,1)\longrightarrow A
       \xrightarrow{\ \chi\ }\mathbb Z^2\longrightarrow0.
 \label{d4i:eq:split-exact-sequence}\tag{D4I28}
\]

For $\omega=(-1+\sqrt{-3})/2$, one has
$\omega+\overline\omega=-1$ and $\omega\overline\omega=1$.
Direct multiplication therefore proves
\[
                 (u+v\omega)(u+v\overline\omega)=u^2-uv+v^2.
 \label{d4i:eq:eisenstein-norm}\tag{D4I29}
\]
The quadratic factor in the original cubic is exactly this
Eisenstein norm, with the preceding coefficient two unchanged.

The original Euclidean metric on these same coordinates is,
by expansion,
\[
 \|\Phi(w,u,v)\|^2
            =12w^2+8w(u+v)+2(u^2+uv+v^2).
 \label{d4i:eq:full-slice-metric}\tag{D4I30}
\]
Write $d=(-3,1,1,1)$ for the direction when used in the metric
calculation. Its squared norm is twelve, and the exact pairing is
\[
          \langle\Phi(w,u,v),d\rangle=12w+4(u+v).
 \label{d4i:eq:direction-pairing}\tag{D4I31}
\]
Hence orthogonal projection away from the real line $\mathbb R d$
is given in the original coordinates by
\[
\begin{aligned}
 \pi\Phi(w,u,v)
  &=\Phi(w,u,v)-\left(w+\frac{u+v}{3}\right)d\\
  &=\left(0,\frac{2u-v}{3},\frac{2v-u}{3},-\frac{u+v}{3}\right).
\end{aligned}
 \label{d4i:eq:orthogonal-projection}\tag{D4I32}
\]
This depends only on the quotient coordinate $\chi(a)=(u,v)$.
Its squared norm is
\[
\begin{aligned}
 \|\pi\Phi(w,u,v)\|^2
  &=\frac{(2u-v)^2+(2v-u)^2+(u+v)^2}{9}\\
  &=\frac23(u^2-uv+v^2).
\end{aligned}
 \label{d4i:eq:projected-norm}\tag{D4I33}
\]
The quotient metric on $A/\mathbb Z d$, realized by this exact
projection, thus has Gram matrix
\[
              \frac13\begin{pmatrix}2&-1\\-1&2\end{pmatrix}
 \label{d4i:eq:quotient-gram}\tag{D4I34}
\]
in the displayed quotient coordinates $(u,v)$.
The retained common-lattice map $a\mapsto\sqrt2(a,a,a)$ multiplies
every squared norm by six, as proved in
\cref{wcl:thm:literal-intersection}. Consequently the corresponding
quotient Gram matrix on the image of $A$ in $\mathcal I$ is
\[
                         \begin{pmatrix}4&-2\\-2&4\end{pmatrix}.
 \label{d4i:eq:common-lattice-quotient-gram}\tag{D4I35}
\]
This is the exact quotient-metric calculation on the indicated
rank-three slice. It neither identifies that slice with the full
$D_4$ nor changes the original common-lattice metric.

% Integration-ready proof module; no canonical file has been edited.
% Requires amsmath, amssymb, amsthm, and hyperref in a standalone host.
% All coordinates and quadratic-form conventions are retained explicitly.

\section{The marked Niemeier lattice, its twelve theta indices, and the umbral shadow}
\label{nus:sec:shadow}

The original objects are
\[
 A=\{(x_0,\ldots,x_5)\in\mathbb Z^6:\sum_jx_j=0\},\qquad
 D=\{d\in\mathbb Z^4:\sum_id_i\equiv0\pmod2\},\qquad
 R=A^{\perp4}\perp D.
\]
The metrics are the displayed Euclidean metrics.  The four copies of $A$
retain their original fibre order $f_i289^j$ for $0\le j<6$, with fibre
labels $(00,10,01,11)$.  Their discriminant generator is
$\lambda=(5,-1,-1,-1,-1,-1)/6$.  The discriminant of $D$ retains
$10\mapsto v=(1,0,0,0)$, $01\mapsto s=(1,1,1,1)/2$, and
$11\mapsto c=(-1,1,1,1)/2$.  Consequently the original norm-valued
quadratic refinement, with codomain $\mathbb Q/2\mathbb Z$, is
\[
 q_R(k_1,k_2,k_3,k_4;y)
 =\frac56\sum_i k_i^2+\mathbf1_{y\ne0}\pmod{2\mathbb Z}.
\]
The original glue and Niemeier lattice are
\[
 \begin{split}
 E&=\{u\in\mathbb F_2^4:\sum_i u_i=0\},\qquad
 \tau(u)=(u_2+u_4,u_3+u_4),\\
 j(a,b)&=(a+b,a+2b,a,b),\qquad C_3=j(\mathbb F_3^2),\\
 H&=\{(3u+2j(a,b);\tau(u)):u\in E,\ (a,b)\in\mathbb F_3^2\},\\
 N&=\{x\in R^\vee:x+R\in H\}.
 \end{split}
\]
Thus $|H|=8\cdot9=72$ and $N/R=H$, with the root system
$X=A_5^4D_4$ proved in the retained lattice calculation.  In particular,
no norm-two vectors are introduced by the nonzero glue classes: their
original minima are four or six.  We use this proved literal root system,
rather than a rank-only identification, throughout the construction below.

\subsection{The primary prescription and its complete coefficient matrix}

Write $e(t)=\exp(2\pi\mathrm i t)$ and $q=e(\tau)$ for
$\tau\in\mathbb H$.  Cheng--Duncan--Harvey, arXiv:1307.5793v2,
Section~3.2 \cite[Sections~3.2--3.3]{esn:CDH2014}, source labels
\texttt{thetexp}, \texttt{transf\_theta},
and \texttt{transf\_theta\_2}, use
\[
 \theta_{m,r}(\tau,z)
 =\sum_{k\equiv r\ (2m)}q^{k^2/(4m)}e(kz),
 \qquad r\in\mathbb Z/(2m).
\]
Their Section~3.3, source label \texttt{def:OmegaMatrices}, defines
\[
 \Omega_m(n)_{r,s}
 =\begin{cases}
 1,&r+s\equiv0\pmod{2n},\quad r-s\equiv0\pmod{2m/n},\\
 0,&\text{otherwise},
 \end{cases}
 \qquad n\mid m.
\]
Their ADE table, source label \texttt{ADE1}, assigns
\[
 \Omega^{A_{m-1}}=\Omega_m(1),\qquad
 \Omega^{D_{m/2+1}}=\Omega_m(1)+\Omega_m(m/2),
\]
and extends by addition over components with the same Coxeter number.
These are the primary authors' prescribed maps.  The following calculation
specializes and proves their values in the present marked coordinates; it
does not assert a new ADE classification.

The six-cycle $e_j\mapsto e_{j+1\bmod6}$ on an $A$ fibre has
characteristic polynomial $(x^6-1)/(x-1)$ on its augmentation space.
Its order is six.  In $D$, take the original simple roots
$\alpha_1=e_1-e_2$, $\alpha_2=e_2-e_3$, $\alpha_3=e_3-e_4$,
$\alpha_4=e_3+e_4$.  Their reflection product
$s_{\alpha_1}s_{\alpha_2}s_{\alpha_3}s_{\alpha_4}$ sends
\[
 e_1\mapsto e_2,\quad e_2\mapsto e_3,\quad
 e_3\mapsto-e_1,\quad e_4\mapsto-e_4.
\]
It has order six and characteristic polynomial
\[
 (x^3+1)(x+1)
 =(x^2-x+1)(x+1)^2
 =\frac{(x^2-1)(x^6-1)}{(x-1)(x^3-1)}.
\]
Thus the original five components have Coxeter number $h=m=6$, and the
primary theta prescription has precisely $2h=12$ indices and denominator
$4h=24$ in its exponents.  The root rank remains $4\cdot5+4=24$;
the glue index remains $72$.

Put $J=\Omega_6(3)$.  Its two congruences are
\[
 r+s\equiv0\pmod6,\qquad r-s\equiv0\pmod4.
\]
The simultaneous solutions modulo twelve consist of the single value
$s\equiv5r\pmod{12}$: this value satisfies both congruences, and the
difference of any two solutions is divisible by both six and four, hence
by twelve.  Therefore
\[
 \boxed{\Omega^X=5I_{12}+J,\qquad J_{r,s}=\mathbf1_{s\equiv5r\ (12)}.}
\]
In the full original index order $0,1,\ldots,11$ this is
\begingroup
\csname c@MaxMatrixCols\endcsname=12\relax
\[
\Omega^X=\begin{pmatrix}
6&0&0&0&0&0&0&0&0&0&0&0\\
0&5&0&0&0&1&0&0&0&0&0&0\\
0&0&5&0&0&0&0&0&0&0&1&0\\
0&0&0&6&0&0&0&0&0&0&0&0\\
0&0&0&0&5&0&0&0&1&0&0&0\\
0&1&0&0&0&5&0&0&0&0&0&0\\
0&0&0&0&0&0&6&0&0&0&0&0\\
0&0&0&0&0&0&0&5&0&0&0&1\\
0&0&0&0&1&0&0&0&5&0&0&0\\
0&0&0&0&0&0&0&0&0&6&0&0\\
0&0&1&0&0&0&0&0&0&0&5&0\\
0&0&0&0&0&0&0&1&0&0&0&5
\end{pmatrix}.
\]
\endgroup
The unary series and the actual shadow are
\[
 S_{6,r}(\tau)=\frac1{2\pi\mathrm i}
  \left.\frac{\partial\theta_{6,r}}{\partial z}\right|_{z=0}
 =\sum_{k\equiv r\ (12)}kq^{k^2/24},
 \qquad S^X=\Omega^X S_6.
\]
All series and their derivatives converge normally on compact subsets of
their stated domains, by Gaussian decay.  Substitution $k\mapsto-k$
proves $S_{6,-r}=-S_{6,r}$, so $S_{6,0}=S_{6,6}=0$.
Let $\mathcal A:\mathbb C^5\to\mathbb C^{12}$ insert zero in positions
zero and six and put $(\mathcal Au)_{12-r}=-u_r$ for $1\le r\le5$.
This retains, and explicitly reconstructs, all twelve components.  Then
\[
 \Omega^X\mathcal A=\mathcal A M^X,\qquad
 \boxed{M^X=\begin{pmatrix}
 5&0&0&0&1\\0&4&0&0&0\\0&0&6&0&0\\
 0&0&0&4&0\\1&0&0&0&5
 \end{pmatrix}.}
\]
Indeed its $(r,s)$ entry is
$\Omega^X_{r,s}-\Omega^X_{r,12-s}$; direct substitution gives the matrix.
In particular the fourth A-type multiplicity and the full D-type
contribution have both been retained.  The diagonal entries
$(5,4,6,4,5)$ are exactly the Coxeter exponent multiplicities:
four copies of $(1,2,3,4,5)$ and one copy of $(1,3,3,5)$.

\subsection{Quadratic refinements and the finite theta transformations}

Let $L=\mathbb Z\ell$ with $(\ell,\ell)=12$.  Its dual is
$L^\vee=\mathbb Z\ell/12$.  Hence
\[
 A_L=L^\vee/L=\mathbb Z/12,\qquad
 q_L(r)=r^2/12\pmod{2\mathbb Z},\qquad
 Q_L(r)=q_L(r)/2=r^2/24\pmod{\mathbb Z},
\]
and the associated bilinear form is $B_L(r,s)=rs/12\pmod{\mathbb Z}$.
The symbols $q_L$ and $Q_L$ deliberately have different codomains and
retain the factor two.

Define the unitary Fourier matrices and reflection by
\[
 (F_\pm)_{r,s}=12^{-1/2}e(\pm rs/12),\qquad
 R_{r,s}=\mathbf1_{s\equiv-r\ (12)}.
\]
Finite geometric-series summation gives
\[
 F_\pm^2=R,\qquad F_-=RF_+=F_+R,\qquad F_+F_-=I.
\]
The substitution $k\mapsto k+12$ shows that $e(k^2/24)$ depends only on
$r=k\bmod12$.  Thus, with
$D_{r,s}=e(r^2/24)\delta_{r,s}$,
\[
 \theta_6(\tau+1,z)=D\theta_6(\tau,z).
\]
Poisson summation of the Gaussian gives
\begin{align*}
 \theta_6(-1/\tau,z/\tau)
 &=\sqrt{-\mathrm i\tau}\,e(6z^2/\tau)F_-\theta_6(\tau,z),\\
 \theta_6(-1/\tau,-z/\tau)
 &=\sqrt{-\mathrm i\tau}\,e(6z^2/\tau)F_+\theta_6(\tau,z).
\end{align*}
Here $\sqrt{-\mathrm i\tau}$ is the holomorphic square root positive
when $\tau$ is positive imaginary.  For completeness, apply the Fourier
transform convention $\widehat f(u)=\int_{\mathbb R}f(x)e(-ux)\,dx$
to $f(x)=\exp(\pi\mathrm i\tau x^2/12+2\pi\mathrm i zx)$.
Completing the square evaluates its transform; the real positive Gaussian
integral is obtained by squaring the integral and using polar coordinates,
and holomorphic continuation gives the expression for $\Im\tau>0$.
Periodizing $f(r+12x)$, its $n$-th Fourier coefficient is
$12^{-1}e(nr/12)\widehat f(n/12)$.  Summing these coefficients at zero
and grouping $n$ by its residue modulo twelve gives the first displayed
transformation.  Absolute Gaussian convergence justifies the integration,
grouping, and Fourier evaluation.  Replacing $z$ by $-z$ and using
$\theta_{6,-r}(\tau,z)=\theta_{6,r}(\tau,-z)$ gives the second.
The second, with its plus-sign Fourier matrix, is precisely CDH's
convention; the first uses the usual generator
$S=\left(\begin{smallmatrix}0&-1\\1&0\end{smallmatrix}\right)$
with elliptic argument $z/\tau$.

The matrices carrying the finite multiplier for this first convention are
\[
 U(T)=D,\qquad U(S)=e(-1/8)F_-.
\]
The square-root weight factor has not been included in $U$.
Differentiation of the theta transformation at $z=0$ gives
\[
 S_6(-1/\tau)=\tau\sqrt{-\mathrm i\tau}F_-S_6(\tau),\qquad
 S_6(\tau+1)=D S_6(\tau).
\]
Equivalently, CDH's reflected convention gives
$S_6(-1/\tau)=-\tau\sqrt{-\mathrm i\tau}F_+S_6(\tau)$.
The equality follows from $RS_6=-S_6$.  Thus the unary series has weight
$3/2$ with the same finite multiplier; theta has weight $1/2$.

The permutation $J$ preserves $Q_L$ because
$((5r)^2-r^2)/24=r^2\in\mathbb Z$.  Since $5^{-1}=5$ modulo twelve,
\[
 (JF_\pm)_{r,s}=12^{-1/2}e(\pm5rs/12)
 =(F_\pm J)_{r,s}.
\]
Therefore $\Omega^X$ commutes with both finite multiplier matrices.
The shadow transformation follows by applying this proved intertwiner.
Its component series vanish at infinity since every nonzero summand has
strictly positive exponent; the modular transformation shows the same
cuspidal behavior at every cusp of the full modular group.

There is also an exact comparison with the original A5 discriminant,
without replacing its form.  Reduction
\[
 p:\mathbb Z/12\longrightarrow\mathbb Z/6,\qquad p(r)=r\bmod6
\]
has kernel $\{0,6\}$ and satisfies
\[
 q_A(p(r))=\frac56r^2=10q_L(r)\pmod{2\mathbb Z}.
\]
Thus it is a surjective quadratic similitude with precisely multiplier
ten, and its bilinear forms satisfy the same multiplier relation.
It is not asserted to be an isometry.  Applying $p$ to four coordinates
and the identity to the retained $D$ discriminant gives
\[
 \widetilde D=(\mathbb Z/12)^4\oplus\mathbb F_2^2
 \longrightarrow D_R,
 \qquad (r;y)\longmapsto(p(r_1),\ldots,p(r_4);y).
\]
The pullback quadratic refinement is exactly
$10\sum_iq_L(r_i)+\mathbf1_{y\ne0}$, its radical is
$\{0,6\}^4\oplus0$, and its quotient by that radical is the original
$(D_R,q_R)$.  To prove the radical claim, pairing with each coordinate
generator forces $5r_i/6\in\mathbb Z$, hence $6\mid r_i$; the
nondegenerate $D$ form forces $y=0$.  The vectors so obtained have
quadratic value $30$ in each nonzero coordinate, hence zero modulo two.
The inverse image of the original $H$ has order $16\cdot72=1152$ and
is isotropic for this pullback form.  This gives the exact relation of the
theta indices to all original glue coordinates, with the form and
multiplier specified on both sides.

\subsection{The full Fourier determinant and the retained order-twelve character}
\label{nus:subsec:determinant}

Let $\chi:\mathrm{SL}_2(\mathbb Z)\to\mathbb Z/12$ be the retained
marking
\[
 \chi(T)=-1,\qquad \chi(S)=3.
\]
These values respect the presentation
$S^4=I$, $(ST)^3=S^2$: in the abelian target the required relations
are $4\cdot3=0$ and $3(3-1)=2\cdot3$ modulo twelve.
They therefore define a character, and its value at $T$ makes it onto.
We now calculate the determinant of the entire twelve-dimensional
finite theta representation, rather than that of a reduced component
list.

Put $\zeta=e(1/12)$.  The Vandermonde determinant gives
\[
 \det F_+=12^{-6}\prod_{0\le j<k\le11}(\zeta^k-\zeta^j).
\]
For every such pair,
\[
 \zeta^k-\zeta^j
 =2\mathrm i\sin\!\left(\frac{\pi(k-j)}{12}\right)
       \exp\!\left(\frac{\pi\mathrm i(j+k)}{12}\right),
\]
whose sine is strictly positive.  There are $\binom{12}{2}=66$ pairs,
and each index occurs in eleven of them, so
\[
 \sum_{j<k}(j+k)=11\sum_{j=0}^{11}j=11\cdot66=726.
\]
The phase of the determinant is consequently
\[
 \mathrm i^{66}\exp(726\pi\mathrm i/12)=(-1)\mathrm i=-\mathrm i.
\]
Its remaining factor is positive.  Unitarity, proved by the finite
geometric sum above, makes its absolute value one.  It follows exactly
that
\[
 \boxed{\det F_+=-\mathrm i,\qquad\det F_-=\mathrm i.}
\]
In particular retaining CDH's reflection changes the Fourier determinant
by $\det R=(-1)^5=-1$; it cannot be dropped.

The sum of the squares of all twelve indices is
\[
 \sum_{r=0}^{11}r^2=\frac{11\cdot12\cdot23}{6}=506.
\]
Consequently
\[
 \boxed{\det U(T)=e(506/24)=e(1/12),\qquad
 \det U(S)=e(-12/8)\det F_-=-\mathrm i=e(-3/12).}
\]
For an explicit check that the finite matrices have the relevant
metaplectic relations, let
$G=\sum_{r=0}^{11}e(r^2/24)$.  Listing the squares modulo twenty-four
gives
\[
 G=1+4e(1/24)+2e(1/6)+2e(3/8)+2e(2/3)+e(1/2).
\]
The first and last terms cancel, as do the terms with arguments $1/6$
and $2/3$.  The elementary sine and cosine values at $\pi/12$ and
$3\pi/4$ give
\[
 4e(1/24)+2e(3/8)=\sqrt6(1+\mathrm i)
 =\sqrt{12}\,e(1/8).
\]
Completing the square in the finite sum proves
\[
 (F_-D)^2=e(1/8)D^{-1}F_-,\qquad
 (F_-D)^3=e(1/8)R.
\]
For the first equality, its $(r,s)$ entry before evaluation is
$12^{-1}e(s^2/24)\sum_te(t^2/24-t(r+s)/12)$; shifting $t$ by
$r+s$ produces exactly the stated Gauss sum and matrix entry.
The second follows by multiplication and $F_-^2=R$.
It follows that
\[
 (U(S)U(T))^3=U(S)^2=-\mathrm iR,\qquad U(S)^4=-I_{12}.
\]
Taking determinants kills this metaplectic central sign because the
dimension twelve is even.  The determinant therefore descends to an
ordinary character of $\mathrm{SL}_2(\mathbb Z)$.  Its two computed
generator values, together with generation by $S,T$, prove
\[
 \boxed{\det U(\gamma)=e(-\chi(\gamma)/12).}
\]
The dual determinant is its inverse and is exactly
$e(\chi(\gamma)/12)$, with the retained signs $T\mapsto e(-1/12)$,
$S\mapsto\mathrm i$.

This is a statement about the determinant multiplier, with its domain
and codomain specified.  It is not a statement that a vector of weight
$1/2$ has scalar weight one after taking a determinant.  The separate
Dedekind-eta identities
\[
 \eta(\tau+1)=e(1/24)\eta(\tau),\qquad
 \eta(-1/\tau)=\sqrt{-\mathrm i\tau}\,\eta(\tau)
\]
show that $\eta^2$ has weight one and multiplier $\det U$, whereas
$\eta^{-2}$ has weight minus one and multiplier $(\det U)^{-1}$.
Indeed $\eta^2(S\tau)=(-\mathrm i)\tau\eta^2(\tau)$ and
$\eta^{-2}(S\tau)=\mathrm i\tau^{-1}\eta^{-2}(\tau)$.
The integer-weight factors $(c\tau+d)^{\pm1}$ obey the cocycle identity
by matrix multiplication, so the two generator formulas prove the
statements for every $\gamma$.  At $-I$ the character and odd integer
weight each contribute a minus sign, whose product is one.

\subsection{The residue Fourier map and what the Weyl quotient retains}

The original marking is
$\kappa(289^j)=j\pmod6$, with the full ordered residue list
\[
 (289^j)_{j=0}^5=(1,289,361,169,121,529)\pmod{840}.
\]
On the six-dimensional complex coordinate carrier define
\[
 E_2:\mathbb C[S_{840}]\longrightarrow\mathbb C[\mathbb Z/12],
 \qquad E_2(e_{289^j})=e_{2j},
\]
and on a separately indexed $\mathbb C^6$ define
\[
 B(e_k)=\frac{e_k+e_{k+6}}{\sqrt2},\qquad0\le k<6.
\]
Both are injective and preserve the indicated coordinate Hermitian
products.  Retain the original negative-sign normalized Fourier matrix
$F_6^-=(6^{-1/2}e(-kj/6))_{k,j}$.  Entry-by-entry calculation gives
\[
 \boxed{F_-E_2=B F_6^-.}
\]
Indeed the $(r,j)$ entry of the left side is
$12^{-1/2}e(-rj/6)$, and the right side has the same value, using
$r\bmod6$.  Thus the original unnormalized sample transform is carried
into this formula by its exact factor $\sqrt6$.
Translation $j\mapsto j+1$ intertwines under $E_2$ with translation
$r\mapsto r+2$ modulo twelve.  This gives a literal Fourier and
translation map, not an identification of different quadratic forms.
In particular $Q_L(6)=3/2=1/2\pmod{\mathbb Z}$, so the quotient
$\mathbb Z/12\to\mathbb Z/6$ does not carry $Q_L$ itself to a
quadratic form on the quotient.  Its exact multiplier-ten comparison
with the original A5 form was proved above.

For the original lattice action, $\rho(289)$ is the same six-cycle in
each of the four $A$ coordinate sets and is the identity on $D$.
Each coordinate transposition of $A$ is reflection in the corresponding
root $e_i-e_j$, so the full residue action is in
$W(R)=W(A)^4\times W(D)$.  It acts trivially on every A discriminant
because
$Q_\sigma\lambda-\lambda=e_{\sigma(0)}-e_0\in A$, and it fixes
the D discriminant.  Consequently
\[
 S_{840}\xrightarrow{\rho}O(N)\longrightarrow O(N)/W(R)
\]
is the trivial homomorphism.  This proves exactly which residue-group
action is killed, without discarding its already constructed Fourier map.

The original marked quotient is
\[
 G=O(N)/W(R)\cong\mathrm{GL}_2(\mathbb F_3).
\]
Here its actual coordinate map uses
$\ell_1=(1,1)$, $\ell_2=(1,2)$, $\ell_3=(1,0)$,
$\ell_4=(0,1)$.  For $A\in\mathrm{GL}_2(\mathbb F_3)$ determine
the unique signs and permutation by
\[
 \ell_iA=\varepsilon_i\ell_{p(i)},\qquad
 (P_pu)_i=u_{p(i)},\qquad B_p\tau(u)=\tau(P_pu).
\]
The four row lines exhaust the projective line, so the signs and
permutation exist uniquely.  The binary equation defines $B_p$ uniquely
since $\tau$ is onto.  Then $(k;y)\mapsto(\varepsilon_i k_{p(i)};B_py)$
preserves precisely the original two glue codes, and its ternary
restriction is $M_gj=jA$.  Conversely a marked glue stabilizer gives this
same $A$.  Multiplication of these identities proves the group
isomorphism, with the original row/column convention retained.

To compare its chosen signed lifts with CDH's simple-root permutation
section, let $Q_{\rm rev}e_j=e_{5-j}$ in an $A$ fibre.  When
$\varepsilon_i=-1$, replace the representative $-I$ by
$-Q_{\rm rev}$; their quotient is $Q_{\rm rev}\in W(A)$ and
$-Q_{\rm rev}(e_j-e_{j+1})=e_{4-j}-e_{5-j}$.
Thus the resulting map permutes the original simple roots by diagram
reversal.  When $\varepsilon_i=1$ no change is needed.
For D, put $\Delta=(\alpha_1\ \alpha_2\ \alpha_3\ \alpha_4)$.
Associate the original classes $v,s,c$ with outer nodes $1,4,3$,
respectively; their fundamental weights are
$e_1$, $(1,1,1,1)/2$, and $(1,1,1,-1)/2$, the last of which differs
from the retained $c$ by $(1,0,0,-1)\in D$.
For $B_p$, let $P_{B_p}$ permute these outer nodes exactly as $B_p$
permutes $v,s,c$ and fix node two.  Then
\[
 d_{B_p}=\Delta P_{B_p}\Delta^{-1}
\]
is the exact original-coordinate diagram representative.  Its
orthogonality follows because permuting the outer nodes preserves the
original Gram matrix.  It preserves $D$, and its action on fundamental
weights proves its discriminant action is exactly $B_p$.
Its quotient with the originally chosen D lift has trivial discriminant
action, hence belongs to the proved Weyl kernel $W(D)$.
This specifies the actual change of section on every summand.

Define the following four class functions using those retained maps:
\[
 a_g=\#\{i:p(i)=i\},\quad
 b_g=\sum_{p(i)=i}\varepsilon_i,\quad
 s_g=\operatorname{sgn}(B_p|_{\{v,s,c\}}),\quad
 c_g=\#\operatorname{Fix}(B_p|_{\{v,s,c\}})-1.
\]
The first two are the traces on the four-component permutation and
signed permutation representations.  The last is the trace on the
two-dimensional sum-zero subspace of the three outer-node permutation
representation, and $s_g$ is its sign character.  Thus these are precisely
CDH's $(\bar\chi^{X_A},\chi^{X_A},\chi^{X_D},\check\chi^{X_D})$;
the remaining $\bar\chi^{X_D}$ is one because the single D component
is fixed.  The comparison of sections above proves this identification
for the actual marked quotient.

CDH Section~5.1 \cite[Section~5.1]{esn:CDH2014} defines the full
twisted coefficient matrix by
\[
 (\Omega_g^{X_A})_{r,s}
 =\bigl(b_g\mathbf1_{r\ \text{even}}+
          a_g\mathbf1_{r\ \text{odd}}\bigr)\delta_{r,s},
\]
and, for its exceptional $m=6$ D4 case, by
\[
 (\Omega_g^{X_D})_{r,s}=
 \begin{cases}
 c_g\delta_{r,s},&r\equiv0,3\pmod6,\\
 \delta_{r,s}+s_gJ_{r,s},&r\equiv1,5\pmod6,\\
 s_g\delta_{r,s}+s_gJ_{r,s},&r\equiv2,4\pmod6.
 \end{cases}
\]
All deltas in this formula have indices modulo twelve.  The exact maps
from the marked glue quotient to the prescribed twisted shadows are
\[
 g\longmapsto\Omega_g^X=\Omega_g^{X_A}+\Omega_g^{X_D},\qquad
 S_g^X=\Omega_g^XS_6.
\]
They are class-function constructions from traces of representations;
the assignment $g\mapsto\Omega_g^X$ is not asserted to be a group
representation.  On the odd five-component presentation the explicit
matrix is
\[
 \boxed{M_g^X=\begin{pmatrix}
 a_g+1&0&0&0&s_g\\
 0&b_g&0&0&0\\
 0&0&a_g+c_g&0&0\\
 0&0&0&b_g&0\\
 s_g&0&0&0&a_g+1
 \end{pmatrix}.}
\]
This follows by the same signed component reconstruction $\mathcal A$:
in rows two and four the two D terms cancel, in row three the contribution
is $c_g$, and rows one and five pair with one another.
For reference, all possible character tuples are
\[
\begin{array}{c|rrrr|r}
 [g]&a_g&b_g&s_g&c_g&\text{number of elements}\\\hline
 1A&4&4&1&2&1\\
 2A&4&-4&1&2&1\\
 2B&2&0&-1&0&12\\
 4A&0&0&1&2&6\\
 3A&1&1&1&-1&8\\
 6A&1&-1&1&-1&8\\
 8A,8B&0&0&-1&0&12
\end{array}
\]
with the CDH class names of their table
\texttt{tab:chars:eul:6}.  The formulas apply to each actual $A$, so
their validity does not depend on a choice of conjugacy-class names.
The finite checker obtains the table from all $48$ invertible matrices
over $\mathbb F_3$ and the original row forms.

The distinguished residue action has quotient $g=1$ and therefore
produces the untwisted shadow, while a product $\rho(r)L_g$ produces
the same $S_g^X$ for every $r$.  The four-component signed action and
the D4 triality quotient survive and enter the displayed matrix.  The
central element $-I\in\mathrm{GL}_2(\mathbb F_3)$, for example,
changes $b_g$ from four to minus four while fixing $a_g,s_g,c_g$;
it changes the even-index odd-space entries with exactly that sign.

The already marked composite
\[
 \mathrm{SL}_2(\mathbb Z)\xrightarrow{\chi}\mathbb Z/12
 \longrightarrow\mathbb Z/6\xrightarrow{j\mapsto289^j}S_{840}
\]
is onto, has kernel $\chi^{-1}(\{0,6\})$, and sends $T$ to
$289^5=529$.  Following it by the Weyl quotient kills it, by the
proved residue calculation.  Following $\chi$ into the dual determinant
multiplier retains its complete order twelve, by the determinant theorem.
These are explicit different outgoing morphisms from the same marking,
including their kernels and signs.

\subsection{The shadow-to-Fricke map with the original scalar seventy-two}

The Coxeter calculation above gives the original formal frame exponent
\[
 \pi_X=-5[1]+[2]-[3]+5[6].
\]
Let $\mathcal E(\sum a_d[d])=\prod_d\eta(d\tau)^{a_d}$.
It is a group homomorphism from addition of formal exponents to
multiplication of nonzero holomorphic functions, by the exponent laws.
The retained scalar function is
\[
 \boxed{T_6(\tau)=\mathcal E(-\pi_X)
 =\frac{\eta(\tau)^5\eta(3\tau)}
        {\eta(2\tau)\eta(6\tau)^5}.}
\]
The minus sign on $\pi_X$ is essential: $\mathcal E(\pi_X)=T_6^{-1}$.

Here is a direct coefficient map from the actual shadow matrix to this
eta function.  Its diagonal difference vector is
$\alpha=(5,4,6,4,5)$, and set $\alpha_0=0$ modulo six.
Define the Lambert-series map at this marked matrix by
\[
 \mathscr L(\Omega^X)(\tau)
 =1+\sum_{n\ge1}\alpha_{n\bmod6}
                    \frac{nq^n}{1-q^n}.
\]
Convergence is absolute and locally uniform since $|q|<1$ and
$\alpha$ is bounded.  The coefficient of $q^N$ is exactly
$\sum_{d\mid N}d\alpha_{d\bmod6}$; thus this specifies every
coefficient, not only its initial terms.
For all positive integers $n$ the retained multiplicities satisfy
\[
 \alpha_{n\bmod6}
 =5-\mathbf1_{2\mid n}+\mathbf1_{3\mid n}-5\mathbf1_{6\mid n}.
\]
Checking the six residues proves this identity.  The product definition
$\eta(\tau)=q^{1/24}\prod_{n\ge1}(1-q^n)$ and termwise logarithmic
differentiation then give
\[
 \boxed{f^X:=\mathscr L(\Omega^X)
 =-q\frac{d}{dq}\log T_6
 =5\lambda_6+\lambda_2-\lambda_3,}
\]
where the original functions on the right are
\[
 \lambda_d=q\frac{d}{dq}\log\frac{\eta(d\tau)}{\eta(\tau)}
 =\frac{d-1}{24}+\sum_{n\ge1}\sigma_1(n)(q^n-dq^{dn}).
\]
The constant is
$[5(6-1)+(2-1)-(3-1)]/24=1$, and the other coefficients agree by
the six-residue identity.  This proves the shadow-to-eta calculation
directly for every coefficient.

The precise primary connecting map is CDH
Section~4.2 \cite[Section~4.2]{esn:CDH2014}:
\[
 \sigma^X(\tau,z)=\sum_{r\ (12)}\overline{S_r^X(\tau)}
                                \theta_{6,r}(\tau,z),\qquad
 \mathscr S_{1,1}(\sigma^X)=f^X.
\]
Its source labels are \texttt{eqn:sigmaX}, \texttt{def:weight2},
\texttt{weight2:Aseries}, and \texttt{weight2\_EZ}.
The displayed Lambert map is its explicit specialization at this
root-system shadow.  It verifies the resulting function independently
of the general Shimura-lift theorem.  In particular the shadow and the
retained Fricke function are connected by this complete coefficient map.

The original modular subgroup here is $\Gamma_0(6)$, and
$[\mathrm{SL}_2(\mathbb Z):\Gamma_0(6)]=12$.  To compute the index,
reduction modulo six gives
$\mathrm{SL}_2(\mathbb Z/6)\cong
\mathrm{SL}_2(\mathbb F_2)\times\mathrm{SL}_2(\mathbb F_3)$ of
order $6\cdot24=144$.  Reduction is onto: elementary upper and lower
unipotent matrices generate over each field, and CRT choices of their
entries realize each factor while acting trivially on the other.
The upper triangular subgroup modulo six has $\varphi(6)\cdot6=12$
elements (choose its first unit diagonal entry and its upper-right entry;
the other diagonal entry is its inverse).  Its inverse image is
$\Gamma_0(6)$, proving index $144/12=12$.
This modular-group index and the twelve theta indices have their
specified constructions; neither replaces the lattice index seventy-two.

Under the exact Fricke transformation $W_6\tau=-1/(6\tau)$,
the eta inversion formula sends each $\eta(dW_6\tau)$ to
$\sqrt{-\mathrm i(6/d)\tau}\,\eta((6/d)\tau)$.
The retained exponent vector at $d=1,2,3,6$ is $(5,-1,1,-5)$;
its reversed vector is its negative and its exponent sum is zero.
The $\tau$ factors therefore cancel.  Since each $6/d$ is positive,
the consistent square-root branch gives the positive scalar
\[
 \prod_{d\mid6}(6/d)^{a_d/2}
 =6^{5/2}3^{-1/2}2^{1/2}=72.
\]
Consequently
\[
 \boxed{T_6(W_6\tau)=72/T_6(\tau).}
\]
For the weight-two slash action of the determinant-six representative,
\[
 (f|_2W_6)(\tau)=(6\tau^2)^{-1}f(-1/(6\tau)).
\]
Differentiating the last identity and using
$d(W_6\tau)/d\tau=(6\tau^2)^{-1}$ gives
\[
 \boxed{f^X|_2W_6=-f^X.}
\]
On the Jacobi side, the primary Eichler--Zagier operator
$\mathcal W_6(6)$ acts on theta by $\Omega_6(6)=R$.
This is verified by its defining congruences, $r+s=0\pmod{12}$,
$r-s=0\pmod2$.  On the odd shadow and on the corresponding
skew-Jacobi form its eigenvalue is minus one.  The specialized coefficient
map just proved carries that negative eigenline to the negative Fricke
eigenline of $f^X$.  This proves an exact equivariant linear comparison
on these lines; it does not evaluate an $\mathrm{SL}_2(\mathbb Z)$
character on the determinant-six matrix.

Finally, the original sheet coordinates remain
\[
 J_6=T_6+5+72/T_6,\qquad Y_6=T_6-72/T_6.
\]
Direct substitution and differentiation prove
\[
 J_6(W_6\tau)=J_6(\tau),\quad
 Y_6(W_6\tau)=-Y_6(\tau),\quad
 Y_6^2=(J_6-5)^2-288,\quad
 dJ_6=Y_6\,dT_6/T_6=-2\pi\mathrm iY_6 f^X\,d\tau.
\]
Thus the original scalar $72$, the additive shift $5$, the branch
constant $288$, and the differential factor $-2\pi\mathrm i$ all
survive the complete shadow-to-sheet calculation.

\subsection{Established scope and reproducibility}

The constructed maps are the ADE coefficient specialization, the exact
finite-theta intertwiner, the multiplier determinant and its dual, the
quadratic similitude to the original discriminants, the marked glue-group
trace prescription, the original six-to-twelve Fourier map, and the
coefficient-level shadow-to-Fricke map.  Their kernels, weights, signs,
and constants were specified and proved above.  They establish concrete
relations even where a direct isometry or an identical group action is
absent.  They do not construct a new graded moonshine module or a new
moonshine theorem.  The name ``umbral shadow'' here refers to precisely
the CDH prescription, with its source provenance retained.

The standard-library program \texttt{check\_shadow.py} uses exact
arithmetic in $\mathbb Q[z]/(z^8-z^4+1)$ with $z=e(1/24)$.
It verifies both full and odd matrices, the Fourier determinant by the
Vandermonde product, the Gauss sum, the character signs, the Fourier
embedding, all $48$ original-code character records, the first $240$
Lambert coefficients, and the squared positive Fricke scalar $72^2$.
The proofs above establish the unbounded identities independently of
those finite checks.

